\def\ChapterCode{PKU}
\chapter
[\CO{[PKU]} Erweiterte Physikalisch-körperliche Untersuchung]
{Erweiterte Physikalisch-körperliche Untersuchung}
\label{chp:PKU}



\section*{Danksagung und Hinweise}

\begin{figure}[H]

\Captionof{figure}{\EE{Danke!} Die Autoren bedanken sich für
die tatkräftige Mitarbeit bei (vordere Reihe von links):  Angela Hausladen,
Manuel Scharrer,
Brigitte Mozga,
Andre Hermann.}
\end{figure}


\noindent Wir danken \ldots
\begin{itemize}
  \item[\ldots] für die zahlreichen Anregungen der
  Kursteilnehmer der vergangenen Jahre.
  \item[\ldots] Stephan Kranl für die Anmerkungen und
  Korrekturen, sowie seinen Einsatz während des stressigen
  Endspurts.
  \item[\ldots] der {\sffamily Arbeitsgemeinschaft {\color{ubuntured}Arbeits- und Ausbildungsstandards} {\color{darkBlue}für den
  Sanitätsdienst}} (\ARGEAASS) für die
  gute Zusammenarbeit und Unterstützung.
\end{itemize}

\paragraph{Hinweise}
Wie jede Wissenschaft ist die Medizin ständigen Ent\-wicklungen
unterworfen. Soweit fachliche Angaben ge\-macht werden, darf der Leser
zwar darauf ver\-trau\-en, dass die Autoren große Sorg\-falt
da\-rauf verwandt haben, dass diese An\-ga\-ben dem Wissens\-stand bei
Ferti\-gstel\-lung des Werkes entsprechen. Für Angaben,
spe\-ziell zu Dosierungsanweisungen und
Ap\-plikations\-formen, kann jedoch keine Ge\-währ übernommen werden. 

Jeder Leser ist angehalten, durch
sorg\-fältige Prü\-fung der Literatur, der Lehrmeinung so\-wie der
Beipackzettel der ver\-wendeten Präparate und gegebenenfalls
mittels Konsultation eines Spezialisten fest\-zustellen, ob
die betreffende Aussage, Maßnahme oder Empfehlung für
Do\-sierungen oder die Beachtung von Kon\-tra\-indi\-ka\-tionen gegenüber der Angabe in diesem
Werk abweicht. Jede Maßnahme, Dosierung oder Ap\-pli\-kation
erfolgt auf eigene Gefahr des Be\-nutzers. 
Die Herausgeber appellieren an jeden Le\-ser, 
auffallende Un\-ge\-nau\-igkei\-ten umgehend mitzuteilen.

Geschützte Waren\-namen und Wa\-ren\-zei\-chen werden nicht besonders
kenntlich ge\-macht. Aus dem Fehlen eines solchen Hin\-weises kann also
nicht ge\-schlossen werden, dass es sich um einen freien Waren\-namen
handelt.
Die in diesem Skriptum verwendeten Personenzeichnungen wie z.\,B.
\enquote{Patient},
\enquote{Untersucher} bzw.
\enquote{Arzt} sind geschlechtsneutral, also
sowohl für weibliche als auch männliche Personen zu verstehen.



\section{Vorwort zur Lehrveranstaltung}

\Text{Kleidung und Zubehör}{%
Bitte erscheinen Sie beim Praktikum in einem
\EE{weißen Mantel} und bringen Sie \EE{immer} das
Kursskriptum und folgende Gegenstände
mit:
\begin{itemize}
  \item einen \EE{Stift},
  \item eine \EE{Uhr
  mit Sekundenzeiger},
  \item ein \EE{Stethoskop},
  \item einen \EE{Perkussionshammer} und
  \item eine \EE{Untersuchungslampe}.
\end{itemize}

Sie sollen die basalen Untersuchungstechniken am lebenden
menschlichen Körper erlernen, d.\,h. Hören, Fühlen, Tasten,
etc.
Dies bedingt, dass Sie einen Mitstudiereden, der
die entsprechenden Körperregionen freimacht, untersuchen, und, dass
Sie sich ebenso untersuchen lassen müssen. Es wird daher dringend
empfohlen, \EE{leichte Sportbekleidung (kurzärmeliges
T-Shirt, kurze Hose)} für das praktische Lernen und Üben
anzulegen. Für Umkleidemöglichkeiten und Sichtschutz in den
Kursräumen ist gesorgt.
}{%
}{}{}{}

\Fazit{Falls von Ihrem Lehrenden Stationsbesuche geplant sind,
nehmen Sie hierzu unbedingt einen  \EE{sauberen} (!),
\EE{geruchsneutralen} (!) weißen Mantel mit. }
\Cave{Vergilbte \enquote{Sezierkursmäntel} sind nicht akzeptabel!}

\Text{Strukturierte Beobachtung / Checklisten / Famulaturpropädeutikum}{%
Im am Ende des Semesters stattfindenden Famulaturpropädeutikum
werden wir Sie bitten,
6 Aufgaben (2 ärztliche Grundfertigkeiten, 2 ärztliche
Gesprächsführungen, 2 physikalische Krankenuntersuchungen)
vorzuführen, die Sie im Laufe dieser Kurse gelernt und geübt haben.
Die Unterrichtenden der Kurse werden Sie dabei beobachten, und Ihre
Leistung mittels einer Checkliste protokollieren. Im Anschluss erhalten
Sie einen Feedbackausdruck über die von Ihnen erbrachte Leistung. 
Die Checklisten zum Untersuchungsgang (siehe Wochenthemen
dieses Kursskriptums) dienen als Grundlage der Beurteilung
für die PKU-Stationen \CC{[A1, B1, B2, C1, C2, D1, D2,
E1, E2]}. Es wird nur Ihr Handeln beurteilt und keine
Theorie abgeprüft.
}{}{}{}{}


\Text{Mitarbeit und Übungszeit}
{%
Da das Famulaturpropädeutikum knapp vor der SIP 2
angesiedelt ist, und die Ersatzleistung in der Durchführung
weiterer Stationenaufgaben besteht und nicht sichergestellt
werden kann, ob Ersatztermine vor Oktober bereitgestellt
werden können, bitten wir, in Ihrem eigenen Interesse im
Praktikum intensiv mitzuarbeiten und sich auf das
Famulaturpropädeutikum ernsthaft vorzubereiten.

Voraussichtlich in den ersten drei Wochen des
Famulaturpropädeutikums werden die Kursräume \VonBis{203}{218} zwischen 12:00 und 20:00 für
\EE{freiwilliges Üben} geöffnet sein. In Ihrem eigenen
Interesse raten wir dringend, von diesem Angebot Gebrauch zu machen. Genauere Infos zu Organisation und
Ablauf der \EE{freien Übungszeiten} entnehmen Sie bitte Ende
März den \enquote{study guides}.
Bitte beachten Sie, dass das Skriptum jedes Jahr
überarbeitet wird. Benutzen Sie daher zur Vorbereitung
unbedingt die \EE{aktuellste Version}.}
{}
{}
{}
{}




\Text{Weitere Möglichkeiten zum Selbststudium}
{Im DEMAW-\I{Computer\-lern\-studio} stehen Ihnen multimediale
Lernprogramme zur physikalischen Krankenuntersuchung zur Verfügung.
Wir empfehlen unter anderem die Lernprogramme:

\begin{itemize}
  \item Der klinische Status (Schwarzmeier)
  \item Wiener Auskultationstrainer (Schmidts)
  \item Klinische Untersuchung in der Pneumologie (Laennec)
\end{itemize}
% 
% \Cave{%
% Beachten Sie aber bitte auch für etwaige aktuelle Änderungen
% des organisatorischen Ablaufs und der Prüfungsmodi die
% Hinweise in den \enquote{study guides} auf der MUW Homepage.
% 
% {(Diplomstudium Humanmedizin oder Diplomstudium Zahnmedizin → 4.
% Semester → Skills Line → Physikalische Krankenuntersuchung).}
% }
}{}{}{}{}

\begin{Literary}
\fullcite{DahmerAnamneseBefund10}

\fullcite{ChecklisteAnamUnt1}

\fullcite{MiddekeAnamnese4}
\end{Literary}








% \cleardoublepage
\section{Die physikalische Krankenuntersuchung}

\subsection{Generelles}

\Text{Arbeitskleidung}{%
Tragen Sie bitte bei allen im Folgenden angeführten
Tätigkeiten eine angemessene, spitalsübliche
Arbeitskleidung (i.\,d.\,R. einen weißen Mantel)
Arbeitsmäntel werden zugeknöpft, mit hochgekrempelten
Ärmeleln getragen. 
Wechseln Sie die Arbeitskleidung nach jeder offensichtlichen
Kontamination, ansonsten -- wenn möglich -- täglich. 
Binden Sie lange Haare zurück und entfernen Sie
jeglichen Schmuck von Ihren Händen. Ebenso dürfen keine
Uhren, Armbänder oder ähnliches am Handgelenk getragen
werden. 

}{}{}{}{}

\Text{Händedesinfektion}
{\index{Händedesinfektion!Krankenuntersuchung}
Die Hände des Personals sind die häufigsten Keimüberträger
im klinischen Betrieb. Konsequente Händehygiene ist deshalb
eine der wirksamsten Maßnahmen zur Vermeidung nosokomialer
Infektionen. Händewaschen entfernt nur Schmutz und führt zu
einer unwesentlichen Keimreduktion der Haut.

\Cave{Die hygienische Händedesinfektion soll transiente
Keime beseitigen (Mikrobenspezies, die aus der Umgebung
aufgenommen werden).}

Hygienisches Arbeiten und einwandfreie, desinfizierte
Gerätschaften z.\,B. Stethoskope zählen selbstverständlich
auch zu den elementaren Voraussetzungen einer sorgfältigen
Krankenuntersuchung.
}{}{}{}{}

\Text{Kontakt mit dem Patienten}{%
Wichtig ist der höfliche, taktvolle Umgang mit dem Patienten
und die Konzentration auf die vom Patienten angegebenen
Beschwerden. Ablenkungen wie z.\,B. Telefonieren oder
Gespräche mit Anderen sind zu unterlassen.
Die \EE{Kontaktaufnahme} mit einem Patienten beinhaltet
folgende Handlungen: 


\begin{itemize}
  \item {Begrüßung}
  \item {Vorstellung}
  \item Handelt es sich um den richtigen Patienten?
  (Armcode falls Patient nicht kontaktierbar)
  \item Eigentliche Kontaktaufnahme (einen Moment
  innehalten, \EE{den Patient als Person
  wahrnehmen}, den Patientenkontakt nicht auf
  die reine Handlung reduzieren)
  \item {Erklärung der Tätigkeit}
\end{itemize}
}{}{}{}{}

\Text{Beim Beenden des Patientenkontakts}{%
\begin{itemize}
  \item Patient ermutigen, Fragen zu stellen und sich
  bei Problemen und Komplikationen rückzumelden
  \item Auf allfällige weitere
  Anweisungen/Informationen nicht vergessen
  \item {Verabschiedung vom Patienten}
\end{itemize}
}{}{}{}{}


\subsection{Zum Untersuchungsablauf}

Anm.: Aus organisatorischen Gründen wird die Erhebung des
Status präsens in Teilbereiche gegliedert, in der Realität
wird der Status in seiner Gesamtheit erhoben!
Das Augenmerk wird dabei nicht auf den effizientesten
Untersuchungsablauf gelegt.

Eine \EE{sorgfältige Anamnese}, gemeinsam mit einer genauen
körperlichen Untersuchung beim Erstkontakt, weist die
Richtung der zunächst vorzunehmenden weiterführenden
Untersuchungen, ermöglicht effiziente Diagnosen und hilft
dadurch dem Arzt und dem Patienten.
Darüber hinaus werden Kosten für u.\,U. unnötige
Untersuchungen vermieden. Als Basis für eine physikalische
Krankenuntersuchung dient die genaue Kenntnis der
Verhältnisse beim Gesunden.

Zum \EE{ersten Patientenkontakt} gehört eine Erhebung der
Vitalparameter (Blutdruck, Puls-, Atemfrequenz,
Körpertemperatur). Der Untersucher macht sich ein
orientierendes Bild vom Patienten und schätzt ab, ob der
Patient akut schwer erkrankt ist und daher dringender
Handlungsbedarf besteht (Notfall). Im weiteren Verlauf der klinischen Untersuchung entscheidet der
Untersucher auch über die weitere Diagnostik (bildgebende
Verfahren, Labor, konsiliarische Begutachtung etc.).

\Text{Patientenführung}{%
Als Untersucher sollte man daran denken, dass der klinische
Status für den Patienten meist keine Routine ist, und auch
immer eine Verletzung der Privatsphäre bedeutet. Deshalb ist
es wichtig, jeden Untersuchungsgang zu erklären, und dem
Patienten klare Anweisungen zu geben. Das Wissen um die
Notwendigkeit bzw. den Ablauf der Untersuchung beugt dem
Gefühl des Kontrollverlustes beim Patienten vor.
}{}{}{}{}

\Text{Dokumentation von Befunden}{%
Bei der Dokumentation sollten keine Diagnosen, sondern
Symptome und eigene Worte zur Beschreibung der
Auffälligkeiten verwendet werden. Also nicht
\Speech{\enquote{Psoriasis oder Mitralstenose}} sondern:
\Speech{\enquote{schuppendes Exanthem bzw. Diastolikum mit punctum
maximum über der Herzspitze}}.
Bei einigen Untersuchungen wurden zu
Illustrationszwecken Beispielbefunde für Normbefunde
(\SymbolFindingNormal), häufige Varianten (\SymbolFindingVariant) 
oder häufige pathologische Befunde
(\SymbolFindingPathological) angegeben.

Beziehen Sie den Patienten bei pathologischen Befunden immer
als Informationsquelle mit ein. Eventuell wurde bereits
anderswo weiterführende Diagnostik durchgeführt, oder der
Patient kennt sich mit seiner Grunderkrankung bereits
bestens aus.
}{}{}{}{}



  
\subsection{Untersuchungstechniken}



\Text{Inspektion}{%
Der Patient muss sich zur \T{Inspektion} vollständig
freimachen (d.\,h. Entkleiden bis auf die Unterwäsche). Bei der
Inspektion wird der Patient aufmerksam beobachtet und alle
Abweichungen werden in beschreibenden, möglichst
anschaulichen Worten dokumentiert.
}{}{}{}{}


\Text{Palpation}{%
Die \T{Palpation} (palpare \Lang{lat.}: streicheln) ist die
Untersuchung mit den Händen, Betasten, Befühlen, Beurteilung
von Organgrenzen, Oberfläche, Konsistenz, Pulsationen,
Fluktuation, Temperaturunterschiede, Druckschmerz und
Fremitus. Je nach Fragestellung kann man oberflächlich oder
tief palpieren, bzw. eine oder beide Hände gebrauchen.

 }{}{}{}{}


\Text{Perkussion}{%
\label{topic:perkussion}
Bei der \T{Perkussion} (percussio \Lang{lat.}: das Klopfen)
versucht man die Gewebe des Patienten durch Beklopfen zu erfassen,
wobei der hervorgerufene Klopfschall beurteilt wird. Als
\enquote{Perkussionshammer} dient der Mittelfinger der dominanten
Hand.

Dieser klopft senkrecht auf das Endgelenk des Mittelfingers
der anderen Hand, und zwar mit kurzem federndem Schlag aus
dem lockeren Handgelenk. Wichtig ist auch der Anpressdruck
des aufliegenden Fingers. Die Stärke des Klopfens hat einen
direkten Einfluss auf die Tiefe der Perkussion und damit der
wahrzunehmenden Strukturen.


\begin{PictureSeries}{}{Korrekte Hand- und Fingerhaltung bei der
Perkussion}
\PictureSeriesRow{2}
{./images/pku/clean-img69.jpg}
{}
{./images/pku/clean-img70.jpg}
{}
{}
{}
{}
{}
\end{PictureSeries}


Je stärker geklopft wird, desto weniger sind kleinere Veränderungen
(dünner Erguss, kleinere Infiltrationen) zu erfassen. Stärkeres
Klopfen dringt aber tiefer ins Gewebe ein.

Bei pathologischen Veränderungen der Lunge ist es bei (seiten-)
vergleichender Perkussion manchmal nicht leicht festzustellen ob
Unterschiede durch hypersonoren Klopfschall der einen oder Dämpfung
der anderen Seite verursacht sind.

Man unterscheidet:
\begin{itemize}
  \item \EE{Abgrenzende Perkussion}: Dabei hält man den
  aufliegenden Finger -- nur die Endphalanx -- parallel zu
  der zu erwartenden Organgrenze und beklopft das Endgelenk; wichtig zur genauen
  Erfassung einer Schallgrenze (Grenze zwischen lufthaltigem und nicht
  lufthaltigem Gewebe)
  \item \EE{Vergleichende Perkussion}: (Erfassung von
  Schallunterschieden): z.\,B. rechts/links
  
  Im Gegensatz zur abgrenzenden Perkussion wird hierbei der gesamte
  Mittelfinger aufgelegt (größeres Schallareal) und die
  Mittelphalanx beklopft.
\end{itemize}
\subparagraph{Schallqualitäten, Klopfschall (KS)}
\index{Klopfschall}
\begin{itemize}
  \item \T{Schenkelschall}: leise (kleine Amplitude), hell
  (höher frequent), kurz über nicht lufthaltigen
  Organen; z.\,B. Leber und großen Ergüssen
  \item \T{Sonorer Klopfschall}: über lufthaltigen Organen
  (Lungenschall) laut (große
  Amplitude), tief (niedere Frequenz), lang
  \item \T{Hypersonorer Klopfschall}: bei Emphysem,
  Pneumothorax
  \item \T{Tympanitischer Klopfschall}: (niederfrequent,
  eher rhythmische Schwingungen), lang, laut über gespannten,
  luftgefüllten Hohlorganen z.\,B. auch beim Gesunden über
  dem Magenfundus, aber auch über dem Areal eines
  Pneumothorax / Spannungspneumothorax
  (Schachtelton)
  \item \T{Dämpfung}: des normalen sonoren Schalls
  (gedämpfter Klopfschall) bei Gewebsverdichtungen:
  \begin{itemize}
    \item z.\,B. bei \I{Pneumonie}
  (pneumonische Infiltrationen werden nur erreicht, wenn
  sie nicht tiefer als etwa 5\CM* unter der Oberfläche
  liegen und eine gewisse Ausdehnung erreichen)
  \item z.\,B. bei
  \I{Pleuraerguss} und \I{Pleuraschwarte} (bei Erguss steigt
  die dorsale Obergrenze der Dämpfung von medial nach lateral hin
  an (\I{Ellis-Damoiseau'sche Linie}))
  \item z.\,B. bei \I{Atelektase} (bei verringertem Luftgehalt des
  Lungengewebes)
  \end{itemize}

\end{itemize}



% \begin{center}
% \includegraphics[width=0.5\linewidth]{./images/pku/clean-img2.jpg}
% \end{center}
% 
% \begin{description}
%   \item[Abgrenzende Perkussion] dabei hält man den
%   aufliegenden Finger -- nur die Endphalanx - parallel zu der zu
%   erwartenden Organgrenze und beklopft das Endgelenk; wichtig zur genauen
%   Erfassung einer Schallgrenze (Grenze zwischen lufthaltigem und nicht
%   lufthaltigem Gewebe)
%   \item[Vergleichende Perkussion] (Erfassung von
%   Schallunterschieden): z.\,B. rechts/links
%   
%   Im Gegensatz zur abgrenzenden Perkussion wird hierbei der gesamte
%   Mittelfinger aufgelegt (größeres Schallareal) und die
%   Mittelphalanx beklopft.
% \end{description}

% 
% \subparagraph{Schallqualitäten}
% \begin{itemize}
%   \item \T{Schenkelschall}: dumpfer, kurzer Schall wie bei
%   Perkussion des Oberschenkels, Leber
%   \item \T{Sonorer Klopfschall} (KS): klingender, tiefer
%   Schall, wie über gesundem Lungengewebe
%   \item H\T{ypersonorer KS}: tiefer, lauter, länger als
%   sonor, wie bei überblähter Lunge
%   \item \T{tympanitischer KS}: trommelartig über gespannten,
%   luftgefüllten Hohlräumen, z.\,B. Magen
%   \item \T{Dämpfung}: sonorer KS wird leiser, dumpfer z.\,B.
%   Flüssigkeitansammlung, Gewebevermehrung
% \end{itemize}

}{}{}{}{}




\Text{Auskultation}{%
Während der \T{Auskultation}
(auscultare \Lang{lat.}: horchen) werden mit Hilfe eines Stethoskops
Geräusche gehört und analysiert. Die meisten Stethoskope sind mit
einem Schalltrichter für nieder- bis mittelfrequente und einem
Membrankopf für höherfrequente Schwingungen ausgestattet, die
abwechselnd verwendet werden können. Andere Stethoskope
verfügen über eine Membran, bei der mittels des ausgeübten
Drucks selektiv höhere oder tiefere Frequenzen verstärkt
werden können. 

Um brauchbare Ergebnisse zu erzielen, müssen die Ohroliven
gut angepasst sein und Schalltrichter bzw. Membran gut
aufliegen. Die \EE{Ohrbügel} sollen dem Verlauf des
Gehörgangs entsprechend \EE{nach vorne} zeigen. Weiters
sollen während der Auskultation die Schläuche nicht berührt
werden.
}{}{}{}{}



Es ist zu empfehlen, sich bei der physikalischen
Untersuchung ein Schema zurechtzulegen, und sich
systematisch vorzuarbeiten. Man sollte sich vom Scheitel zu
den Sohlen durcharbeiten, meist folgt man der Reihenfolge:
Inspektion -- Perkussion -- Palpation -- Auskultation.


%%%%%%%%%%%%%%%%%%%%%%%%%%%%%%%%%%%%%%%%%%%%%%%%%%%%%%%%%%%
%%%%%%%%%%%%%%%%%%%%%%%%%%%%%%%%%%%%%%%%%%%%%%%%%%%%%%%%%%%
%%%%%%%%%%%%%%%%%%%%%%%%%%%%%%%%%%%%%%%%%%%%%%%%%%%%%%%%%%%
%%%%%%%%%%%%%%%%%%%%%%%%%%%%%%%%%%%%%%%%%%%%%%%%%%%%%%%%%%%
%%%%%%%%%%%%%%%%%%%%%%%%%%%%%%%%%%%%%%%%%%%%%%%%%%%%%%%%%%%


\section{Allgemeinstatus}

\Text{Allgemeinzustand, Ernährungszustand}{%
Die einleitende Inspektion besteht aus der Erhebung

\begin{itemize}
  \item des \EE{Allgemeinzustands}
  (gut, krank, (sehr) reduziert, im Schock, Konstitutionstyp),
  \item des \EE{Ernährungszustands}
  (kachektisch, normal, übergewichtig, adipös),
  \item der
  \EE{Bewusstseinslage}
  (wach, somnolent,
  soporös, komatös; zeitl./räuml./situativ/zur Person
  orientiert)
  \item sowie einer
  \EE{Begutachtung}
  des gesamten Organs Haut (siehe Blickdiagnosen).
\end{itemize}

}{%
}{}{}{}



\Text{Allgemeinzustand (AZ)}
{Wichtig
ist der erste Eindruck des Patienten! 
Achte auf:

\begin{itemize}
  \item Global: normaler, reduzierter, sehr reduzierter
  Allgemeinzustand
  \item Tonus, Körperhaltung
  \item psychisch: normal, nervös, verwirrt, eingeengt
  \item Mimik: z.\,B. starr bei M. Parkinson, 
  \item Sprache: normal, heiser, verwaschen
  \item Konstitutionstyp: athletisch, pyknisch, leptosom
  \item Bewusstseinslage: 
%   \begin{multicols}{2}
  \begin{itemize}
    \item Quantitativ: 
    \begin{compactitem}
      \item wach
    \item \enquote{verlangsamt}
    \item agitiert
    \item somnolent\footnote{\I{Somnolenz}: Patient wacht
    durch leichte Reize auf}, soporös\footnote{\I{Sopor}:
    Patient benötigt starke Reize (Schmerzreiz) um zu
    erwachen, bzw. klart auch durch starke Reize nicht
    vollständig auf.}
    \item bewusstlos, komatös
    \end{compactitem}%\columnbreak
    \item Qualitativ: Orientierung zu 
    \begin{compactitem}
      \item Person
      \item Ort
      \item Zeit
      \item Situation
    \end{compactitem}
  \end{itemize}
%   \end{multicols}
\end{itemize}
}
{}
{}
{}
{}




\Text{Ernährungszustand (EZ)}{%

\begin{itemize}
  \item {normal, reduziert, über-, untergewichtig, kachektisch, adipös}
  \item Größe/Gewicht: BMI\,=\,Gewicht /
  Körperlänge\textsuperscript{2}
  [\MetricUnit{\nicefrac{kg}{m\textsuperscript{2}}}]
  \begin{itemize}
    \item Untergewicht: {\textless} 18
    \item Normalgewicht: \VonBis{18}{25}
    \item Übergewicht: \VonBis{26}{29}
    \item Adipositas: \VonBis{30}{40}
    \item Adipositas permagna: {\textgreater} 40
  \end{itemize}
\end{itemize}

{AZ und EZ sollen, wenn nicht normal, differenziert
beschrieben werden.} }{}{}{}{}


\Text{Vitalfunktionen}{%
(Erwachsene)

\begin{itemize}
  \item \EE{Bewusstsein}
  \item \EE{Atemwege}
  \begin{itemize}
    \item {Stridor: hörbares (pfeifendes) Geräusch bei
    behinderter Einatmung durch Verengung der oberen Luftwege}
  \end{itemize}
  \item \EE{Atmung}
  \begin{itemize}
    \item {normale Atemfrequenz: \VonBis{12}{16}\PerMin*
    (orthopnoisch)}
    \begin{compactitem}
      \item {Tachypnoe: {\textgreater} 16\PerMin*}
      \item {Bradypnoe: {\textless} 12\PerMin*}
    \end{compactitem}
    \item Zeichen der Atemnot (Dyspnoe):
    \begin{itemize}
      \item Patient ringt um Luft
      \item Pausen beim Sprechen notwendig (Sprechdyspnoe)
      \item Zyanose
    \end{itemize}
  \end{itemize}
  \item \EE{Kreislauf}
  \begin{itemize}
    \item Puls: normal \VonBis{60}{100}\PerMin*
    
    Abweichungen:
    \begin{compactitem}
      \item {Tachykardie {\textgreater} 100\PerMin*}
      \item {Bradykardie {\textless} 60\PerMin*}
      \item {rhythmisch, arrhythmisch (Extrasystolen, Vorhofflimmern)}
    \end{compactitem}
    \item Blutdruck: Messung bei Statuserhebung immer notwendig!
  \end{itemize}
\end{itemize}
}{}{}{}{}



\section{Kopf, Hals}

\Text{Setting}
{Sofern nicht anders angegeben, werden die Untersuchungen im
Sitzen durchgeführt. Der Patient und der Untersucher sitzen
sich gegenüber, die Beine beider Personen sind auf die
jeweils linke Seite gedreht.

}
{}
{}
{}
{}

\Text{Inspektion des Gesichtshautzustandes}{%
Hier können Phänomene wie Basaliome, Spider naevi,
Xanthelasmen, Ödeme und v.\,a.\,m. beobachtet werden. Siehe
Blickdiagnosen!
\bigskip

\bigskip
}{%
}{}{}{}


\Text{Palpation der Nasennebenhöhlen und des Mastoids}{%
Eine Entzündung der \I{Nasennebenhöhlen} kann
sich durch Schmerzen beim Beklopfen äußern (klopfdolent). Bei der
Untersuchung klopft der Untersucher mit seinen Fingerspitzen über die
Regionen der \I[Sinus!frontales]{Sinus frontales} und 
\I[Sinus!maxillares]{Sinus maxillares} (Stirn- und
Kiefer-Höhlen), sowie der \I{Prozessus mastoidei}.

}{%
}{}{}{}

\begin{PictureSeriesWide}{}{Palpation der Nasennebenhöhlen und des Mastoids}
\PictureSeriesRowWide{3}
{./images/pku/gabriel-sebastian/Nebenhoehlen-Palp-001}
{}
{./images/pku/gabriel-sebastian/Nebenhoehlen-Palp-002}
{}
{./images/pku/gabriel-sebastian/Mastoid}
{}
{}
{}
\end{PictureSeriesWide}

\Text{Trigeminus-Austrittspunkte}
{Die drei Hauptäste des \T[Nervus!trigeminus]{N.~trigeminus}
(V\textsubscript{1}: \I[Nervus!ophthalmicus]{N.~ophthalmicus}, 
V\textsubscript{2}: \I[Nervus!maxillaris]{N.~maxillaris},
V\textsubscript{3}: \I[Nervus!mandibularis]{N.~mandibularis}) 
treten an drei verschiedenen Stellen aus dem
Schädelknochen, nämlich durch das
\I[Foramen!supraorbitale]{Foramen supraorbitale}, das
\I[Foramen!infraorbitale]{Foramen infraorbitale} und das
\I[Foramen!mentale]{Foramen mentale}. Diese sollten im
Rahmen der \EE{Palpation} durch leichten Druck (Daumen) auf
Schmerzhaftigkeit untersucht werden. Die Austrittspunkte
liegen annähernd auf einer gedachten vertikalen Linie
untereinander.} 
{}
{}
{}
{}

\begin{PictureSeriesWide}{}{}
\PictureSeriesRowWide{3}
{./images/pku/gabriel-sebastian/Trigeminus-V1.jpg}
{Innerhalb der Augenbrauen 
(Foramen supraorbitale; V\textsubscript{1}:
N.~ophthalmicus ) }
{./images/pku/gabriel-sebastian/Trigeminus-V2.jpg}
{Unterrand der Orbita 
(Foramen infraorbitale; V\textsubscript{2}:
N.~maxillaris) }
{./images/pku/gabriel-sebastian/Trigeminus-V3.jpg}
{seitlich des Kinn 
(Foramen mentale; V\textsubscript{3}:
N.~mandibularis)}
{}
{}
\end{PictureSeriesWide}



\TextSlide{Mimik (N. facialis): }{%
Bei der Prüfung des \I[Nervus!facialis]{N. facialis}
(\I[VII. Hirnnerv]{VII}\index{Hirnnerv!VII}) wird besonders
auf den Seitenvergleich geachtet. Bei Vorliegen einer Parese, hat die
Unterscheidung, ob es sich um eine \EE{zentrale oder
periphere Störung} handelt, einen großen Stellenwert. Bei
Vorliegen einer zentralen Störung ist die Funktion des
Versorgungsgebiete der oberen Äste (Stirn) weitgehend
intakt, d.\,h. der Patient kann die Stirn
runzeln.\footnote{Stirnrunzeln bei zentraler
Facialisparese: Die oberen Versorgungsgebiete werden von
beiden Seiten versorgt, sodass bei einem zentralen Ausfall
einer Seite die andere weitgehend kompensieren kann: 
Z.\,B. kommt es bei Schädigung der \E{linken} morotischen
Großhirnrinde zu einer Parese der \E{rechten}
Gesichtshälfte mit Aussparung der Stirnmuskulatur und des Lidschlusses.}
Man bittet den
Patienten
\ldots

}{%
\begin{PkuBefunde}{}
\item[\SymbolPkuNormal] \PkuBefund{Seitengleich, keine Ausfälle}
\item[\SymbolFindingPathological] Seitenungleichheit, Parese
\end{PkuBefunde}
}{}{}{}

\begin{PictureSeriesWide}{}{N. facialis}
\PictureSeriesRowWide{3}
{./images/pku/clean-img9.jpg}
{{\ldots} die Augenbrauen hochzuziehen und die Stirne zu
runzeln, }
{./images/pku/clean-img10.jpg}
{{\ldots} die Augen zusammen zu kneifen,}
{./images/pku/clean-img11.jpg}
{{\ldots} die Zähne zu zeigen,}
{}
{}
\PictureSeriesRowWide{3}
{./images/pku/clean-img12.jpg}
{{\ldots} die Backen aufzublasen,}
{./images/pku/clean-img13.jpg}
{{\ldots} zu Pfeifen oder eine Kerze auszublasen.}
{empty}
{Freiraum!}
{}
{}
\end{PictureSeriesWide}



\Text{Begutachtung der Augen}{%
\begin{itemize}
  \item Lider: Lidödem, Lideinlagerungen (Xanthelasmen)
  \item Konjunktiven: blass, gerötet, blaurot, Blutungen, Verfärbungen
  \item Skleren: Gefäßzeichnung, Gelbfärbung (ikterisch)
\end{itemize}

Siehe Blickdiagnosen!
}{%
}{}{}{}


\TextSlide{Grobe Augenmotorik}{%
Es wird hierbei die Funktion der Hirnnerven \I{III}, \I{IV}
und \I{VI} untersucht 
(\I[Nervus!oculomotorius]{N.~oculomotorius}\index{Hirnnerv!III},
\I[Nervus!trochlearis]{N.~trochlearis}\index{Hirnnerv!IV},
\I[Nervus!abducens]{N.~abducens}\index{Hirnnerv!Vi}). 
Der Untersucher sitzt dem Patienten gegenüber, der Patient
folgt mit seinen Augen dem Zeigefinger des Untersucheres
(WICHTIG: keine Mitbewegung des Kopfes!). Der Finger des
Untersuchers wird seitlich sowie nach oben und unten bewegt.
Die Augen des Patienten sollten in einer gleichmäßigen,
flüssigen Bewegung den Zeigefinger verfolgen.

Ein Fehlen derselben kann auf eine Störung der Augenmuskulatur, der
sie versorgenden Nerven oder auf zentrale Störungen
zurückzuführen sein. Bei Lähmung eines Augenmuskels bleibt der
Bulbus zurück und es resultieren Doppelbilder in der Blickrichtung
des Funktionsausfalles des betroffenen Augenmuskels. Die
häufigste Lähmung ist die \I{Abduzensparese}: Der Bulbus
kann hierbei nicht mehr nach lateral bewegt werden, bzw. es
kommt zum Einwärtsschielen.

}{%
\begin{PkuBefunde}[]
  \item[\SymbolPkuNormal] \PkuBefund{Augenmotorik intakt}
%   \item[\SymbolPkuVariante] \PkuBefund{}
%   \item[\SymbolFindingPathological] \PkuBefund{}
\end{PkuBefunde}
}{}{}{}


\begin{PictureSeriesWide}{}{Augenmotorik}
\PictureSeriesRowWide{3}
{./images/pku/clean-img14.jpg}
{}
{./images/pku/clean-img15.jpg}
{}
{./images/pku/clean-img16.jpg}
{}
{}
{}
\PictureSeriesRowWide{3}
{./images/pku/clean-img17.jpg}
{}
{./images/pku/clean-img18.jpg}
{}
{./images/pku/clean-img19.jpg}
{}
{}
{}
\PictureSeriesRowWide{3}
{./images/pku/clean-img20.jpg}
{}
{}
{}
{}
{}
{}
{}
\end{PictureSeriesWide}



\begin{table}[H]
\Captionof{table}{Bsp.: Nystagmus und Adduktionsschwäche}
\FormatTable
\includegraphics[width=\linewidth]{./images/pku/Nystagmus}

% \begin{flushleft}

\begin{tabularx}{\linewidth}{XXXX}
\centering  Nystagmus des rechten Auges
beim Blick nach rechts
&
\centering  Adduktions-schwäche beim
Blick nach rechts
&
\centering  Adduktions-schwäche beim
Blick nach links
&
\centering Nystagmus des
linken Auges beim Blick nach links
\\
\end{tabularx}
\end{table}




% \end{flushleft}

Die \II{Naheinstellung} der Augen sollte geprüft
werden, indem der Zeigefinger des Untersuchers von der Ferne bis auf
ca. 10\CM* an den Patienten
angenähert wird und die Augen des Patienten eine Konvergenzbewegung
durchmachen, die von einer Pupillenverengung begleitet wird.

\begin{PictureSeries}{}{}
\PictureSeriesRow{2}
{./images/pku/clean-img22.jpg}
{}
{./images/pku/clean-img23.jpg}
{}
{}
{}
{}
{}
\end{PictureSeries}



\TextSlide{Gesichtsfelduntersuchung (Fingerperimetrie)}
{Mit dem \I{Konfrontationsversuch} (\I{Fingerperimetrie})
ist eine grobe Beurteilung des Gesichtsfeldes möglich. Die
Methode ist allerdings sehr grob und nur bei
\I{Quadrantenausfall} oder \I{Hemianopsie} diagnostisch
verwertbar. Wegen der großen diagnostischen und
prognostischen Bedeutung von Gesichtsfeldausfällen sollte
schon bei Verdacht auf Ausfälle eine fachärztliche
Begutachtung durchgeführt werden.

\EE{Durchführung}: Es werden die Augen einzeln untersucht, 
das jeweils andere Auge wird abgedeckt.
Der Patient fixiert die Nase des vor ihm
sitzenden Untersuchers. Der Untersucher bringt seine Hand zunächst in
den Raum außerhalb des Gesichtsfeldes des Patienten (z.\,B.
seitlich neben dessen Kopf), dann wird die Hand mit sich
langsam bewegenden Fingern kontinuierlich an die
Gesichtsfeldgrenzen und weiter bis in das Gesichtsfeld des
Patienten geführt. Der Patient gibt an, ab welcher Position
er die sich bewegenden Finger des Untersuchers
wahrnimmt.~\cite{Zeiler2001}

Die Prüfung des rechten und linken Gesichtsfeldes in
mittlerer Höhe reicht zur Beurteilung nicht aus, da auf
diese Weise allfällige Quadrantenanopsien leicht übersehen
werden können; es ist erforderlich, alle vier Quadranten der
Gesichtsfelder einzeln zu überprüfen (horizontal, vertikal,
diagonal, jeweils inklusive der nasalen Anteile).~\cite{Zeiler2001}


}{%
\begin{PkuBefunde}[]
  \item[\SymbolPkuNormal] \PkuBefund{Gesichtsfeld fingerperimetrisch unauffällig}
%   \item[\SymbolPkuVariante] \PkuBefund{}
  \item[\SymbolFindingPathological] \PkuBefund{Skotom {\ldots} / 
  Hemianopsie {\ldots} }
\end{PkuBefunde}
}{}{}{}


% \begin{table}[H]
% \begin{tabularx}{\linewidth}{XX}
% \includegraphics[width=\linewidth]{./images/pku/clean-img24.jpg}
% &
% \includegraphics[width=\linewidth]{./images/pku/clean-img25.jpg}
% \\
% \end{tabularx}
% \end{table}


\TextSlide{Pupillen und Pupillenreflex}{%
\E{Vor} Prüfung des Pupillenreflexes muss die
\EE{Beschaffenheit} der Pupillen beurteilt werden:


\begin{itemize}
  \item gleichweit (\I{isokor}),
  \item mittelweit,
  \item rund vs. entrundet,
  \item weit (\I{Mydriasis}) oder
  \item eng gestellt (\I{Miosis}; Medikamenten- oder
  Drogengebrauch?)
\end{itemize}

Den \II{Pupillenreflex}
prüft man am besten mit Hilfe einer Taschenlampe.
Hierbei fixiert der Patient einen Punkt in der Ferne. 
Der Untersucher schirmt durch seine am Nasenrücken des Patienten liegende
nicht dominante Hand die Augen des Patienten von einander
ab. Er betrachtet dann abwechselnd die Lichtreaktion der
ipsilateralen und kontralateralen Pupille (\I{swinging
flashlight}).

Beim Beleuchten eines Auges sollte
eine konsensuelle Reaktion sowohl in der beleuchteten, als
auch in der gegenüberliegenden Pupille 
hervorgerufen werden. Beurteilt wird die Promptheit der
Reaktion (prompt, verlangsamt, fehlend) und die
Gleichseitigkeit (\I{isokor}, \I{anisokor}). 
Liegt eine anisokore Reaktion vor, so ist eine genau
Beschreibung wichtig, um eine motorische Störung von einer
sensorischen abgrenzen zu können.



}{\begin{PkuBefunde}[]
  \item[\SymbolPkuNormal] \PkuBefund{Pupillen rund,
  mittelweit, isokor, prompte konsensuelle Lichtreaktion}
  \item[\SymbolFindingPathological] 
  \PkuBefund{rund, anisokor, rechts mittelweit, links
  geweitet, fehlende motorische Lichtreaktion links,
  Lichtsensorik bds. intakt}
\end{PkuBefunde}}{}{}{}





\begin{PictureSeriesWide}{}{}
\PictureSeriesRowWide{3}
{./images/pku/clean-img26.jpg}
{}
{./images/pku/clean-img27.jpg}
{}
{./images/pku/clean-img28.jpg}
{}
{}
{}
\end{PictureSeriesWide}


\TextSlide{Inspektion des Mundes}{%
Die Inspektion erfolgt unter Zuhilfenahme einer Lampe und
eines Spatels.\footnote{\EE{Spatel nicht fixieren}: Der
Spatel soll mittels Daumen und Zeigefinger so gehalten werden, dass er nicht
komplett fixiert ist, sondern bei Anstossen an
Strukturen nach hinten ausweichen kann.}\footnote{Der Mund
trocknet rasch aus, bei längerer Untersuchung soll der
Patient zwischendurch den Mund schließen.}

\begin{itemize}
  \item Farbe der Lippen und der \I{Gingiva}
  (Mundschleimhaut)
  \item Veränderungen der Schleimhaut wie Bläschen,
  \I{Ulzerationen} (Geschwüre) oder \I{Rhagaden}
  (spaltförmige Einrisse der Haut)
  \item Rötung oder Schwellung der Schleimhäute als Zeichen einer
  Entzündung
  \item Foetor ex ore?
\end{itemize}

Der prinzipielle Aspekt der Mundhöhle und
die Begutachtung der Vollständigkeit und des Zustandes des Gebisses
stellen weiters einen wichtigen Teil der Untersuchung dar, und geben
außerdem Indizien auf die Hygiene des Patienten
(Gebiss saniert, vollständig, sanierungsbedürftig, teilweise
erhalten, Prothese). 

}{%
\begin{PkuBefunde}[]
  \item[\SymbolPkuNormal] \PkuBefund{Mundschleimhaut bland,
  Gebiss saniert}
  \item[\SymbolPkuVariante] \PkuBefund{Zahnteilprothese
  Oberkiefer / Zahnvollprothese Unterkiefer / Gebiss
  lückenhaft}
  \item[\SymbolFindingPathological] 
  \PkuBefund{Mundschleimhaut: weißliche Stippchen, Gebiss kariös,
  sanierungsbedürftig}
\end{PkuBefunde}}{}{}{}




\TextSlide{Zunge}{%
\begin{itemize}
  \item Größe
  \item Farbe
  \item Oberfläche (pathologischer Belag (z.\,B. im Rahmen einer Pilzinfektion),
  Exsikkose) 
  \item \EE{Zungenmotorik}: 
  Der Untersucher bittet den Patienten, die Zunge gerade, dann nach rechts
und links hinauszustrecken. Eine ungewollte Abweichung der Zunge nach
einer Seite weist auf eine Störung des ipsilateralen (gleichseitigen)
XII. Hirnnervs (\I[Nervus!hypoglossus]{N. hypoglossus}) hin;
d.\,h. \EE{die Zunge weicht zur Seite der Läsion hin}.
\end{itemize}


}{%
\begin{PkuBefunde}[]
  \item[\SymbolPkuNormal] \PkuBefund{Zunge feucht, ohne
  Belag, Motorik seitengleich}
  \item[\SymbolPkuVariante] \PkuBefund{Zunge weißlich
  belegt}
  \item[\SymbolFindingPathological] \PkuBefund{Zunge
  trocken und borkig mit gelblichem Belag / Zunge weicht
  nach \ldots ab / Zunge atroph}
\end{PkuBefunde}}{}{}{}






\TextSlide{Untersuchung des Pharynx und der Tonsillen}{%
\index{Uvula}%
\index{Gaumenzäpfchen}%
Vorgehensweise: Mit dem Spatel fixieren Sie die Zunge, der Patient soll
\Speech{\enquote{Aaaaaaah!}} sagen -- dabei achten Sie auf
die Uvulastellung (Gaumenzäpfchen). Eine Seitenabweichung weist auf
eine Läsion des \I[Nervus!glossopharyngeus]{N.
glossopharyngeus} hin.

\begin{multicols}{2}
\includegraphics[width=\linewidth]{./images/pku/clean-img29.jpg}

\includegraphics[width=\linewidth]{./images/pku/clean-img30.jpg}
\end{multicols}
}{%
\begin{PkuBefunde}[]
  \item[\SymbolPkuNormal] \PkuBefund{Pharynx bland,
  Tonsillen bland}
  \item[\SymbolPkuVariante] \PkuBefund{Tonsillen hypertroph}
  \item[\SymbolFindingPathological] \PkuBefund{Pharynx gerötet
  / Eiterstraße in Pharynx / Tonsillen mit
  weißlich-gelblichen Stippchen belegt}
\end{PkuBefunde}}{}{}{}


 


\TextSlide{Lymphknoten}
{Die \EE{Palpation} erfolgt \EE{von hinten} mit warmen
Händen, beidhändig, gleichzeitig rechts und links
vergleichend. Der Patient neigt den Kopf leicht vor.
Man beurteilt die Lymphknoten hinsichtlich:

\begin{itemize}
  \item Lokalisation
  \item Verschieblichkeit
  \item Konsistenz
  \item Größe
  \item Schmerzhaftigkeit
\end{itemize}

\noindent Eine Vergrößerung der Lymphknoten kann
auf Entzündungen oder tumoröse Prozesse hinweisen:

\begin{compactitem}
  \item Akut entzündete Lymphknoten: meist weich, schmerzhaft und
  verschieblich
  \item \II{Lymphknotenmetastasen} oft hart, indolent und
  \EE{unverschieblich}
\end{compactitem}
}{%
\begin{PkuBefunde}[]
  \item[\SymbolPkuNormal] \PkuBefund{Halslymphknoten nicht
  palpabel}
%   \item[\SymbolPkuVariante] \PkuBefund{}
  \item[\SymbolFindingPathological]
  \PkuBefund{3\,$\times$\,2\CM* großer, verschieblicher,
  indolenter Lymphknoten jugulär}
\end{PkuBefunde}
}{}{}{}



\bigskip

Es werden im Bereich Kopf-Hals routinemäßig folgende Lymphknoten
palpiert:

\begin{WidePar}
\noindent\Parallel
{\begin{enumerate}
  \item[{\sffamily\textcircled{\scriptsize 1}}] \EE{Präauriculär}
  \item[{\sffamily\textcircled{\scriptsize 2}}] \EE{Retroauriculär}
  \item[{\sffamily\textcircled{\scriptsize 3}}] \EE{Submandibulär}
  \item[{\sffamily\textcircled{\scriptsize 4}}] \EE{Submental}
  \item[{\sffamily\textcircled{\scriptsize 5}}] \EE{Prä-/paralaryngeal}
  \item[{\sffamily\textcircled{\scriptsize 6}}] \EE{Jugulär}
  \item[{\sffamily\textcircled{\scriptsize 7}}] \EE{Nuchal}
  \item[{\sffamily\textcircled{\scriptsize 8}}] \EE{Supraclaviculär}
\end{enumerate}
\Captionof{figure}{Schema der Halslymphknoten.
\SourcePicture{Piero Lercher, aus: Thurnher et al.:
\emph{HNO-Heilkunde: Ein symptomorientiertes Lehrbuch.} \cite{ThurnherHNO1}}} } 
{\includegraphics[width=\linewidth]{./images/pku/Lymphknoten-Schema-Grasl}}
\end{WidePar}






\TextSlide{Die Schilddrüse}{%
\subparagraph{Inspektion} Ein ausgeprägter Kropf
(\I{Struma}) ist eventuell schon bei \EE{zurück geneigtem Kopf} an der Asymmetrie des Halses zu
erkennen. Die Schilddrüse ist bei Normalbefund in Ruhe \EE{nicht
sichtbar}. 

\subparagraph{Bimanuelle Palpation} Bei der bimanuellen
Schilddrüsenpalpation steht der Untersucher hinter dem
sitzenden Patienten. Der Patient hält den Kopf gerade und
neigt ihn leicht vor, die Finger palpieren suchend vom
Ringknorpel nach lateral. Dann legt der Untersucher die Finger flach auf. Der Patient wird gebeten zu schlucken
(eventuell mit Hilfe eines Glases Wasser), die Schilddrüse gleitet
spürbar unter den Fingern des Untersuchers vorbei. 

Pathologische Befunde sind im Gegensatz zur regelrechten SD relativ
leicht zu tasten (z.\,B. Knoten, Vergrößerung, Asymmetrie, fehlende
\I{Schluckverschieblichkeit}).

}{
\begin{PkuBefunde}[]
  \item[\SymbolPkuNormal] \PkuBefund{Schilddrüse
  schluckverschieblich, nicht vergrößert}
%   \item[\SymbolPkuVariante] \PkuBefund{}
%   \item[\SymbolFindingPathological] \PkuBefund{}
\end{PkuBefunde}

}{}{}{}

\begin{PictureSeries}{}{}
\PictureSeriesRow{2}
{./images/pku/clean-img34.jpg}
{}
{./images/pku/clean-img35.jpg}
{}
{}
{}
{}
{}
\end{PictureSeries}

\begin{Literary}
\fullcite{ThurnherHNO1}


~~
\end{Literary}

% \Text{{\protect\dsliterary} Weiterführende Literatur}{%
% \begin{labeling}{mmm}
%   \item[] Thurnher, Grasl, Erovic und Lercher:
%  \emph{HNO-Heilkunde: Ein symptomorientiertes Lehrbuch.}
% \end{labeling}
% }{}{}{}{}

\cleardoublepage
\subsection{Blickdiagnosen} 

Quellen
der Abb.: nicht eruierbar


\begin{WidePar}

\begin{multicols}{3}

\noindent\begin{minipage}{\linewidth}
\includegraphics[width=\linewidth]{./images/pku/gabriel-sebastian/pupillendifferenz}

\Captionof{figure}{\I{Anisokorie} der Pupillen.
\SourcePicture{Gabriel}}
\end{minipage}






\end{multicols}

\end{WidePar}






%%%%%%%%%%%%%%%%%%%%%%%%%%%%%%%%%%%%%%%%%%%%%%%%%%%%%%%%%%%
%%%%%%%%%%%%%%%%%%%%%%%%%%%%%%%%%%%%%%%%%%%%%%%%%%%%%%%%%%%
%%%%%%%%%%%%%%%%%%%%%%%%%%%%%%%%%%%%%%%%%%%%%%%%%%%%%%%%%%%
%%%%%%%%%%%%%%%%%%%%%%%%%%%%%%%%%%%%%%%%%%%%%%%%%%%%%%%%%%%
%%%%%%%%%%%%%%%%%%%%%%%%%%%%%%%%%%%%%%%%%%%%%%%%%%%%%%%%%%%
\clearpage


\section{Wirbelsäule/Thorax/Lunge}



\TextSlide{Beurteilung der Wirbelsäule: }
{%
Beurteilt wird die Symmetrie im aufrechten Stehen und im
Herunterbeugen sowie die Klopfschmerzhaftigkeit. Die
Inspektion erfolgt dabei im \EE{Stehen}, die \EE{Schuhe
müssen ausgezogen} werden.
 
Normale, physiologische Wirbelsäulenkrümmungen:
\begin{itemize}
  \item konvexe Brustwirbelsäulenkrümmung = Brustkyphose
  \item konkave Lendenwirbelsäulenkrümmung = Lendenlordose
\end{itemize}

\noindent Krümmungen können \EE{verstärkt} (z.\,B. bei
osteoporotischer Hyperkyphose) oder \EE{vermindert} sein:

\begin{itemize}
  \item \T{Skoliose}: eine pathologische Seitenabweichung
  der Wirbelsäule liegt vor
  \item \T{Kyphose}: Rundrücken
  (Wirbelsäulenkrümmung in sagittaler Ebene dorsal-konvex)
  \item \T{Lordose}: Wirbelsäulenkrümmung
  dorsal-konkav
  \item \T{Gibbus} (Buckel bzw. Höcker z.\,B. bei
  Osteoporose mit Wirbeleinbruch)
  \item \T{Kyphoskoliose}: starke Verkrümmung der
  Wirbelsäule mit Atem- und Kreislaufbehinderung
\end{itemize}}
{
\begin{PkuBefunde}[]
  \item[\SymbolPkuNormal] \PkuBefund{Wirbelsäule
  inspekorischn unauffällig, keine Klopfdolenz}
%   \item[\SymbolPkuVariante] \PkuBefund{}
%   \item[\SymbolFindingPathological] \PkuBefund{}
\end{PkuBefunde}
}
{}
{}
{}


\TextSlide{Inspektion des Thorax (Form, Symmetrie, inkl. Haut)}
{\subparagraph{Inspektion}


\begin{itemize}
  \item Spuren von traumatisch bedingten Veränderungen (Rippenbrüche,
  Abschürfungen, etc.)
  \item anderen Auffälligkeiten
\end{itemize}

\subparagraph{Form des Thorax} 


\begin{itemize}
  \item \T{Trichterbrust} (\T{Pectus excavatum}): unteres
  Sternum eingedellt, meist angeboren,
  \item \T{Hühnerbrust}
  (\T{Pectus carinatum}): Sternum ragt wie ein Kiel hervor;
  oft angeboren oder bei \I{Rachitis} (\IX[Vitamin D]{Vitamin-D}-Mangel)
  \item \T{Fassthorax}: Oft bei Lungenemphysem
  (Lungenüberblähung), die Supraklavikulargruben sind hierbei oft vorgewölbt.
\end{itemize}




Siehe Blickdiagnosen:
\begin{itemize}
  \item \T{Blue bloater} übergewichtiger,
  zyanotischer Patient,
  \item \T{Pink puffer} eher magerer Patient mit
  rosa Haut.
\end{itemize}}
{
\begin{PkuBefunde}[]
  \item[\SymbolPkuNormal] \PkuBefund{Thorax
  symmetrisch, inspektoisch unauffällig}
%   \item[\SymbolPkuVariante] \PkuBefund{}
%   \item[\SymbolFindingPathological] \PkuBefund{}
\end{PkuBefunde}
}
{}
{}
{}






\subparagraph{Skoliose}

Klinische Skoliosezeichen bei BWS-rechtskonvexer und LWS-linkskonvexer
Skoliose:


\begin{tabularx}{\linewidth}{XX}\FormatTable
a)

Schulterhochstand rechts,
Schulterblatthochstand rechts,
abgeflachtes Taillendreieck links und
vertieftes rechts, Beckenhochstand rechts

b)

bei Vorneigung sind die Torsionssymptome 
(Rippenbuckel und Lendenwulst) 
deutlicher zu erkennen.


Quelle der Abb.: nicht eruierbar
&
\noindent\begin{minipage}[m]{\linewidth}
\includegraphics[width=\linewidth]{./images/pku/Skoliose}
\end{minipage}

\\
\end{tabularx}

\begin{PictureSeriesWide}{}{}
\PictureSeriesRowWide{3}
{./images/pku/clean-img59.jpg}
{Überprüfung des Beckengeradestandes durch gleichzeitige Palpation
der Beckenkämme.}
{./images/pku/clean-img60.jpg}
{Überprüfung möglicher Seitenabweichungen oder
Unregelmäßigkeiten der Dornfortsätze (Patient: aus dem Stand in
die Rumpfbeuge). Auch auf Asymmetrie von Hautfalten achten (Skoliose).}
{./images/pku/clean-img61.jpg}
{}
{}
{}
\PictureSeriesRowWide{3}
{./images/pku/clean-img62.jpg}
{Klopfschmerz: Überprüfung von kranial nach kaudal durch
\enquote{Perkussion} mit der Faust. CAVE:
gefühlvolles Klopfen!! (Osteoporose = Frakturneigung)}
{./images/pku/clean-img63.jpg}
{}
{empty}
{}
{}
{}
\end{PictureSeriesWide}


\Z{Der Klopfschmerz der Niere kann an dieser
Stelle mit erhoben werden.}


\clearpage
\Text{Orientierung/Orientierungslinien}
{\mbox{}

% \noindent\begin{minipage}{\linewidth}
% \includegraphics[width=\linewidth]{./images/pku/Thorax-Orientierungslinien.pdf}
% 
% \Captionof{figure}{
% }
% \end{minipage}
% 
% \noindent\begin{minipage}{\linewidth}
% \includegraphics[width=\linewidth]{./images/pku/Thorax-horizontal-Orientierungslinien.pdf}
% 
% \Captionof{figure}{
% }
% \end{minipage}
% 
% 
% 
% \noindent\begin{minipage}{\linewidth}
% \includegraphics[width=\linewidth]{./images/pku/Thorax-Ruecken-Orientierungslinien.pdf}
% 
% \Captionof{figure}{
% }
% \end{minipage}

} 
{}
{}
{}
{}


\begin{PictureSeriesWide}{}{Orientierunglinien am Thorax}
\PictureSeriesRowWide{2}
{./images/pku/Thorax-Orientierungslinien.pdf}
{\EE{Ventrale Orientierungslinien} am
Thorax: 
\\
Vordere Medianlinie bzw. Sternalinie
{\sffamily\textcircled{\scriptsize a}},
\\
Parasternallinie {\sffamily\textcircled{\scriptsize b}},
\\
Medioklavikularlinie (MCL;
{\sffamily\textcircled{\scriptsize c}}),
\\
Vordere Axillarlinie {\sffamily\textcircled{\scriptsize d}}}
{./images/pku/Thorax-horizontal-Orientierungslinien.pdf}
{\EE{Horizontale Orientierungspunkte} am
Thorax:
\\
Interkostalräume (\I{ICR}, schwarze Pfeile),
\\
Rippen (1.~tastbare Rippe unter der Clavicula (Ellipse)
ist die 2.~Rippe),
\\
Rippenbogen
}
{./images/pku/Thorax-Ruecken-Orientierungslinien.pdf}
{}
{empty}
{}
\PictureSeriesRowWide{2}
{./images/pku/Thorax-Ruecken-Orientierungslinien.pdf}
{\EE{Dorsale Orientierungslinien} am
Thorax: 
\\
Hintere Medianlinie bzw. Vertebralinie
{\sffamily\textcircled{\scriptsize a}},
\\
Paravertebrallinie {\sffamily\textcircled{\scriptsize b}} , 
\\
Scapularlinie (Ang.
inf. scapulae); {\sffamily\textcircled{\scriptsize c}},
\\
hintere Axillarlinie {\sffamily\textcircled{\scriptsize d}},
\\
Zählen von Wirbelkörpern und Dornfortsätzen bezogen auf den
Dornfortsatz des 7. HWK}
{empty}
{}
{empty}
{}
{empty}
{}
\end{PictureSeriesWide}





\TextSlide{Thoraxkompression}{%
Zur Überprüfung von Instabilität (Rippenfraktur) und
Druckdolenz wird mit beiden Händen symmetrisch der Thorax
seitlich und frontal \EE{diskret} komprimiert: auf
Schmerzreaktion des Patienten achten!
}{
\begin{PkuBefunde}[]
  \item[\SymbolPkuNormal] \PkuBefund{Kein
  Thoraxkompressionsschmerz}
%   \item[\SymbolPkuVariante] \PkuBefund{}
%   \item[\SymbolFindingPathological] \PkuBefund{}
\end{PkuBefunde}
}{}{}{}

\TextSlide{Atemexkursionen}{%
Feststellung nicht seitengleicher Atembewegung: Auflegen der
Hände frontal an der unteren Hälfte des Brustkorbes und
Beobachtung, ob sie sich beim Atmen seitengleich verschieben. 
Achten Sie auf die Bewegungen von Thorax und Abdomen.
}{
\begin{PkuBefunde}[]
  \item[\SymbolPkuNormal] \PkuBefund{Seitengleiche
  Atemexkursionen von unauffälliger Tiefe}
%   \item[\SymbolPkuVariante] \PkuBefund{}
  \item[\SymbolFindingPathological] \PkuBefund{Paradoxe Atmung}
  
  \PkuBefund{seichte Atmung / geringe Atemexkursionen}
\end{PkuBefunde}
}{}{}{}



\TextSlide{Auskultation}
{Die Auskultation wird bei tiefer Atmung durch den
\E{geöffneten Mund} durchgeführt.
Achten Sie dabei darauf, dass die Auskultation
nicht zu lange dauert, sonst wird der
Patient plötzlich schwindelig (wegen
respiratorischer Alkalose durch
Hyperventilation mit cerebraler 
Vasokonstriktion), ggfs. den Patienten zwischendurch Pausen
machen lassen.

Am Ende der Auskultation immer \EE{2\,$\times$ husten}
lassen, dabei befindet sich das
Stethoskop über den mittleren 
Lungenpartien jeweils rechts und links.
(Erst dabei wird die Menge Schleim in
den Bronchien erkannt)

\subparagraph{Auskultationspunkte}

immer im direkten Seitenvergleich auskultieren

\noindent\Parallel
{\includegraphics[width=\linewidth]{./images/pku/Auskultationspunkte}}
{\includegraphics[width=\linewidth]{./images/pku/clean-img67.jpg} 

\includegraphics[width=\linewidth]{./images/pku/clean-img68.jpg}}






\subparagraph{Orientierung/Orientierungslinien}
\begin{itemize}
  \item ventral: vordere \I{Medianlinie},
  \I[Sternallinie]{Sternal-}, \I{Parasternalinie},
  \I{Medioklavikularlinie} (\I{MCL}) rechts und links,
  vordere/mittlere \I{Axillarlinie}
  \item horizontal: \I{Interkostalräume} (\I{ICR}), Rippen
  (1.
  tastbare Rippe unter der Klavikula ist die 2. Rippe),
  Rippenbogen
  \item dorsal: hintere \I{Axillarlinie}, \I{Scapularlinie}
  (Ang.
  inf. Scapulae), \I[Vertebrallinie]{Vertebral-},
  \I{Paravertebralinie}, hintere \I{Medianlinie}; Zählen von Wirbelkörpern und
  Dornfortsätzen bezogen auf den 7. HWK und den 11. BWK
\end{itemize}

\subparagraph{Normale Atmung oder Luftnot} (Dyspnoe)


\begin{compactitem}
  \item Belastungsdyspnoe nur bei Anstrengung
  \item Ruhedyspnoe (\Ca{Notfall!})
\end{compactitem}

\subparagraph{Atemfrequenz:} (normal: \VonBis{12}{16}\PerMin*)


\begin{compactitem}
  \item[$\uparrow$] \T{Tachypnoe}
  \item[$\downarrow$] \T{Bradypnoe}
\end{compactitem}




\subparagraph{Atemtyp}


\begin{itemize}
  \item \EE{normal}: Einatmung (\I{Inspiration}) etwas
  kürzer als Ausatmung (\I{Exspiration})
  
  \begin{itemize}
    \item I\,:\,E unter 1
  \end{itemize}
  \item bei obstruktiven Erkrankungen z.\,B. Asthma, ist das Exspirium deutlich
  verlängert
  
  \begin{itemize}
    \item {I\,:\,E deutlich unter 1}
  \end{itemize}
  \item {bei inspiratorischer Dyspnoe ist dies umgekehrt}
  
  \begin{itemize}
    \item {I\,:\,E über 1}
  \end{itemize}
\end{itemize}


\subparagraph{Atemgeräusche}


\begin{itemize}
  \item \T{Vesikuläratmen}: normales Atemgeräusch,
  Inspiration ist lauter als Exspiration
  \item \T{Bronchialatmen}: wie über Trachea oder oberen
  Lungenpartien mit großen Bronchien. In- und Exspiration sind
  verschärft, z.\,B. über verdichtetem Lungengewebe
  \item \T{Trockene Rasselgeräusche} (RG): Giemen und
  Brummen, exspiratorisches Pfeifen
  \item \T{Feuchte RG}: unterteilt in
  
  \begin{itemize}
    \item \I{feinblasig} (eher Knistern)
    \item \I{grobblasig} (Blubbern)
  \end{itemize}
  \item \T{Stridor}: inspiratorisches
  (gelegentlich exspiratorisches), auf Distanz (ohne Stethoskop)
  hörbares Stenosegeräusch (Pfeifen) bei Verengung der großen
  Atemwege durch obstruktive Lungenerkrankungen, entzündliche oder
  ödematöse Prozesse (\I{Anaphylaxie}) oder Fremdkörper.
\end{itemize}
}{
\begin{PkuBefunde}[]
  \item[\SymbolPkuNormal] \PkuBefund{Vesikuläres
  Atemgeräusch beidseits ohne Atemnebengeräusche,
  normales I\,:\,E-Verhältnis, Eupnoe}
%   \item[\SymbolPkuVariante] \PkuBefund{}
  \item[\SymbolFindingPathological] \PkuBefund{Feuchte
  Rasselgeräusche bds. basal, sonst vesikuläres
  Atemgeräusch / verlängertes Exspirium / Tachypnoe}
\end{PkuBefunde}
}{}{}{}

\TextSlide{Stimmfremitus \& Bronchophonie}{%
\begin{description}
  \item[Stimmfremitus (99)] Palpatorische Erfassung feiner
  Vibrationen mit beidseits über den Lungenpartien leicht aufgelegten Fingern und Handflächen. Der Patient
  spricht mit tiefer Stimme \Speech{\enquote{99}}:
  
  Verstärkte Vibration über Infiltraten, vermindert über
  Pleuraergüssen, Pleuraschwarten und Pneumothorax.
  \item[Bronchophonie (66)] Patient flüstert
  \Speech{\enquote{66}}
  Prinzipiell deutet eine abgeschwächte oder fehlende
Bronchophonie auf lufthaltige Räume in der Lunge hin (z.B.
Pneumothorax, Lungenemphysem).
Eine verstärkte Bronchophonie findet man generell bei
Verdichtungen des Lungengewebes (z.B. Pneumonie). 
 
Bei Atelektasen ist die Bronchophonie etwas diffiziler zu
betrachten und es kann keine generelle
Aussage getätigt werden.
%   
%   \begin{itemize}
%     \item besser hörbar über verdichtetem Lungengewebe (z.\,B. Pneumonie,
%     Atelektase).
%     \item vermindert hörbar z.\,B. über Pleuraerguss oder Pleuraschwarte.
%   \end{itemize}
\end{description}




% Resorptions- oder Verschlussatelektase: Bronchophonie fehlt
% bzw. abgeschwächt 
% 
% Kompressionsatelektase von innen: 
% Bronchophonie Kompressionsatelektase von außen (z.B. durch
% Pleuraerguß): Bronchophonie fehlt bzw. abgeschwächt Generell
% können auch Atelektasen kleineren Ausmaßes zu einer
% verstärkten Bronchophonie führen.



}{
\begin{PkuBefunde}[]
  \item[\SymbolPkuNormal] \PkuBefund{Stimmfremitus und Bronchophonie unauffällig}
%   \item[\SymbolPkuVariante] \PkuBefund{}
%   \item[\SymbolFindingPathological] \PkuBefund{}
\end{PkuBefunde}
}{}{}{}

\begin{PictureSeries}{}{}
\PictureSeriesRow{2}
{./images/pku/clean-img64.jpg}
{Stimmfremitus}
{./images/pku/clean-img65.jpg}
{Bronchophonie}
{}
{}
{}
{}
\end{PictureSeries}







\TextSlide{Perkussionsuntersuchung}{%
\label{topic:perkussionsuntersuchung}
Die Technik der Perkussion wird im Detail unter
\prettyref{topic:perkussion} auf
Seite \pageref{topic:perkussion} vorgestellt.


\begin{itemize}
  \item \EE{Dorsal}:
  In der Scapularlinie verlaufen die Lungengrenzen entlang der 9. Rippe,
  und in der Vertebrallinie am 11. Brustwirbel (welcher durch den
  Dornfortsatz bestimmbar ist).
  \item \EE{Ventral (bei speziellen Fragestellungen)}:
  Man perkutiert die Lungenflügel entlang der Medioclavicularlinie (MCL,
  s.\,o.) von cranial (oben) nach caudal (unten). Die
  Lungen-Leber-Grenze (= der Übergang des lungentypischen
  sonoren Klopfschalles in die
  \enquote{Leberdämpfung}) sollte rechts in
  Höhe der 6. Rippe verlaufen.
  
  Links wird die Grenze durch die Herzdämpfung und den Übergang in den
  tympanitischen Schall des Magens überlagert.
  In der vorderen Axillarlinie verlaufen die Lungengrenzen etwa in Höhe
  der 8. Rippe.
  \item \EE{Perkussion der
  Krönig{\textquoteright}schen Schallfelder}
  (bei speziellen
  Fragestellungen):
  Man perkutiert ca. 3 QF breit in Schulterhöhe über den Lungenspitzen
  (vom Hals zum Akromion).
  Eine Erweiterung spricht für ein Emphysem, während eine Dämpfung
  in diesem Bereich auf z.\,B. eine Lungenspitzentuberkulose hinweisen
  kann.
\end{itemize}


}{\MoveDown}{}{}{}







\TextSlide{Vergleichende Perkussion}{%

}{
\MoveDown
}{}{}{}


\begin{PictureSeries}{}{}
\PictureSeriesRow{2}
{./images/pku/clean-img71.jpg}
{}
{./images/pku/clean-img72.jpg}
{}
{}
{}
{}
{}
\PictureSeriesRow{2}
{./images/pku/clean-img73.jpg}
{}
{./images/pku/clean-img74.jpg}
{}
{}
{}
{}
{}
\end{PictureSeries}



\TextSlide{Prüfung der Atemverschieblichkeit}{%
Bei der Testung der Atemverschieblichkeit perkutiert man die basalen
Grenzen am Rücken bei maximaler Inspiration (den Patienten bitten,
kurz die Luft anzuhalten) und vergleicht sie mit jenen bei maximaler
Exspiration.

Die Lungengrenzen sollten bei Inspiration in der
Vertebrallinie um ca. 3\CM*, in der Scapullarlinie und in
der hinteren Axillarlinie um ca. \VonBis{4}{6}\CM* tiefer stehen als bei Exspiration.

Der normale Befund über der Lunge ist ein beidseitiger sonorer
Klopfschall in den oben beschriebenen Grenzen und eine
Atemverschieblichkeit der Grenzen um ca. \VonBis{4}{6}\CM*
(in der Scapularlinie), was auch mit \VonBis{2}{3} QF
(Querfingern) angegeben wird.
}{
\begin{PkuBefunde}[]
  \item[\SymbolPkuNormal]
  \PkuBefund{Sonorer seitengleicher Klopfschall
  über beiden Lungenflügeln, unauffällige
  Atemverscheblichkeit}
%   \item[\SymbolPkuVariante] \PkuBefund{}
%   \item[\SymbolFindingPathological] \PkuBefund{}
\end{PkuBefunde}

}{}{}{}

\begin{PictureSeries}{}{}
\PictureSeriesRow{2}
{./images/pku/clean-img75.jpg}
{}
{./images/pku/clean-img76.jpg}
{}
{}
{}
{}
{}
\end{PictureSeries}



\bigskip
\subsubsection{Typische Auskultationsbefunde}

{\small{
\begin{WidePar}
\noindent (Selbststudium) Beispiele: I = Inspektion und
Beschwerden, P = Palpation, Pk = Perkussion, A = Auskultation,
\begin{longtable}{p{0.23\linewidth}p{0.23\linewidth}p{0.23\linewidth}p{0.23\linewidth}}
\toprule\addlinespace
\TabHeading{Inspektion} & \TabHeading{Palpation} & \TabHeading{Perkusion} & \TabHeading{Auskultation}
\\\addlinespace\midrule\addlinespace\endhead
\multicolumn{4}{c}{{\TabSub{Pneumothorax}}}
\\\addlinespace\midrule\addlinespace
Dyspnoe, Thorax-, Schulterschmerzen,
  Interkostalräume können vorgewölbt sein
&Atemexkursion abgeschwächt, Stimmfremitus abgeschwächt
&hypersonorer bis tympanitischer Klopfschal, evtl. Schachtelton
&abgeschwächte bis fehlende Atemgeräusche und Bronchophonie
%   \item[»] (Cave: Spannungspneumothorax, akut
%   lebensbedrohlicher Notfall!)
\\\addlinespace\midrule\addlinespace
\multicolumn{4}{c}{{\TabSub{Pleuritis sicca}}}
\\\addlinespace\midrule\addlinespace
evtl. Schonhaltung (Patient liegt auf kranker Seite), geringere Atemexkursionen,
& 
Atemexkursionen abgeschwächt, Nachhängen der kranken Seite
& 
verminderte Basenverschieblichkeit
& 
trockene Reibegeräusche, \enquote{Lederknarren}, Verschwinden bei Auftreten eines Ergusses
%   \item[»]
\\\addlinespace\midrule\addlinespace
\multicolumn{4}{c}{{\TabSub{Pleuraerguss}}}
\\\addlinespace\midrule\addlinespace
&
Atemexkursionen abgeschwächt, 

evtl. asymmetrisch, 

Stimmfremitus abgeschwächt
& 
Dämpfung, Ellis-Damoisau'sche Linie,
Atemverschieblichkeit aufgehoben, im Sitzen nach oben abnehmende Dämpfung
& 
Atemgeräusch abgeschwächt
\\\addlinespace\midrule\addlinespace
\multicolumn{4}{c}{{\TabSub{Pleuraschwarte}}}
\\\addlinespace\midrule\addlinespace
evtl. Einziehung des Interkostalraumes
& 
Stimmfremitus (99) vermindert, 

abgeschwächte Atemexkursion
& 
KS gedämpft
& 
abgeschwächte Atemgeräusche
\\\addlinespace\midrule\addlinespace
\multicolumn{4}{c}{{\TabSub{Lungenstauung, -ödem bei kardialer Links-Dekompensation}}}
\\\addlinespace\midrule\addlinespace
Dyspnoe und/oder Orthopnoe (Luftnot beim flachen
Liegen, die sich durch aufrechte Körperhaltung bessert)
&
& 
eher gedämpft
& 
feinblasige Rasselgeräusche über Basis, bis
  zunehmend über ganzer Lunge; bei Übergang ins Lungenödem
  grobblasige RG, evtl. sogar mit bloßem Ohr hörbar 

\\\addlinespace\midrule\addlinespace
\multicolumn{4}{c}{{\TabSub{Bronchitis}}}
\\\addlinespace\midrule\addlinespace
trockener Husten
&
&
& 
trockene Rasselgeräusche
\\\addlinespace\midrule\addlinespace
\multicolumn{4}{c}{{\TabSub{Asthma bronchiale (Anfall)}}}
\\\addlinespace\midrule\addlinespace
Dyspnoe, beschleunigte Atmung, sichtbare
  Behinderung der Exspiration, Atemhilfsmuskulatur
&
& 
(nur im Anfall) hypersonor, tiefstehende Basen
& 
verlängertes Exspirium, trockene, giemende,
  pfeifende Rasselgeräusche; Cave: \I{Silent lung} bei
  zunehmender Schwere (lebensbedrohlich!)

\\\addlinespace\midrule\addlinespace
\multicolumn{4}{c}{{\TabSub{Chronische Bronchitis (Raucher)}}}
\\\addlinespace\midrule\addlinespace
Verlauf mit intermittierenden Verschlechterungen – Dyspnoe, Husten
&
& 
Von unauffällig bis hypersonorer Klopfschall mit
  tiefstehenden Lungenbasen und verminderter
  Atemverschieblichkeit bei Übergang zum Emphysem
& 
trockene RG, verlängertes Exspirium
\\\addlinespace\midrule\addlinespace
\multicolumn{4}{c}{{\TabSub{Emphysem}}}
\\\addlinespace\midrule\addlinespace
Fassthorax, horizontaler Rippenverlauf
&
& 
hypersonorer KS, Basen tiefstehend, unverschieblich
& 
abgeschwächtes Atemgeräusch, verlängertes Exspirium, trockene Rasselgeräusche
\\\addlinespace\midrule\addlinespace
\multicolumn{4}{c}{{\TabSub{Pneumonie (bakteriell)}}}
\\\addlinespace\midrule\addlinespace
Fieber, evtl. beschleunigte Atmung, Husten, bei schweren Fällen Dyspnoe
& 
Stimmfremitus (99) verstärkt , evtl. abgeschwächte Atemexkursionen
& 
evtl. Dämpfung abhängig von Größe, Abstand zur Oberfläche (<5cm)
& 
Bronchialatmen, Knisterrasseln, feinblasige
  Rasselgeräusche, Bronchophonie (66)
\\\addlinespace\midrule\addlinespace
\multicolumn{4}{c}{{\TabSub{Atelektase}}}
\\\addlinespace\midrule\addlinespace
gelegentlich asymmetrische Atmung mit
  inspiratorischer Einziehung der Zwischenrippenräume (nur
  bei großen Defekten)
& 
Stimmfremitus abgeschwächt
& 
Evtl. Dämpfung abhängig von Größe und Abstand
zur Oberfläche (<\,5\CM*)
& 
Atemgeräusche abgeschwächt
% \\\addlinespace\midrule\addlinespace
% \multicolumn{4}{c}{{\TabSub{}}}
% \\\addlinespace\midrule\addlinespace
% 
% 
\\\addlinespace\bottomrule\addlinespace
\end{longtable}
\end{WidePar}
}
}



% \bigskip
\section{Mamma}

Vorraussetzung für die Untersuchung der Mammae ist das
Vorliegen bzw. Erheben einer spezifischen Anamnese. Wichtig
sind besonders die folgenden Punkte:

\begin{compactitem}
  \item Letzte Mammographie
  \item Bisherige Erkrankungen, Verletzungen
  oder Behandlungen der Brust
  \item Familienanamnese
\end{compactitem}

Die Untersuchung der Brust betrifft nicht nur weibliche
Patientinnen, sie ist auch bei Männern
durchzuführen.\footnote{Die Inzidenz von Brustkreps bei
Männern beträgt in Österreich ca. 1 pro 100\,000 (Quelle:
Statistik Austria, Krebsinzidenz 
nach Lokalisationen und Geschlecht, Österreich seit 198).}
Auch Schwangere und Stillende sollen untersucht werden. 

\TextSlide{Inspektion der Mammae}{%
Die Inspektion der Mammae erfolgt sitzend sowohl mit locker
herabhängenden, als auch mit hochgestreckten Armen.
Die Mamma wird zur Untersuchung eingeteilt:
\begin{itemize}
\item in 4 Quadranten 
\item die retroareoläre Region.
\end{itemize}

Inspektion auf: 
\begin{itemize}
  \item grobe Seitenunterschiede
  \item Einziehungen
  \item Konturveränderungen
  \item Hautveränderungen (z.\,B. Rötung, peau
  d'orange\footnote{\I{Peau d'orange} \Lang{franz.}:
  \I{Orangenschalenhaut}. Im Rahmen der Cellulite
  auftretende Veränderung des subkutanen Fettgewebes.
  Die Haut wirkt beim Zusammenziehen induriert, die
  Follikelöffnungen sind vergrößert, im Bereich der
  Brust evtl. \EE{Zeichen eines Mammakarzinoms}
  \cite{Pschy259}.}, Hautveränderungen nach
  Irradiatio\footnote{\I{Irradiation} \Lang{lat.}:
  Bestrahlung}, \ldots)
\end{itemize}




\noindent\Parallel
{\includegraphics[width=\linewidth]{./images/pku/clean-img90}}
{41\PC* aller Karzinome treten im oberen äußeren Quadranten
auf. Häufig werden Knoten und Resistenzen von den Patientinnen selbst
entdeckt, darum ist die Selbstkontrolle eine
wichtige Vorsorgemaßnahme.

Eine Vergrößerung der Brust (Drüsen- und/oder Fettgewebe) bei
Männern nennt man \T{Gynäkomastie}.

\\
\Captionof{figure}{Schema der Lymphknoten im Brust- und
Achselbereich. \Source{Quelle nicht eruierbar.}}
}
}{
\begin{PkuBefunde}[]
  \item[\SymbolPkuNormal] \PkuBefund{Mammae inspektorisch unauffällig}
%   \item[\SymbolPkuVariante] \PkuBefund{}
%   \item[\SymbolFindingPathological] \PkuBefund{}
\end{PkuBefunde}
}{}{}{}


\TextSlide{Palpation der Mammae}
{Bei der Untersuchung ist darauf zu achten, dass die
Brustmuskeln der Patientin entspannt sind. Im Hinblick
auf die folgende Untersuchung der Achsellymphknoten
empfiehlt sich aus hygienischen Gründen das Tragen von
Untersuchungshandschuhen.\footnote{Zugegeben, auf den Fotos trägt die Untersucherin keine Handschuhe. Die Welt bietet regelmäßig Möglichkeiten zur Prozessoptimierung.}

Die Untersuchung
erfolgt \EE{bimanuell im Uhrzeigersinn mit mittlerem,
in die Tiefe gehenden, Druck und mit nach zentral
streichenden Bewegungen}, beginnend im rechten oberen Quadranten. Zuletzt
wird gesondert \EE{retromammillär} getastet.
Durch die streichenden Bewegungen kann es zu einer
Sekretion kommen, welche pathologisch sein kann.
% Die Mamille
% wird durch vorsichtiges Auspressen auf Sekretion überprüft.

}{
\begin{PkuBefunde}[]
  \item[\SymbolPkuNormal] \PkuBefund{Keine Resistenzen palpabel}
%   \item[\SymbolPkuVariante] \PkuBefund{}
%   \item[\SymbolFindingPathological] \PkuBefund{}
\end{PkuBefunde}
}{}{}{}

\begin{PictureSeriesWide}{}{}
\PictureSeriesRowWide{3}
{./images/pku/clean-img91.jpg}
{}
{./images/pku/clean-img92.jpg}
{}
{./images/pku/clean-img93.jpg}
{}
{}
{}
\PictureSeriesRowWide{3}
{./images/pku/clean-img94.jpg}
{}
{./images/pku/clean-img95.jpg}
{}
{./images/pku/clean-img96.jpg}
{}
{}
{}
\end{PictureSeriesWide}




\TextSlide{Lymphknoten (LK)}{%
LK sind nur bei Vergrößerung tastbar. Es handelt sich also um eine
orientierende Palpation der entsprechenden Regionen. LK können
reaktiv entzündlich vergrößert sein (druckschmerzhaft, weich
und verschieblich), im Rahmen von malignen oder systemischen
Erkrankungen (Lymphom, Metastasen, Sarkoidose etc: Lymphknotenpakete,
meist verhärtet, wenig druckschmerzhaft). Siehe auch Wochenthema A.

\subparagraph{Palpation der axillären LK}

Der Patient lässt die Arme locker am Körper hängen. Der
Untersucher schiebt die Hände mit geschlossenen Fingern tief in die
Axilla und palpiert dann mit mäßigem Druck gegen die Rippen nach
kaudal. Die Untersuchung wird weiter ventral und dorsal wiederholt.
Für einen vollständigen Befund sollen, den Lymphabflusswegen
entsprechend, auch die \EE{supra-} und
\EE{infraklavikulären}, sowie die \EE{Halslymphknoten}
untersucht werden.

Oft kann man reaktiv vergrößerte, verschiebliche Lymphknoten
(infolge von Rasur) tasten, welche ohne pathologische Bedeutung sind.
Suspekt sind nicht verschiebliche, sehr große oder
verbackene Lymphknoten.

}{
\begin{PkuBefunde}[]
  \item[\SymbolPkuNormal] \PkuBefund{Axilläre Lymphknoten nicht palpabel}
%   \item[\SymbolPkuVariante] \PkuBefund{}
%   \item[\SymbolFindingPathological] \PkuBefund{}
\end{PkuBefunde}
}{}{}{}



\begin{PictureSeries}{}{Palpation der axillären Lymphknoten}
\PictureSeriesRow{2}
{./images/pku/orig/gabriel-sebastian/Achsellymphknoten-001-00800.jpg}
{}
{empty}
{}
{}
{}
{}
{}
\end{PictureSeries}









\cleardoublepage
\subsection{Blickdiagnosen} 
(Quellen der Abb.: nicht
eruierbar)

\begin{WidePar}
\begin{multicols}{3}

\noindent\begin{minipage}{\linewidth}
\includegraphics[width=\linewidth]{./images/pku/clean-img97}

\Captionof{figure}{\I{Trichterbrust}}
\end{minipage}

\noindent\begin{minipage}{\linewidth}
\includegraphics[width=\linewidth]{./images/pku/clean-img105}

\Captionof{figure}{\I{Blue bloater}}
\end{minipage}

\noindent\begin{minipage}{\linewidth}
\includegraphics[width=\linewidth]{./images/pku/clean-img106.jpg} 

\Captionof{figure}{\I{Pink puffer}}
\end{minipage}

\noindent\begin{minipage}{\linewidth}
\includegraphics[width=\linewidth]{./images/pku/clean-img98} 

\Captionof{figure}{\I{Skoliose}}
\end{minipage}

\noindent\begin{minipage}{\linewidth}
\includegraphics[width=\linewidth]{./images/pku/clean-img103.jpg}

\Captionof{figure}{\I{Mamma Carcinom}}
\end{minipage}

\noindent\begin{minipage}{\linewidth}
\includegraphics[width=\linewidth]{./images/pku/clean-img104.jpg}

\Captionof{figure}{\I{Ablatio mammae}}
\end{minipage}









\end{multicols}
\end{WidePar}


%%%%%%%%%%%%%%%%%%%%%%%%%%%%%%%%%%%%%%%%%%%%%%%%%%%%%%%%%%%
%%%%%%%%%%%%%%%%%%%%%%%%%%%%%%%%%%%%%%%%%%%%%%%%%%%%%%%%%%%
%%%%%%%%%%%%%%%%%%%%%%%%%%%%%%%%%%%%%%%%%%%%%%%%%%%%%%%%%%%
%%%%%%%%%%%%%%%%%%%%%%%%%%%%%%%%%%%%%%%%%%%%%%%%%%%%%%%%%%%
%%%%%%%%%%%%%%%%%%%%%%%%%%%%%%%%%%%%%%%%%%%%%%%%%%%%%%%%%%%
\clearpage

\section{Herz}

\subsubsection{Vorbereitung des Untersuchungsthemas}

Die Auskultation des Herzens wird nach dem 4 + 2 Schema von
Prof.~H.~Weber durchgeführt. Als Vorbereitung auf den
Auskulationsnachmittag wird dringend empfohlen, die Seiten
\VonBis{1}{19} des Skriptums \emph{Auskultation und
Perkussion des Herzens} (Download in den study guides)
gründlich zu studieren. 

\subsubsection{Theorie}

% \Parallel
{Die Herzaktion lässt sich in 2 Phasen einteilen: 

\begin{itemize}
  \item \T{Systole}: Anspannungs- und Austreibungsphase
  \item \T{Diastole}: Entspannungs- und Füllungsphase 
\end{itemize}}
% {\includegraphics[width=\linewidth]{./images/pku/clean-img77.jpg}}


\TextSlide
{Herztöne (HT)}
{~~

\noindent\Parallel
{~~

\T{1. Herzton} (\T{S\textsubscript{1}}):
\begin{itemize}
  \item Am lautesten über der \EE{Herzspitze}
  \item Niederfrequent
  \item Entsteht in der Anspannungsphase durch Anspannung
  der Ventrikelmuskulatur
  %     \footnote{\EE{Entstehung des
  %     1. Herztons}: Die Theorie, dass der S\textsubscript{1}
  %     durch den Klappenschluss der Mitra- und
  %     Trikuspidalklappe entsteht, gilt als überholt.}
  \item Schluss der AV-Klappen und Schwingung des ganzen Systems
  \item Carotispuls ist (mit leichter
  Zeitverzögerung) synchron tastbar
\end{itemize}
\T{2. Herzton} (\T{S\textsubscript{2}}):
\begin{itemize}
  \item Am lautesten über der \EE{Herzbasis}
  \item Etwas höher frequent
  \item Entsteht in der Entspannungsphase durch Schluss
  der Taschenklappen und Schwingung des Systems
  \item Es sind Anteile der Aorten- und Pulmonalklappe
  unterscheidbar (\T{A\textsubscript{2}},
  \T{P\textsubscript{2}}).
\end{itemize}
}
{\includegraphics[width=\linewidth]{./images/pku/herzzyklus.png}}


Auskultatorisch ist die Systole die Phase zwischen 1. und 2. HT. 
Die Diastole ist somit die Phase zwischen 2. und nächstem 1.
HT. Da sich die Diastole mit zunehmender Herzfrequenz
verkürzt, ist die Unterscheidung von Systole und Diastole
nicht immer einfach, die Orientierung anhand der
Pulspalpation ist hilfreich.}
{\MoveDown}
{}
{}
{}

 
\TextSlide
{Herzgeräusche}
{% Die Lautstärke wird in Sechstel angegeben:
% 
% \begin{labeling}[~:]{mmm}
%   \item[1/6] leise, kaum hörbar
%   \item[2/6] mittellaut
%   \item[3/6] leise, gut hörbar
%   \item[4/6] laut
%   \item[5/6] sehr laut (fortgeleitet auch außerhalb der Herzgegend zu
%   hören)
%   \item[6/6] Distanzgeräusch, mit freiem Ohr zu hören
% \end{labeling}
%
Geräusche entstehen physikalisch durch Übergang einer
laminaren in eine turbulente Strömung aus vielfachen Gründen:
Hindernis (Stenose, Verkalkungen), Flussumkehr
(Insuffizienz), etc. Herzgeräusche werden je nach Ihrem
zeitlichen Auftreten  in \T{Systolika} und \T{Diastolika}
unterschieden.

% 
% % ex, weil bei Befunden:
% \subparagraph{Charakterisierung}
%
% \begin{itemize}
%   \item \EE{zeitliche Zuordnung}: systolisch -
%   diastolisch
%   \item \EE{zeitlicher Verlauf}: crescendo-decrescendo,
%   gleichförmig, bandförmig, spindelförmig
%   \item \EE{Charakter}: hoch-, niederfrequent
%   \item \EE{Tonhöhe}: z.\,B. hell, laut
%   \item \EE{Ausstrahlung/Fortleitung}: Axilla, Carotiden
%   \item \EE{Veränderung durch Atmung, Lage}:
%   
%   z.\,B. physiologische Spaltung des 2. Herztones, Umlagerung oder
%   Bewegung
%   
%   z.\,B. akzidentielle Geräusche
% \end{itemize}
}
{\MoveDown}
{}
{}
{}




\TextSlide
{Auskultationspunkte}
{% Herzspitze: Mitralklappe: 5. ICR medial der MCL o. Palpation des Herzspitzenstoßes
% Trikuspidalklappe: 5. ICR rechts parasternal o. über Sternum
% Herzbasis: Aortenklappe: 2. ICR rechts parasternal
% Herzbasis: Pulmonalklappe: 2. ICR links parasternal
% Erb´scher Punkt: 3. ICR links parasternal (grundsätzlich alle Töne und Geräusche hier gut hörbar)
% 
Es gibt fünf klassische Auskultationspunkte, wobei vier
davon die vier Herzklappen repräsentieren:

\noindent\Parallel
{\begin{enumerate}
  \item \T[Mitralklappe!Auskultationspunkt]{Mitralklappe}:
  5. ICR medial der MCL bzw. über Herzspitzenstoß
  \item \T[Trikuspidalklappe!Auskultationspunkt]{Trikuspidalklappe}: 4. ICR rechts parasternal (oder über Sternum)
  \item \T[Erb'scher Punkt!Auskultationspunkt]{Erb'scher
  Punkt}: 3. ICR links parasternal (grundsätzlich alle Töne und Geräusche hier gut
  hörbar)
  \item \T[Aortenklappe!Auskultationspunkt]{Aortenklappe}:
  2. ICR rechts parasternal
  \item \T[Pulmonalklappe!Auskultationspunkt]{Pulmonalklappe}: 2. ICR links parasternal
\end{enumerate}
}
{\includegraphics[width=\linewidth]{./images/pku/Auskultationspunkte-Herz}}


% \begin{table}[H]
% \begin{WidePar}
% \FormatTable
% \begin{tabularx}{\linewidth}{XXX}
% \includegraphics[width=\linewidth]{./images/pku/clean-img80.jpg}
% &
% \includegraphics[width=\linewidth]{./images/pku/clean-img81.jpg}
% &
% \includegraphics[width=\linewidth]{./images/pku/clean-img82.jpg}
% \\\addlinespace
% Erb'scher Punkt
% &
% Pulmonalklappe
% &
% Mitralklappe
% \\\addlinespace
% \includegraphics[width=\linewidth]{./images/pku/clean-img83.jpg}
% &
% \includegraphics[width=\linewidth]{./images/pku/clean-img84.jpg} 
% &
% \includegraphics[width=\linewidth]{./images/pku/clean-img85.jpg} 
% \\\addlinespace
% Trikuspidalklappe
% &
% Aortenklappe
% &
% Auskultation auch in \EE{Linksseitenlage}.
% \\\addlinespace
% \end{tabularx}
% \end{WidePar}
% \end{table}
}
{\MoveDown}
{}
{}
{}



\TextSlide
{Fortleitungspunkte}
{\begin{itemize}
  \item Halsgefäße rechts und links (z.\,B.
  Aortenklappenstenose)
  \item Axilla links (z.\,B. bei Mitralinsuffizienz)
\end{itemize}
}
{\MoveDown}
{}
{}
{}

\begin{PictureSeriesWide}{}{Fortleitungspunkte}
\PictureSeriesRowWide{3}
{./images/pku/clean-img86.jpg}
{A. carotis rechts \ldots}
{./images/pku/clean-img87.jpg}
{{\ldots} und links}
{./images/pku/clean-img88.jpg}
{Axilla}
{}
{}
\end{PictureSeriesWide}



\subsubsection{Untersuchungsgang}

\TextSlide{Palpation des Herzspitzenstoßes}{%
Dieser entsteht während der Systole durch
das Anschlagen des linken Ventrikels an der Brustwand, und seine
Lokalisation ist von der Herzgröße abhängig. Man tastet im
Normalfall den Herzspitzenstoß in der
Medioclavicularlinie im 5. ICR li (in Linksseitenlage deutlicher).


Bei adipösen Patienten kann der Herzspitzenstoß abgeschwächt
sein, während er bei jungen, schlanken Personen deutlich zu spüren
ist. Eine Verlagerung des Herzspitzenstoßes nach lateral (seitlich)
kann im Rahmen einer Hypertrophie (Vergrößerung) des linken
Ventrikels erfolgen; weiters kann er bei Rechtsherzbelastung proximal
abweichen.
}
{\MoveDown}
{}
{}
{}

\begin{PictureSeries}{}{Palpation des Herzspitzenstoß}
\PictureSeriesRow{2}
{./images/pku/clean-img79.jpg}
{}
{empty}
{}
{}
{}
{}
{}
\end{PictureSeries}

\TextSlide{Mehrmalige Auskultation der Auskultationspunkte}
{Das Herz wird über allen fünf Auskultationspunkten mehrmals
abgehört, zuerst orientierund, danach differenzierend.

\begin{itemize}
  \item Mitralklappe
  \item Trikuspidalklappe
  \item Erb'scher Punkt
  \item Aortenklappe
  \item Pulmonalklappe
\end{itemize}

\bigskip}
{\MoveDown}
{}
{}
{}


\TextSlide{$\oplus$ Beurteilung S\textsubscript{1} über Herzspitze} 
{(Mitral-/ Trikuspidalklappe)\bigskip
\bigskip}
{\MoveDown}
{}
{}
{}

\TextSlide{$\oplus$ Beurteilung S\textsubscript{2} über Herzbasis} 
{Beurteilung des S\textsubscript{2} getrennt nach
A\textsubscript{2} und P\textsubscript{2} über der
Herzbasis(Aorten/Pulmonalklappe)

Falls Spaltung des S\textsubscript{2}:
Differenzierung 
\begin{itemize}
  \item Variabel: physiologisch (häufig)
  \item Fixe Spaltung: Verdacht auf
  \I{ASD~II}\footnote{\I{ASD}: Atrium-Septum-Defekt.}
\end{itemize}
}
{\MoveDown}
{}
{}
{}

\TextSlide{$\oplus$ Extratöne}
{Meist  telesystolisch/protodiastolisch (kurz vor oder nach
S\textsubscript{2}); 
falls vorhanden, versuche eine Differenzierung:

\noindent\begin{minipage}{\linewidth}
\FormatTable
\begin{tabularx}{\linewidth}{XXXX}
\toprule\addlinespace
\noindent\TabHeading{HT/ET}
&
\noindent\TabHeading{Auskult. Punkt}
&
\noindent\TabHeading{Laut/Leise}
&
\noindent\TabHeading{Frequenz}
\\\addlinespace\midrule\addlinespace
\TabSub{Gespaltener S\textsubscript{2}}
&
Basis
&
laut
&
hoch
\\\addlinespace
\TabSub{S\textsubscript{3}}
&
Erb
&
leise
&
nieder
\\\addlinespace
\TabSub{MÖT}
&
Spitze
&
laut
&
hoch
\\\addlinespace
\TabSub{Telesystol. Click}
&
Erb
&
laut
&
hoch
\\\addlinespace\bottomrule
\end{tabularx}
\end{minipage}

}
{\MoveDown}
{}
{}
{}


\TextSlide{$\oplus$ Herzgeräusche}
{Falls ja, versuche eine Differenzierung:

\begin{itemize}
  \item Aortenstenose (AS): niederfrequentes Systolikum,
  Fortleitung Carotiden
  \item Mitralinsuffizienz (MI): hochfrequentes 
  Systolikum, Fortleitung nach lateral
  \item Mitralstenose (MS):
  niederfrequentes Diastolikum
  \item Aorteninsuffizienz (AI):
  hochfrequentes Diastolikum
\end{itemize}

\EE{Merke}: Stenosen sind niederfrequent, Insuffizienzen sind hochfrequent.
}
{\MoveDown}
{}
{}
{}

\TextSlide{Bei V.\,a. Rechtsherzinsuffizienz: Inspektion der
Jugularvenen bds.}
{Hierfür liegt der Patient in 45{\textdegree}
Oberkörperhochlage auf der Untersuchungsliege und neigt seinen Kopf
etwas zur Seite (von Untersucher weg!). Lateral des entspannten
M. Sternocleidomastoideus
(Kopfkissen benutzen) und oberhalb der Clavicula kann der
Jugularvenenpuls (2 Spitzen/Herzschlag) inspektorisch begutachtet
werden. Evtl. mit der Untersuchungslampe von tangential her beleuchten
-- Seitenvergleich. Auf Stauungszeichen (Rechtsherzinsuffizienz)
achten.
}{\MoveDown}{}{}{}

\TextSlide{Bei V.\,a. Aorteninsuffizienz: Kapillarpuls am Fingernagel testen}
{Es wird ein leichter Druck auf den Fingernagel ausgeübt,
sodass sich eine klare Grenze zwischen rotem
(durchbluteten) und weißem (\enquote{abgedrückten}) Anteil
darstellt. Wandert diese Grenze rhythmisch hin- und her
(bei gleichbleibendem Druck auf den Nagel) ist der
Kapillarpuls positiv (typisch für Wasserhammerpuls bei Aorteninsuffizienz).
}
{\MoveDown}
{}
{}
{}




\TextSlide{Auskultationsbefund erstellen}
{
\begin{description}
  \item[Normalbefund]
  \Speech{\enquote{Der erste HT ist {\ldots} , der 2. HT ist {\ldots} ,
  keine Spaltung. Keine ET, keine Geräusche.}}
  \item[Auffälliger Befund] mit
  Besonderheiten:
  \begin{itemize}
    \item HT paukend oder Leise, S\textsubscript{2} fix
    gespalten, Extratöne vorhanden 
    
    Charakterisierung von Herzgeräuschen:
    \begin{itemize}
      \item Zeitliche Zuordnung: systolisch -- diastolisch
      \item Zeitlicher Verlauf: crescendo-decrescendo,
      gleichförmig, bandförmig, spindelförmig
      \item Charakter: hoch-, niederfrequent
      \item Ausstrahlung/Fortleitung: Axilla, Carotiden
    \end{itemize}
  \end{itemize}
  Nach der systematischen Beurteilung Versuch einer
  Verdachtsdiagnose (s.\,u.)
\end{description}
}
{
\begin{PkuBefunde}[]
  \item[\SymbolPkuNormal] \PkuBefund{Cor: rhythmisch, rein,
  normfrequent, keine Spaltung}
  \item[\SymbolPkuVariante] \PkuBefund{Variable Spaltung}
  \item[\SymbolFindingPathological] \PkuBefund{arrhythmisch,
  Systolikum mit punctum maximum über Aortenklappe mit
  Fortleitung in die Carotiden}
\end{PkuBefunde}
}
{}
{}
{}


\Text{Erstellen einer Verdachtsdiagnose}
{Siehe \prettyref{topic:herz-auskultation-befunde}.}
{}
{}
{}
{}




\bigskip
\subsubsection{Typische Befunde und Interpretation:
Herzauskultation} 

\label{topic:herz-auskultation-befunde}

% 
% \Text{}{%
% \begin{labeling}[:]{mm}
%   \item[Leitsymptom]
%   \item[I]
%   \item[P]
%   \item[Pk]
%   \item[A]
%   \item[»]
% \end{labeling}
% }{}{}{}{}
% 
\begin{WidePar}
% \thispagestyle{plain}
% \loadgeometry{FullPage}
% \begin{landscape}
{\small{
\noindent\begin{longtable}{p{0.3\linewidth}p{0.3\linewidth}p{0.3\linewidth}}
\toprule\addlinespace
\TabHeading{Leitsymptome} & \TabHeading{Inspektion, Palpation, Perkusion} & \TabHeading{Auskultation}
\\\addlinespace\midrule\addlinespace\endhead
\multicolumn{3}{c}{{\TabSub{Aortenklappenstose}}}
\\\addlinespace\midrule\addlinespace
Belastungsschwindel / Synkope, Angina
  pectoris, Dyspnoe
&
Pulsus parvus et tardus (kleiner Puls, träger Anstieg der Pulswelle)
& 
Raues, spindelförmiges Systolikum,
  p.\,m. (punktum maximum)  Aortenklappe, Fortleitung in die
  Carotiden

(DD: Aortenklappensklerose macht ebenfalls Systolikum,
  wird aber nicht in Carotiden fortgeleitet, DD Carotiden können auch wegen Carotisstenose selbst ein Strömungsgeräusch erzeugen)

\\\addlinespace\midrule\addlinespace
\multicolumn{3}{c}{{\TabSub{Mitralklappeninsuffizienz (MI)}}}
\\\addlinespace\midrule\addlinespace
Oft Vorhofflimmern: unregelmäßiger
  Puls, Palpitationen, Dyspnoe wg. Herzinsuffizienz
& 
Evtl. Facies mitralis (Zyanose im Wangenbereich)

Oft unregelmäßiger Puls
& 	
1. HT: leise
 hochfrequentes lange anhaltendes Systolikum, p.\,m. über der  Herzspitze fortgeleitet in linke Axilla
\\\addlinespace\midrule\addlinespace
\multicolumn{3}{c}{{\TabSub{Mitralklappenstenose (MS)}}}
\\\addlinespace\midrule\addlinespace
Oft Vorhofflimmern: unregelmäßiger
  Puls, Palpitationen, Dyspnoe wg. Herzinsuffizienz
& 
Facies mitralis

Oft unregelmäßiger Puls

& 
1.~HT: paukend (verzögerte Kammerfüllung,
ruckartiger Klappenschluss)

3.~HT: evtl P\textsubscript{2} gleichlaut oder  lauter als
A\textsubscript{2} (pulmonale Hypertension) mitralen
Öffnungston (MÖT) -- je kürzer der Abstand, desto
höhergradiger die Stenose -- niederfrequentes, Diastolikum
nach MÖT p.\,m. Herzspitze, schwer zu finden -- keine
Fortleitung

\\\addlinespace\midrule\addlinespace
\multicolumn{3}{c}{{\TabSub{Aortenklappeninsuffizienz (AI)}}}
\\\addlinespace\midrule\addlinespace
Dyspnoe
& 
pulsierende Gefäße bei großem Schlagvolumen (A.
  carotis) -- \I{Musset'sches Zeichen}

Positiver Kapillarpuls bei leichtem Druck auf
  Fingernagel, Vergrößerte Blutdruckamplitude, Pulsus celer et altus \enquote{Wasserhammerpuls}

& 
2.~HT: abgeschwächt bis fehlend; 
Hochfrequentes leises, \enquote{gießendes} Decrescendo-Diastolikum,
p.\,m. in der Mitte einer gedachten Linie zwischen
Aortenklappe und Herzspitze; Strömungsgeräusche auch über
großen Gefäßen, z.\,B. A.~brachialis oder A.~femoralis

\\\addlinespace\midrule\addlinespace
\multicolumn{3}{c}{{\TabSub{Trikuspidalinsuffizienz}}}
\\\addlinespace\midrule\addlinespace
Rechtsherz-Dekompensation
& 
Pulsation der Halsvenen (positiver Jugularvenenpuls)

Lebervenenpuls (Hand auf vergrößerte Leber legen)

& 
Wie MI, nur p.\,m. über Trikuspidalklappe.
  Geräusch wird bei Inspiration lauter

\\\addlinespace\midrule\addlinespace
\multicolumn{3}{c}{{\TabSub{Pulmonalstenose}}}
\\\addlinespace\midrule\addlinespace
& 

& 
Wie AS, p.\,m., jedoch über pulmonalis -- keine
  Fortleitung


\\\addlinespace\midrule\addlinespace
\multicolumn{3}{c}{{\TabSub{Ventrikelseptumdefekt}}}
\\\addlinespace\midrule\addlinespace
& 



& 
Protosystolikum mit p.\,m. Erb -- je leiser desto
hämodynamisch wirksamer = gefährlicher

\\\addlinespace\midrule\addlinespace
\multicolumn{3}{c}{{\TabSub{Perikardgeräusche}}}
\\\addlinespace\midrule\addlinespace
& 
& 
Herzsynchrones, ohrnahes Reibegeräusch bei trockener
Perikarditis, auch Systole/Diastole überlagernd, oft
flüchtig, verschwindet bei Auftreten eines Ergusses.
% 
% 
% \\\addlinespace\midrule\addlinespace
% \multicolumn{5}{c}{{\TabSub{}}}
% \\\addlinespace\midrule\addlinespace
% 

\\\addlinespace\bottomrule\addlinespace
\end{longtable}
}}
% \thispagestyle{fancy}
% \end{landscape}
\end{WidePar}
% \restoregeometry





\clearpage
\subsection{\CO{[C2]} Pulsstatus}



\TextSlide{Pulsstatus}{%
Der Pulsstatus wird im Seitenvergleich erhoben, wobei die
Palpation mit dem 2. und 3. Finger erfolgt. Das Tasten der
Pulse muss immer mit mindestens 2 Fingern einer Hand
erfolgen, da sonst der Eigenpuls getastet werden könnte. Man
beurteilt die Pulse im direkten
Seitenvergleich. 
% 
% \subparagraph{Rhythmus
% (Pulsfrequenz) und Pulsqualität}

Normal ist ein regelmäßiger Herzschlag mit einer Frequenz
um etwa \VonBis{60}{100}\PerMin*. Normalerweise zählt man
die Pulsschläge in 15 Sekunden und multipliziert sie mit 4, während
bei einer Arrhythmie oder einer Pulsfrequenz unter
50\PerMin* (Bradykardie) eine volle Minute gezählt werden sollte. Die
bei jungen (und bei \enquote{vegetativ}
übererregbaren) Menschen gefundene respiratorische Arrhythmie mit
Frequenzzunahme bei Inspiration und Verringerung bei Exspiration hat
keinen Krankheitswert.


Die Pulsfrequenz ist vom Alter und v.\,a. vom
Trainingszustand abhängig; bei einem durchtrainierten, jungen Mann
ist eine Frequenz von 50 nicht außergewöhnlich, während diese
bei einer älteren, über Schwindel klagenden Patientin zur
Abklärung animieren sollte. Außerdem ist die Feststellung eines
regelmäßigen oder unregelmäßigen Rhythmus diagnostisch
wichtig.


Die Pulsqualität beschreibt die
Blutdruckamplitude (also die Differenz zwischen systolischem und
diastolischem Blutdruck) und unterscheidet z.\,B. einen hohen und harten
Puls (\I[Pulsus!celer et altus]{Pulsus celer et altus} oder
\enquote{\I{Wasserhammerpuls}}) bei
Aorteninsuffizienz (unvollständiger Schluss der Aortenklappe) von
einem niedrigen und trägen Puls (\I[Pulsus!parvus
et tardus]{Pulsus parvus
et tardus}) bei \I[Puls!Aortenstenose]{Aortenstenose}
(Verengung der Aortenklappe).

}{
\begin{PkuBefunde}[]
  \item[\SymbolPkuNormal] \PkuBefund{Pulse der {\ldots}
  seitengleich und gut palpabel, normfrequent}
%   \item[\SymbolPkuVariante] \PkuBefund{}
%   \item[\SymbolFindingPathological] \PkuBefund{}
\end{PkuBefunde}
}{}{}{}



\begin{PictureSeriesWide}{}{Übersicht Pulsstatus}
\PictureSeriesRowWide{3}
{./images/pku/clean-img141.jpg}
{A. carotis}
{./images/pku/clean-img133.jpg}
{A. radialis}
{./images/pku/clean-img134.jpg}
{A. brachialis }
{}
{}
\PictureSeriesRowWide{3}
{./images/pku/clean-img135.jpg}
{A. axillaris}
{./images/pku/clean-img136.jpg}
{A. femoralis }
{./images/pku/clean-img137.jpg}
{A. poplitea }
{}
{}
\PictureSeriesRowWide{3}
{./images/pku/clean-img138.jpg}
{A. tibialis posterior}
{./images/pku/clean-img139.jpg}
{A. dorsalis pedis }
{empty}
{}
% {./images/pku/clean-img140.jpg}
% {Aorta abdominalis }
{}
{}
\end{PictureSeriesWide}








%%%%%%%%%%%%%%%%%%%%%%%%%%%%%%%%%%%%%%%%%%%%%%%%%%%%%%%%%%%
%%%%%%%%%%%%%%%%%%%%%%%%%%%%%%%%%%%%%%%%%%%%%%%%%%%%%%%%%%%
%%%%%%%%%%%%%%%%%%%%%%%%%%%%%%%%%%%%%%%%%%%%%%%%%%%%%%%%%%%
%%%%%%%%%%%%%%%%%%%%%%%%%%%%%%%%%%%%%%%%%%%%%%%%%%%%%%%%%%%
%%%%%%%%%%%%%%%%%%%%%%%%%%%%%%%%%%%%%%%%%%%%%%%%%%%%%%%%%%%
\clearpage
\section
[Wochenthema \CO{[D]}]
{Wochenthema \CC{[D]}}

\HeadingSub{Abdomen, inguinale und rektale Untersuchung}

\subsection{Übersicht: Untersuchungsgang}
\Text{Beginn}
{\begin{itemize}
  \item Hygienische Händedesinfektion
  \item Untersucher stellt sich vor
  \item Patient macht sich frei
\end{itemize}}
{}
{}
{}
{}

\TextSlide{\CO{[D1]} Abdomen}
{
\begin{itemize}
  \item \EE{Inspektion im
  Stehen} (Hernien, Varizen, Ödeme, inkl. Haut)
  \item \EE{Inspektion des Abdomens} in Rückenlage
  (im / über / unter unter Thoraxniveau, Haut, Narben, Behaarung)
  \item Ausführliche \EE{Auskultation des Abdomens} in allen 4 Quadranten (30\Sec* pro
  Quadrant, Strömungs-/ Darmgeräusche)
  \item \EE{Perkussion und Palpation} des
  \EE{rechten oberen Quadranten}
  \begin{itemize}
    \item \EE{Leber}, \EE{Lebergrenzen}, Gallenblase
  \end{itemize}
  \item \EE{Perkussion und Palpation} des \EE{linken
  oberen Quadranten}
  \begin{itemize}
    \item Magen, bimanuelle Palpation der \EE{Milz}
  \end{itemize}
  \item \EE{Perkussion und Palpation} des \EE{linken
  unteren Quadranten}
  \begin{itemize}
    \item Sigma, Harnblase, Loslassschmerz
  \end{itemize}
  \item \EE{Perkussion und Palpation} des
  \EE{rechten unteren Quadranten}
  \begin{itemize}
    \item Colon ascendens, \EE{Appendixdruckpunkte}
  \end{itemize}
  \item \EE{Stosswellenperkussion} eines potentiellen
    Aszites 
  \item \EE{Flankendämpfung}
%   \item \EE{Perkussion des Abdomens} (Magen,
%   Darmschlingen, Flankendämpfung, Ascites)
%   \item \EE{Palpation des Abdomens} (in Rückenlage)
%   \begin{itemize}
%     \item Oberflächliche Palpation (auf den Schmerz
%     zu, Abwehrspannung, Loslassschmerz)
%     \item Tiefe Palpation (inklusive
%     Appendixdruckpunkte)
%   \end{itemize}
%   \item \EE{Palpation der Leber}
%   \begin{itemize}
%     \item \EE{Kratz-Auskultation} der Leber
%     \item \EE{Leberperkussion} (oberer/unterer Rand)
%   \end{itemize}
  \item \EE{Perkussion beider Nierenlager} (Klopfschmerz)
\end{itemize}

}
{
\begin{PkuBefunde}[]
  \item[\SymbolPkuNormal] \PkuBefund{Abdomen im
  Thoraxniveau, lebhafte Darmgeräusche in allen 4
  Quadranten, Leber unter Rippenbogen, Milz nicht palpabel,
  Abdomen weich, kein Druckschmerz, kein Loslassschmerz
  (McBurney und Lanz indolent), kein Hinweis auf Aszites,
  keine Umbilical palpabel,
  Nierenlager beidseits indolent}
  \item[\SymbolPkuVariante] \PkuBefund{Blande Narbe {\ldots}
  / spärliche, leise Darmgeräsuche; Leber nicht palpabel}
%   \item[\SymbolFindingPathological] \PkuBefund{}
\end{PkuBefunde}
}
{}
{}
{}

\TextSlide{\CO{[D2]} Inguinale und rektale Untersuchung}
{%
\begin{itemize}
  \item \EE{Inspektion der Leistenregion} (Hernien)
  \item \EE{Palpation der inguinalen Lymphknoten}
%   \item \EE{Hodenuntersuchung}
  \item \EE{Rektale
  Untersuchung} (in Seitenlage am Modell)
  \begin{itemize}
    \item \EE{Inspektio}n (Ruhe, Pressen)
    \item \EE{Palpation}
    \begin{itemize}
      \item des äußeren Sphinkters
      \item der Ampulla recti
      \item der Prostata / der
      Portio (Verschiebeschmerz, Peritoneale Reizung)
      \item Inspektion der Handschuhe und gegebenenfalls Durchführung des
      Hämoccult
    \end{itemize}
    \item \EE{Hämoccult} auswerten
  \end{itemize}
\end{itemize}}
{\begin{PkuBefunde}[]
  \item[\SymbolPkuNormal] \PkuBefund{Keine Inguinalhernie
  oder ing. Lymphknoten palpabel; Digital-rektale
  Untersuchung: Analregion unauffällig, Ampullenschleimhaut
  und Prostata palp.
  unauffällig, Ampulle frei, Hämoccult negativ}
%   \item[\SymbolPkuVariante] \PkuBefund{}
%   \item[\SymbolFindingPathological] \PkuBefund{}
\end{PkuBefunde}}
{}
{}


\Text{Ende}
{
\begin{itemize}
  \item Patient kleidet sich wieder an
  \item Offene Fragen beantworten
  \item Verabschiedung
  \item Hygienische Händedesinfektion
\end{itemize}
}
{}
{}
{}
{}


\bigskip
% \clearpage
\subsection{\CO{[D1]} Abdomen}

\subsubsection{Allgemeines}

Die Untersuchung des Abdomens sollte im Liegen durchgeführt
werden und umfasst \EE{Inspektion}, \EE{Auskultation},
\EE{Palpation} und \EE{Perkussion} (Reihenfolge beachten!)
mehrerer Organe.

Die Entspannung des Patienten ist bei der Untersuchung des
Abdomens besonders wichtig, deshalb wird der Patient mit
einem Kissen unter dem Kopf auf dem Rücken gelagert,
eventuell können die Beine angewinkelt werden.

Bei der Untersuchung zu beachten: mit dem Patienten reden,
ablenken, unterhalten, tief ein- und ausatmen lassen,
Untersuchung mit warmen Händen (kalte Hände führen zu einer
Abwehrspannung).
Grundsätzlich soll man feinfühlig vorgehen, sich Zeit nehmen
und Ruhe ausstrahlen.

\Text{Einteilung des Abdomens in Quadranten}
{Orientierend denkt man sich 2 senkrecht aufeinander stehende Linien
durch den Nabel, die das Abdomen in 4 Quadranten
einteilen.\footnote{Eine weitere Einteilung unterscheidet
das obere Drittel (\I{Epigastrium}, Oberbauch) vom mittleren
Drittel (\I{Mesogastrium}, Mittelbauch) und unterem Drittel (\I{Hypogastrium}, Unterbauch).}
Die Region um den Nabel wird als \T{periumbilical}
bezeichnet.

}
{}
{}
{}
{}

\begin{PictureSeries}{}{Bauchmuskeln und Einteilung des Abdomens
in Quadranten. \SourcePicture{Quelle nicht eruierbar.}}
\PictureSeriesRow{2}
{./images/pku/clean-img107}
{}
{./images/pku/Abdomen-Einteilung}
{}
{}
{}
{}
{}
\end{PictureSeries}






\subsubsection{Untersuchungsgang}

\TextSlide{Inspektion im Stehen}
{Der erste Schritt ist wie immer die
\E{Inspektion}.
Hierbei sollte man auf eventuelle Narben
und Vorwölbungen achten und diese notieren. Eine umschriebene
Vorwölbung kann Hinweise auf einen malignen Prozess oder Hernien
geben, eine Vorwölbung des gesamten Abdomens kann auf einer
Adipositas oder auf freier Flüssigkeit im Abdomen beruhen (Aszites).

Weitere Besonderheiten können das Auftreten von Striae (z.\,B.
geriffelte Streifen durch Überdehnung) oder ein Caput medusae
(verdickte, geschlängelte Venen, z. B. bei Leberzirrhose) sein.}
{\MoveDown}
{}
{}
{}

\TextSlide{Inspektion in Rückenlage}
{Die Inspektion in Rückenlage erfolgt analog zur Inspektion
im Stehen, im Liegen kann man speziell das Abdominal- und
Thoraxniveau miteinander vergleichen.
}
{\begin{PkuBefunde}[]
  \item[\SymbolPkuNormal] \PkuBefund{Abdomen im
  Thoraxniveau, inspektorisch unauffällig}
  \item[\SymbolPkuVariante] \PkuBefund{Blande Narbe {\ldots}}
%   \item[\SymbolFindingPathological] \PkuBefund{}
\end{PkuBefunde}}
{}
{}
{}


\TextSlide{Ausführliche Auskultation des Abdomens in allen 4
Quadranten} 
{Die Peristaltik des Dünndarms erzeugt charakteristische Geräusche.
Mit dem Stethoskop hört man ca. alle \VonBis{5}{10}~Sekunden
glucksende, gurgelnde oder knarrende Geräusche. Das Vorhandensein dieser
Geräusche wird in allen 4 Quadranten je ca. \EE{30~Sekunden}
lang überprüft. Bei verstärkter Peristaltik (z.\,B. Hunger, Enteritis,
Malabsorption) sind die Geräusche andauernd in erhöhter
Lautstärke hörbar.


\subparagraph{Pathologische Darmgeräusche} Verminderte oder
metallische Darmgeräusche weisen auf eine Pathologie hin (mechanischer Ileus). 
Das Fehlen jeglicher
Darmgeräusche (\enquote{\I{Grabesstille}}) über einen Zeitraum von
mindestens 3 Minuten spricht in Verbindung mit der
Anamnese für einen paralytischen Ileus.\footnote{Eine Paralyse mit verminderten oder
fehlenden Darmgeräuschen kann auch durch Medikamente,
z.\,B. im Rahmen einer akuten Schmerztherapie
(beispielsweise mit Butylscopolamin
(Buscapina{\texttrademark})), beabsichtigt verursacht
werden.} Besteht der Verdacht auf einen Ileus muss
dieser zeitnah weiter abgeklärt werden.

}
{\begin{PkuBefunde}[]
  \item[\SymbolPkuNormal] \PkuBefund{Lebhafte Darmgeräusche in allen 4
  Quadranten}
  \item[\SymbolPkuVariante] \PkuBefund{Spärliche, leise
  Darmgeräusche in allen 4 Quadranten}
  
  \PkuBefund{Fehlende Darmgeräusche (nach Gabe von
  Butylscopolamin vor 1\Hour*)}
  \item[\SymbolFindingPathological] \PkuBefund{Fehlende
  Darmgeräusche}
  
  \PkuBefund{Metallische Darmgeräusche}
\end{PkuBefunde}}
{}
{}
{}


\begin{PictureSeries}{}{Auskultation des Abdomens}
\PictureSeriesRow{2}
{./images/pku/clean-img109.jpg}
{}
{./images/pku/clean-img110.jpg}
{}
{}
{}
{}
{}
\PictureSeriesRow{2}
{./images/pku/clean-img111.jpg}
{}
{./images/pku/clean-img112.jpg}
{}
{}
{}
{}
{}
\end{PictureSeries}



\Text{Anmerkungen zur Palpation des Abdomens}
{Bei der Perkussion und Palpation wird quadrantenweise
vorgegangen. Beide Methoden dienen dazu, die \EE{Größe und
Lage der in den Quadranten befindlichen Organe orientierend zu
untersuchen}. Weiters wird immer die \EE{Bauchdecke}
(weich, hart) und das Vorhandensein von \EE{Druckschmerzen}
beurteilt. Die Palpation wird erst oberflächlich und dann
tief (beidhändig) vorgenommen, um einerseits eine
Abwehrspannung, andererseits oberflächliche Veränderungen zu
erkennen.  \EE{Zuerst wird oberflächlich, und
erst zuletzt tief palpiert.} Es empfiehlt sich, mit jeder
Exspiration des Patienten tiefer zu drücken.

Manchmal können Resistenzen ertastet werden, die pathologischer Natur
oder auch im Rahmen der Norm sein können (z.\,B. \I{Skybala}
(Kotballen), eine als Tumor im Unterbauch imponierende
prall gefüllte Harnblase). 

\begin{description}
  \item[Loslassschmerz]
Erst wird die Bauchwand für einige Sekunden eingedrückt, was einen
mit der Zeit nachlassenden Druckschmerz verursachen kann. Ein bei
plötzlichem Loslassen auftretender Schmerz ist Anzeichen einer
Peritonitis (z.\,B. im Rahmen einer Appendizitis). 
\item[Der schmerzhafte Bauch]
Gibt der
Patient Spontanschmerzen an, so sollte man erst die
nicht betroffenen Teile des Abdomens untersuchen, und sich
erst nach und nach zu dem schmerzhaften Bereich
vorzuarbeiten.
\end{description}



}{}{}{}{}


\begin{PictureSeries}{}{}
\PictureSeriesRow{2}
{./images/pku/clean-img114.jpg}
{Handhaltung}
{./images/pku/clean-img115.jpg}
{}
{}
{}
{}
{}
\end{PictureSeries}

\TextSlide{Perkussion und Palpation des rechten oberen Quadranten} 
{Die Technik der Perkussion wird im Detail unter
\prettyref{topic:perkussion} auf
Seite \pageref{topic:perkussion} vorgestellt.
Bei der Perkussion und Palpation wird quadrantenweise
vorgegangen. Beide Methoden dienen dazu, die \EE{Größe und
Lage der in den Quadranten befindlichen Organe orientierend zu
untersuchen}.
Interessante Strukturen im Rechten oberen Quadranten sind
besonders die \EE{Leber} und die \EE{Gallenblase}. 

\begin{description}
  \item[Perkussion der Lebergrenzen]
Man perkutiert im Verlauf der \E{Medioklavikularline}.
Bei der Perkussion achtet man auf den Übergang vom sonoren
Klopfschall der Lunge zur Dämpfung durch den oberen Rand der
Leber (Höhe ca. 6. Rippe) sowie den Übergang zum sonoren bis
tympanitischen Schall im Mesogastrium.
\item[Palpation der Leber]
Eine nicht vergrößerte Leber ist auch bei Inspiration nicht oder
nur schwer unterhalb des Rippenbogens tastbar. Bei pathologischen
Veränderungen (Vergrößerung) ist bei der Palpation auf
Oberfläche und Konsistenz (Härte) zu achten. Das Ausmaß der
Vergrößerung wird in Querfingern\footnote{Besser, aber
weniger gebräuchlich: Angabe in Zentimetern.} unterhalb des
Rippenbogens angegeben.
\item[Palpation der Gallenblase]
Mit Übung lässt sich evtl. die gefüllte
Gallenblase palpieren. Man taucht hierzu caudal des
Lebervorderrandes medial der Medioklavikularlinie mit den
Händen unter die Leber ab. Es ist auf Druckschmerz zu achten
(\I{Murphy-Zeichen}). 
\item[Kratz-Auskultation der Leber]
Aufsetzen des Stethoskops im Epigastrium (unter das Xiphoid). Man kratzt
mit einem Holzspatel entlang der Medioklavikularlinie beginnend auf
Höhe der 5. Rippe in engen Abständen über das Leberareal nach
kaudal. Befindet sich der Spatel über der Leber, so kann man das
Kratzen mit dem Stethoskop hören. Kranial und kaudal der Leber hört
man kaum etwas (ungenaue Bestimmung der Lebergrenzen).
\end{description}

}{
\begin{PkuBefunde}[]
  \item[\SymbolPkuNormal] \PkuBefund{Leber unter
  Rippenbogen; Abdomen: rechter oberer Quadrant weich, kein
  Druckschmerz, kein Loslassschmerz}
%   \item[\SymbolPkuVariante] \PkuBefund{}
%   \item[\SymbolFindingPathological] \PkuBefund{}
\end{PkuBefunde}
}{}{}{}


\begin{PictureSeries}{}{}
\PictureSeriesRow{2}
{./images/pku/clean-img118.jpg}
{\TabSub{Konventionelle Palpation der Leber und der
Gallenblase}
\\
Man legt die Fingerkuppen \VonBis{3}{5}\CM* unterhalb des
Rippenbogens in der Medioklavikularlinie auf den Mittelbauch.
\\
Während der Patient tief einatmet (forcierte Bauchatmung) wandert der
kaudale Leberrand nach unten auf die Fingerkuppen zu (Gleitpalpation).}
{./images/pku/clean-img119.jpg}
{\TabSub{Alternative Palpation der Leber}
\\
Man steht auf der rechten Thoraxseite des Patienten.
\\
Die Finger beider Hände drücken nach dorsal und kranial unter den
Rippenbogen.
\\
Lassen Sie den Patienten tief einatmen und tasten Sie mit den
Fingerkuppen den kaudalen Leberrand.}
{}
{}
{}
{}
\PictureSeriesRow{2}
{./images/pku/gabriel-sebastian/Gallenblase-Palpation-001-00800.jpg}
{Palpation der Gallenblase}
{./images/pku/gabriel-sebastian/Gallenblase-Palpation-002-00800.jpg}
{Palpation der Gallenblase}
{}
{}
{}
{}
\end{PictureSeries}

\begin{PictureSeries}{}{Kratzauskulatation der Leber}
\PictureSeriesRow{2}
{./images/pku/clean-img120.jpg}
{}
{./images/pku/clean-img121.jpg}
{}
{}
{}
{}
{}
\end{PictureSeries}

\TextSlide{Perkussion und Palpation des linken oberen Quadranten} 
{Interessante Strukturen im linken oberen Quadranten sind
der Magen und die Milz. Meistens lässt sich leicht eine
Luftblase im Magen perkutieren. 
\begin{description}
  \item[Palpation der Milz]
Die Milz ist beim gesunden Erwachsenen
nicht palpabel. Eine palpable, vergrösserte Milz kommt
v.\,a. im Rahmen von hämatologischen Erkankungen vor und ist
ein wichtiges Symptom.
Die rechte Hand des Untersuchers fasst dorsal an die linke
Flanke des Patienten in Höhe des Oberbauches. Mit der linken
Hand versucht man nun die Milz zu tasten und nutzt dabei
die rechte Hand als Widerlager.
\end{description}

}
{
\begin{PkuBefunde}[]
  \item[\SymbolPkuNormal] \PkuBefund{Milz nicht palpabel,
  Abdomen im linken oberen Quadranten weich, kein
  Druckschmerz, kein Loslassschmerz }
%   \item[\SymbolPkuVariante] \PkuBefund{}
%   \item[\SymbolFindingPathological] \PkuBefund{}
\end{PkuBefunde} 
}
{}
{}
{}

\begin{PictureSeries}{}{Palpation der Milz}
\PictureSeriesRow{2}
{./images/pku/gabriel-sebastian/Milzpalpation-001-00800.jpg}
{}
{}
{}
{}
{}
{}
{}
\end{PictureSeries}

\TextSlide{Perkussion und Palpation des linken unteren Quadranten} 
{Hier lässt sich häufig das Sigma gut palpieren, besonders
wenn es mit Stuhl gefüllt ist.

Es soll auch auf Loslassschmerzen untersucht werden; im
Rahmen einer Appendizitis kommt es mitunter zu einem
\E{kontralateralen} Loslassschmerz.

Im Unterbauch kann man oft auch eine gefüllte Harnblase
auffinden.

}
{
\begin{PkuBefunde}[]
  \item[\SymbolPkuNormal] \PkuBefund{
  Abdomen li. unterer Quadrant: weich, kein Druckschmerz,
  kein Loslassschmerz}
%   \item[\SymbolPkuVariante] \PkuBefund{}
%   \item[\SymbolFindingPathological] \PkuBefund{}
\end{PkuBefunde}
}
{}
{}
{}

\Text{Perkussion und Palpation des rechten unteren Quadranten}
{Im rechten Unterbauch lässt sich das Colon ascendens
auffinden. Weiter liegen hier die klinischen bedeutsamen
\T{Appendixdruckpunkte}.

\begin{description}
  \item[Appendixdruckpunkte]
  Hinweis für das
  Vorliegen einer Appendizitis ist die Schmerzhaftigkeit an spezifischen
  Punkten:
  \begin{itemize}
    \item \T{McBurney Punkt}:
    Mitte zwischen rechter Spina iliaca anterior superior und Nabel
    \item \T{Lanz'scher Punkt}:
    Zwischen rechtem und mittlerem Drittel zwischen beiden
    Spinae iliacae anteriores superiores
  \end{itemize}
\end{description}

}
{}
{}
{}
{}

\begin{PictureSeries}{}{}
\PictureSeriesRow{2}
{./images/pku/clean-img116.jpg}
{\T{McBurney Punkt}
(Mitte zwischen rechter Spina iliaca anterior superior und Nabel)}
{./images/pku/clean-img117.jpg}
{\T{Lanz'scher Punkt}
(Zwischen rechtem und mittlerem Drittel zwischen beiden
Spinae iliacae anteriores superiores).}
{}
{}
{}
{}
\end{PictureSeries}

\TextSlide{Besonderheit: Stosswellenperkussion}
{Man legt eine Hand auf die Flanke des Patienten und klopft mit den
Fingerkuppen der anderen Hand kurz und scharf auf die
gegenüberliegende Flanke. Spürt man ein Anschlagen einer
Flüssigkeitswelle (Undulationsphänomen) an der tastenden Hand so
liegt Aszites vor.

}
{
\begin{PkuBefunde}[]
  \item[\SymbolPkuNormal] \PkuBefund{Kein Hinweis auf
  Aszites}
%   \item[\SymbolPkuVariante] \PkuBefund{}
%   \item[\SymbolFindingPathological] \PkuBefund{}
\end{PkuBefunde}
}
{}
{}
{}


\begin{PictureSeriesWide}{}{Stowwellenperkussion}
\PictureSeriesRowWide{3}
{./images/pku/gabriel-sebastian/Stosswellenperkussion-001-00800.jpg}
{}
{./images/pku/gabriel-sebastian/Stosswellenperkussion-002-00800.jpg}
{}
{./images/pku/gabriel-sebastian/Stosswellenperkussion-003-00800.jpg}
{}
{}
{}
\end{PictureSeriesWide}


\TextSlide{Besonderheit: Flankendämpfung}
{Größere Flüssigkeitsmengen im Abdomen bewirken in Rückenlage
eine Flankendämpfung des Klopfschalls. Den Übergang vom
tympanitischen zum gedämpften Schall markiert man und bittet den
Patienten sich in Seitenlage zu drehen. Danach sollte bei vorliegendem
Ascites ein Richtung Nabel verschobener Übergang der
Schallqualitäten auftreten (Flüssigkeit folgt der Schwerkraft, und
die Luft steigt auf).}
{
\begin{PkuBefunde}[]
  \item[\SymbolPkuNormal] \PkuBefund{Kein Hinweis auf
  Aszites}
%   \item[\SymbolPkuVariante] \PkuBefund{}
%   \item[\SymbolFindingPathological] \PkuBefund{}
\end{PkuBefunde}
}
{}
{}
{}




\Text{Perkussion der Nierenlager}{%
Um überflüssige Lageveränderungen des Patienten zu vermeiden, wird
die Perkussion der Nierenlager sinnvollerweise im Rahmen der
Thoraxuntersuchung (Wochenthema B) vorgenommen. 


Der Patient sitzt und man klopft mit der Handkante zunächst vorsichtig
in die Flanken (kostovertebraler Winkel). Normalerweise soll eine
Erschütterung aber kein Schmerz ausgelöst werden. Patienten mit
entzündlichen Prozessen (z.\,B. Pyelonephritis) geben bereits bei
leichter Erschütterung in dieser Gegend starke Schmerzen an.
}{
\begin{PkuBefunde}[]
  \item[\SymbolPkuNormal] \PkuBefund{Nierenlager beidseits
  indolent}
%   \item[\SymbolPkuVariante] \PkuBefund{}
%   \item[\SymbolFindingPathological] \PkuBefund{}
\end{PkuBefunde}
}{}{}{}


\begin{PictureSeries}{}{Perkussion der Nierenlager}
\PictureSeriesRow{2}
{./images/pku/clean-img124.jpg}
{}
{./images/pku/clean-img125.jpg}
{}
{}
{}
{}
{}
\end{PictureSeries}





\subsubsection{Epidemiologie der akuten abdominellen Schmerzen}

Es existieren mehr als 1\,000 Ursachen (sic!) Mehr als 
80\PC* werden allerdings durch \E{drei} Diagnosen
verursacht:
\begin{itemize}
  \item Nonspecific Abdominal Pain (NSAP):  \VonBis{34}{53,7}\PC* 
  \item Appendizitis: \VonBis{18,9}{28,1}\PC*
  \item Cholezystitis: \VonBis{7,0}{9,7}\PC*
\end{itemize}

Die anderen, z.\,T. \enquote{klassische} Diagnosen sind eher im
niedrig-einstelligen Prozentbereich angesiedelt, als
Beispiel seien genannt: 

\begin{itemize}
  \item Dünndarmileus: \VonBis{3,1}{4,1}\PC*
  \item Perforiertes peptisches Ulkus:
  \VonBis{2,5}{3,1}\PC*
  \item Pankreatitis: \VonBis{1,6}{2,9}\PC*
  \item Divertikulitis: \VonBis{1,3}{2}\PC*
  \item Urologische Probleme: 5,5\PC*
  \item Gynäkologische Probleme: \VonBis{4,0}{4,9}\PC*
  \item Andere: \VonBis{1,2}{13,3}\PC*
\end{itemize}





\bigskip
% \cleardoublepage
\subsection{\CO{[D2]} Inguinale und rektale Untersuchung}

\Text{Inspektion der Leistenregion}{%
Die wesentlichen Befunde der Leistenregion sind Hernien,
vergrößerte Lymphknoten sowie der Puls der Arteria femoralis.
}{}{}{}{}




\begin{table}[H]
\includegraphics[width=\linewidth]{./images/pku/Hernien} 

\Captionof{figure}{Lokalisation häufiger Hernien:
Femoralishernie, Inguinalhernie, Umbilicalhernie,
epigastrische Hernie. \SourcePicture{Quelle nicht
eruierbar.}}
\end{table}







\TextSlide{Palpation der inguinalen Lymphknoten}{%
Lymphknoten sind nur bei Vergrößerung tastbar. Es handelt
sich also um eine orientierende Palpation der entsprechenden Regionen. Sie
können reaktiv entzündlich vergrößert sein (druckschmerzhaft, weich
und verschieblich), im Rahmen von malignen oder systemischen
Erkrankungen sind sie meist verhärtet und indolent. Siehe auch
Wochenthema A.

Bei der Palpation der Leiste wird man häufig einige, maximal
erbsengroße, nicht druckdolente Lymphknoten tasten können.
Überprüfen Sie ob die Lymphknoten druckschmerzhaft, verschieblich,
ein- oder beidseitig vergrößert sind.

}{
\begin{PkuBefunde}[]
  \item[\SymbolPkuNormal] \PkuBefund{Inguinale Lymphknoten
  bds. nicht palpabel}
%   \item[\SymbolPkuVariante] \PkuBefund{}
%   \item[\SymbolFindingPathological] \PkuBefund{}
\end{PkuBefunde}
}{}{}{}

\begin{PictureSeries}{}{Inguinale Lymphknoten}
\PictureSeriesRow{2}
{./images/pku/LK-ing}
{}
{./images/pku/clean-img129.jpg}
{}
{}
{}
{}
{}
\end{PictureSeries}


\TextSlide{{\SignErwaehnt} Leistenhernien}{%
Man legt den Zeige-, Mittel- und Ringfinger an den
Leistenkanal und lässt den Patienten husten. Bei einer
Leistenhernie spürt man den Peritonealsack, welcher durch
die abdominelle Druckerhöhung beim Husten tiefer tritt.
\cite{DahmerAnamneseBefund10} 

}{
\begin{PkuBefunde}[]
  \item[\SymbolPkuNormal] \PkuBefund{Leistenhernien bds.
  nicht palpabel}
%   \item[\SymbolPkuVariante] \PkuBefund{}
%   \item[\SymbolFindingPathological] \PkuBefund{}
\end{PkuBefunde}}{}{}{}





% \Text{Klinische Untersuchung des
% \IX[Skrotum]{Skrotums}} 
% {Die klinische Untersuchung besteht
% aus Inspektion und Palpation, wobei die Genitalorgane und auch der Gesamthabitus unter Berücksichtigung
% der sekundären männlichen Geschlechtsmerkmale beurteilt werden.
% 
% \subparagraph{Inspektion} Bei der Inspektion achtet man auf
%   andrologische Störungen, wie beispielsweise Hodenatrophie (bei Störung der Androgenbildung) oder
% testikuläre Feminisierung. 
% 
% Bei den Genitalorganen beurteilt man Abweichungen in Größe und
% Form (anatomische Veränderungen, Fehlbildungen wie Mikrophallus),
% Behaarung (z.\,B. horizontale Schamhaargrenze), Operationsnarben, Tumore,
% Schwellungen der regionären Lymphknoten, entzündliche
% Veränderungen im Bereich des äußeren Vorhautblattes (Phimose,
% Hypospadie und Epispadie) und artifizielle Veränderungen, wie
% beispielsweise \IX[Piercing!Intim-]{Intim-Piercings}.
% 
% \subparagraph{Palpation} Der Patient sollte bei der Untersuchung sizen oder liegen.
% 
% Bei der palpatorischen Untersuchung des Skrotums werden die Skrotalhaut
% und der Skrotalinhalt beurteilt. Dabei wird der Hoden sanft zwischen
% Daumen und Zeige- und Mittelfinger gerollt und abgetastet und die
% Größe verglichen. Meist ist ein Hoden von Natur aus etwas
% größer und tiefer gelegen als der andere. 
% 
% Wichtig ist auch, dass die Rückseite des Skrotums durch Anheben
% beurteilt wird.
% 
% \subparagraph{Auf folgende Veränderungen sollten geachtet
% werden:}
% \begin{itemize}
%   \item {Schmerzlose Vergrößerung eines Hodens}
%   \item {Tastbarer Knoten in einem der beiden Hoden}
%   \item {Schweregefühl oder ziehender Schmerz im Hoden }
%   \item {Berührungsempfindlichkeit eines Hodens}
%   \item {Flüssigkeitsansammlung im Hodensack}
%   \item {Ausfluss aus dem Penis}
% \end{itemize}
% 
% 
% Bei {\textquotedbl}leerem Skrotum{\textquotedbl} (z.\,B. bei
% Kryptorchismus oder Anorchie) sind alle Regionen zu untersuchen, in
% denen der Hoden liegen könnte (z.\,B. Maldeszensus testis bei Lage im
% Abdomen oder im Leistenkanal, oder Hodenektopie bei Nachweis der Hoden
% in der cruralen, perinealen oder penilen Region). 
% 
% Auch bei primär unauffälligem Skrotalinhalt sind Hoden und
% Nebenhoden palpatorisch genau zu überprüfen.
% 
% Eine Vermehrung des Skrotalinhaltes tritt auf bei Hydrozelen,
% Skrotalhernien, Spermatozelen, Entzündungen und Malignomen der Hoden
% und Nebenhoden, sowie bei Varikozelen. Mit einer Durchleuchtung des
% Skrotalinhaltes von hinten (Diaphanoskopie) kann man eine einfache
% klinische Differentialdiagnose zwischen soliden und zystischen
% Prozessen durchführen. 
% 
% Bei Verdacht auf \I{Skrotalhernien} sollte eine Auskultation
% auf Darmgeräusche im Skrotum durchgeführt werden. 
% 
% Eine \I{Varikozele} kann häufig bereits prima vista
% diagnostiziert werden. Der Skrotalinhalt wird auf den Nachweis eines vergrößerten plexus
% pampiniformis im Stehen und im Liegen palpiert. Der Befund wird
% zusätzlich unter
% \I[Valsalva!Varikozele]{Valsalva}-Druckbedingungen erhoben.
% }{}{}{}{}



\TextSlide{Rektale Untersuchung}{%
Die rektale Untersuchung ist ein
Standardbestandteil des physikalischen Status und sollte als solche bei
stationärer Aufnahme durchgeführt werden.
}{
\begin{PkuBefunde}[]
  \item[\SymbolPkuNormal] \PkuBefund{Digital-rektale
  Untersuchung: Analregion unauffällig, Ampullenschleimhaut
  und Prostata palp.
  unauffällig, Ampulle frei, Hämoccult negativ}
%   \item[\SymbolPkuVariante] \PkuBefund{}
%   \item[\SymbolFindingPathological] \PkuBefund{}
\end{PkuBefunde}

Keine Inguinalhernie
  oder ing. Lymphknoten palpabel; }{}{}{}

\minisec{Material}

\noindent\Parallel
{\begin{itemize}
  \item Nierentasse
  \item Gleitmittel
  \item \EE{3} (!) Untersuchungshandschuhe
  \item Abwischtücher
  \item Haemoccult (Test zum Nachweis von okkultem Blut im Stuhl)
  \item Unterlage
\end{itemize}}
{\includegraphics[width=\linewidth]{./images/pku/Rektale-Untersuchung}}





\clearpage
\begin{WidePar}
\minisec{Durchführung}
\begin{multicols}{2}
\begin{enumerate}
  \item \EE{Kontaktaufnahme} mit dem Patienten
  \begin{itemize}
    \item Hier gilt ganz besonders, dass die
    \EE{Intimsphäre} des Patienten gewahrt bleibt, also Abdecken/Entblößen nur soweit erforderlich und
    Besucher hinausschicken (auch im Folgenden darauf achten, dass niemand
    stört).
    \item Vergessen Sie nicht: Für Sie wird diese Tätigkeit bald zur Routine,
    für den Patienten stellt es aber evtl. eine
    \EE{Ausnahmesituation} dar, die ihm Angst macht, in der
    Grenzen überschritten und vielleicht Schmerzen zugefügt werden.
  \end{itemize}
  \item \EE{Trennwand} aufstellen
  \item \EE{Patientenlagerung}
  \begin{itemize}
    \item Patient entkleidet sein Gesäß und
    nimmt Linksseitenlage mit zu der Brust angezogenen Knien
    ein; das Gesäß soll sich an der Bettkante befinden
    \item Zum Patienten ans Bett setzen
  \end{itemize}
  \item \EE{Anlegen} der
  \EE{Handschuhe}, an der Untersuchungshand 2 Stück
  \item \EE{Inspektion der Analregion}
  \begin{itemize}
    \item in Ruhe und beim Pressen
    \item[\eye] Stuhlreste, Condylome, Hämorrhoiden,
    Marisquen\footnote{\I{Marisquen}: zu Hautlappen
    reorganisierte Perianalthrombosen}, Fissuren, Ekzeme \dots
  \end{itemize}
  \item Befeuchtung des Untersuchungsfingers mit
  \EE{Gleitmittel}
  \item \EE{Palpation des äußeren Sphinkters}
  \begin{itemize}
    \item Untersuchungsfinger ruhig auf den Anus legen
    \item Warten bis äußerer Sphinkter entspannt, Patient
    zum leichten Pressen auffordern
    \item Erst jetzt mit drehender Bewegung gefühlvoll
    eindringen
    \item {Patient beobachten: Schmerzen?}
  \end{itemize}
  \item \EE{Palpation der Ampulla recti} / Darmschleimhaut,
  
  erst dorsal, dann ventral
  \begin{itemize}
    \item Normaler Befund: weich, verschieblich, keine tastbaren
    Geschwulste oder Ulcerationen
  \end{itemize}
  \item \EE{Prostata} palpieren
  \begin{itemize}
    \item {Walnuss- bis rosskastaniengroß, prallelastisch,
    zweigelappt}
    \item {Achten auf:}
    \begin{compactitem}
      \item fehlende Abgrenzbarkeit, höckrige Oberfläche oder isolierter Knoten,
      härtere Konsistenz (Karzinom)
      \item Fluktuation (Abszess)
      \item Druckschmerz (Entzündung)
    \end{compactitem}
  \end{itemize}
  \item \EE{Portio (Zervix Uteri)}
  \begin{itemize}
    \item {Bei entzündlichen Prozessen kann ein Portioschiebeschmerz ausgelöst
    werden (z.\,B. Adnexitis).}
  \end{itemize}
  \item Untersuchung abschließen, \EE{Inspektion} des
  Fingerlings
  \begin{itemize}
    \item {Teerstuhl? Blut?}
  \end{itemize}
  \item \EE{Haemoccult} durchführen
  \begin{itemize}
    \item {Stuhlprobe von unterschiedlichen Stellen des Untersuchungsfingers auf
    Untersuchungsfelder auftragen}
    \item Briefchen schließen
  \end{itemize}
  \item Anus des Patienten abwischen
  \begin{itemize}
    \item weitere Abwischtücher beim Pat. lassen
  \end{itemize}
  \item Ausziehen der Handschuhe
  \item Patient zudecken
  \item Beendigung des Patientenkontakts
  \begin{itemize}
    \item {Bei Verdacht auf abnormalen Befund den Patienten keinesfalls
    verunsichern}
  \end{itemize}
  \item {kontaminationsfreie Materialentsorgung}
  \item {Haemoccult auf Station auswerten (Vorgehen laut Herstellerangabe)}
  \item {Verdacht auf abnormalen Befund dem diensthabenden Stationsarzt
  mitteilen}
  \item {Händedesinfektion}
\end{enumerate}
\end{multicols}
\end{WidePar}

\noindent\Parallel
{Die Angabe von Tastbefunden erfolgt gemäß dem
Uhrschema (siehe Zeichnung).} {\includegraphics[width=\linewidth]{./images/pku/Rektal-Uhr}}



\subsection{Blickdiagnosen}

\Layout{57}{}{%

}{}
{%
\begin{tabularx}{\linewidth}{XXX}
\includegraphics[width=\linewidth]{./images/pku/clean-img131}
&
\includegraphics[width=\linewidth]{./images/pku/clean-img132.jpg}
&

\\
\end{tabularx}
}
{Aszites}
{Quellen der Abb.: nicht eruierbar}







%%%%%%%%%%%%%%%%%%%%%%%%%%%%%%%%%%%%%%%%%%%%%%%%%%%%%%%%%%%
%%%%%%%%%%%%%%%%%%%%%%%%%%%%%%%%%%%%%%%%%%%%%%%%%%%%%%%%%%%
%%%%%%%%%%%%%%%%%%%%%%%%%%%%%%%%%%%%%%%%%%%%%%%%%%%%%%%%%%%
%%%%%%%%%%%%%%%%%%%%%%%%%%%%%%%%%%%%%%%%%%%%%%%%%%%%%%%%%%%
%%%%%%%%%%%%%%%%%%%%%%%%%%%%%%%%%%%%%%%%%%%%%%%%%%%%%%%%%%%
\clearpage
\section
[Wochenthema \CO{[E]}]
{Wochenthema \CC{[E]}}

\HeadingSub{Extremitäten/Gelenke, Neurostatus}

% Das Wochenthema D kann nach Übereinkunft mit dem Kursleiter bereits
% während der ersten drei Kurswochen mit geübt werden z.\,B.:
% 
% \begin{itemize}
% \item {Exremitäten bei Wochenthema A}
% \item Pulstatus bei Wochenthema B
% \item {Neurostatus bei Wochenthema C}
% \end{itemize}

\subsection{Vorbereitung des Wochenthemas}

Im \enquote{kleinen Neurostatus} orientieren wir uns am
3\,$\times$\,3\,$+$\,3 Schema des Line-Elements
\enquote{Neurostatus} im 6. Semester. 
Als Grundlage für das Famulaturpropädeutikum beschränken wir
uns auf einen Teilbereich der  ersten 6 Untersuchungsgänge
(der ersten 2 Dreiergruppen), die hauptsächlich zur
differentialdiagnostischen Abklärung einer Parese dienen.
Zur Durchführung der Untersuchung gibt es ein
\EE{Lehrvideo} von Prof. Müller:
\begin{itemize}\small
  \item Teil 1 -- Kraft, Tonus, Trophik:
  \\
  \url{http://www.meduniwien.ac.at/Neurologie/MCW/Folien/le_ns_1.wmv}
  \item Teil 2 -- Motorik/Feinmotilitöt, Sehnenreflexe,
  Pyramidenzeichen:
  \\
  \url{http://www.meduniwien.ac.at/Neurologie/MCW/Folien/le_ns_2.wmv}
  \item Teil 3 -- Zerebelläre Zeichen, Frontale Zeichen,
  Bewegungsversuche:
  \\
  \url{http://www.meduniwien.ac.at/Neurologie/MCW/Folien/le_ns_3.wmv}
  \item Teil 4 -- Meningismus, Laseque, Hirnnerven:
  \\
  \url{http://www.meduniwien.ac.at/Neurologie/MCW/Folien/le_ns_4.wmv}
  \item Teil 5 -- Sensibilität:
  \\
  \url{http://www.meduniwien.ac.at/Neurologie/MCW/Folien/le_ns_5.wmv}
  \item Teil 6 -- Orientierende Neuropsychologische
  Untersuchungen:
  \\
  \url{http://www.meduniwien.ac.at/Neurologie/MCW/Folien/le_ns_6.wmv}
\end{itemize}
Wir bitten Sie,
dieses Video (Kapitel 1 und 2) sowie den unten folgenden
Untersuchungsgang vor Besuch des entsprechenden
Praktikumstags gründlich zu studiereren.

\clearpage
\subsection{Übersicht: Untersuchungsgang}


\Text{Beginn}
{\begin{itemize}
  \item Hygienische Händedesinfektion
  \item {Untersucher stellt sich vor}
  \item {Patient macht sich frei}
\end{itemize}}
{}
{}
{}
{}



\Text{\CO{[E1]} Inspektion der Extremitäten und grobe Funktionsprüfung der Gelenke}
{
\begin{itemize}
  \item \EE{Inspektion Skelett und
  Gelenke} (im Stehen: Fehlstellungen, Varizen, Ödeme)
  \item \EE{Grobe Funktionsprüfung der Gelenke}
  (Beweglichkeit, seitengleich)
  \item \EE{Inspektion und Palpation der Haut} (v.\,a.
  untere Extremitäten) (inklusive Temperaturvergleich, art./venöse
  Ulzera, Ödeme, Varizen, Thrombosen)
\end{itemize}
}
{}
{}
{}
{}

\Text{\CO{[E2]} Neurostatus}
{
\begin{itemize}
  \item \TextCircled{1} \EE{Kraftprüfung} (Beugung und
  Streckung gegen Kraft des Untersuchers oder Schwerkraft)
  \begin{itemize}
    \item \EE{Vorstrecken} der Arme (mit geschlossenen Augen)
    \item \EE{Positionsversuch} (mit geschlossenen
    Augen)

\begin{multicols}{2}[\item \EE{Kennmuskel} der
\EE{oberen Extremität}]\begin{compactitem}
  \item Faustschluß
  \item Fingerspreizen
  \item Beugen im Handgelenk
  \item Strecken im Handgelenk\columnbreak
  \item Beugen im Ellenbogen
  \item Stecken im Ellenbogen
  \item Abduktion im Schultergelenk
  \item Adduktion im Schultergelenk
\end{compactitem}\end{multicols}
\begin{multicols}{2}[\item \EE{Kennmuskel} der \EE{unteren
Extremität}]\begin{compactitem}
  \item Plantarflexion der Zehen
  \item Dorsalextension der Zehen
  \item Plantarflexion der Füße
  \item Dorsalextension der Füße\columnbreak
  \item Streckung im Kniegelenk
  \item Beugung im Kniegelenk
  \item Abduktion im Hüftgelenk
  \item Adduktion im Hüftgelenk
\end{compactitem}
\end{multicols}
  \end{itemize}
  \item \TextCircled{2} \EE{Muskeltonus}
  \begin{itemize}
    \item Untersuchung der OE (Ellenbogen und Handgelenk)
\item Untersuchung der UE (Knie und Hüftgelenk)
  \end{itemize}
  \item \TextCircled{3} \EE{Trophik}
  \begin{itemize}
    \item \EE{Inspektion der Muskulatur} (seitengleich, Atrophie, Kontraktionen)
\item Thenar und Hypothenarmusklen im SV
  \end{itemize}
  \item \TextCircled{4} \EE{Feinmotorik} 
  \begin{itemize}
    \item Fingertapping
  \end{itemize}
  \item \TextCircled{5} \EE{Sehnenreflexe}
  \begin{compactitem}
    \item {Bizepssehnenreflex}
    \item {Radiusperiost-Reflex}
    \item Trizepssehnenreflex
    \item Patellarsehnenreflex
    \item Patellarsehnenreflex mit
    Jendrassik'schem-Handgriff
    \item Achillessehnenreflex
  \end{compactitem}
  \item \TextCircled{6} \EE{Pyramidenbahnzeichen} 
  \begin{itemize}
    \item Babinski
  \end{itemize}
  \item \EE{Nervendehnungsschmerz} (Lasegue und Kernig
  bds.)
  \item Überprüfung eines möglichen \EE{Meningismus}
  (im Liegen Kopf anheben, verschiedene Richtungen)
\end{itemize}
}
{}
{}
{}
{}

\Text{Ende}
{
\begin{itemize}
  \item {Patient kleidet sich wieder an}
  \item {Offene Fragen beantworten}
  \item {Verabschiedung}
  \item Hygienische Händedesinfektion
\end{itemize}
}
{}
{}
{}
{}

\bigskip
\subsection{\CO{[E1]} Inspektion der Extremitäten und grobe Funktionsprüfung der Gelenke}

\Text{Inspektion Skelett und Gelenke}{%
Im Rahmen der Inspektion der Extremitäten sollte man besonders auf
Fehlstellungen, Achsenabweichungen, Längenunterschiede (z.\,B. nach
Hüftoperationen), Schwellungen im Bereich der Gelenke (z.\,B.
Gelenksergüsse), Hämatome, Ödeme etc. achten. Erkrankungen des
Hüft-, Knie- oder Sprunggelenks verursachen Hinken und Schonhaltung
eines Beines. An den oberen Extremitäten sind Veränderungen im
Rahmen von rheumatischen Erkrankungen oder Sehnenverkürzungen (z.\,B.
\I{Dupuytren'sche Kontraktur}, schnellende Finger) rasch
erkennbar. Bei Schmerzen im Bereich der Hände wird der Patient Sie
schon bei der Begrüßung vor zu starkem Händedruck warnen.

Man sollte den Patienten beim Ausziehen der Kleidung, sowie bei
einfachen Bewegungen beobachten, dabei wirken sich schmerzhafte
Bewegungseinschränkungen besonders aus.
}{}{}{}{}

\Text{Grobe Funktionsprüfung der Gelenke (Selbststudium)}{%
Die Beweglichkeit aller wichtigen Gelenke sollte aktiv (der Patient
bewegt das Gelenk selber) und passiv (der Untersucher bewegt das
Gelenk) geprüft werden.

\begin{itemize}
\item {Schulter-, Ellbogen-, Handgelenke}
\item {Hüft-, Knie-, Sprunggelenke}
\end{itemize}

\subparagraph{Beurteilung der Beweglichkeit}

\begin{itemize}
  \item {Flexion -- Extension}
  \item {Abduktion -- Adduktion}
  \item {Innen- \& Außenrotation}
\end{itemize}

Weiterführende Gelenksuntersuchungen werden im unfallchirurgischen und
orthopädischen Status durchgeführt.
}{}{}{}{}

\Text{Inspektion und Palpation der Haut (v.\,a. untere Extremitäten)}{% 
Es fallen minder durchblutete Extremitäten durch ihre
Hautverfärbung und unterschiedliche Temperatur auf
(Seitenvergleich!). Weiße oder bläulich verfärbte Finger
oder Zehen weisen auf eine Minderperfusion hin. Weiters
sollte man auf Varizen, Ödeme und Ulzera achten.

\EE{Ulzera} arteriellen Ursprungs finden sich im Gegensatz
zu venösen immer ganz peripher und schreiten im weiteren
Verlauf nach proximal fort.

\EE{Extremitätenödeme} werden im Seitenvergleich beurteilt,
Beinödeme am besten am Knöchel und prätibial. In subkutanen
Flüssigkeitsansammlungen hinterlässt der drückende Finger
eine Delle, die sich erst langsam wieder zurückbildet.
Lymphödeme weisen dieses Phänomen nicht auf. Ödeme bilden
sich bei Herzinsuffizienz in den tieferliegenden
Körperpartien (lageabhängig - meist Beine).

Bei Verdacht auf eine \EE{tiefe Venenthrombose} muss die
gesamte Extremität untersucht werden. Man achtet auf Differenzen des
Umfangs, Temperatur- und Farbunterschiede und palpiert auf
eine vermehrte Konsistenz, Verhärtungen und Druckschmerz.
}{}{}{}{}



\bigskip
\subsection{\CO{[E2]} Neurostatus}

\TODO{GABS}{2011-02-23}{Fehlt: Verwendung eines
Reflexhammers}{}

\subsubsection{Allgemeines und Klinik}

\Text{Der Status \enquote{3 $\times$ 3 $+$ 3}}
{{\ldots} ist eine mögliche (und hier verwendete)
Ausgangsbasis für die Herangehensweise für die neurologische
Untersuchung. Dieses System zielt primär darauf ab, Paresen
zu erkennen und sie zu zentralen oder peripheren Störungen
zuzuordnen, sowie die wichtigsten neurologischen Symptome
zu erkennen.

Bei der Untersuchung der groben \TextCircled{1} Kraft wird
festgestellt, ob  und in welchem Areal eine Parese vorliegt. Die
Untersuchungen \TextCircled{2} Tonus, \TextCircled{3} Trophik, \TextCircled{4} Feinmotilität,
\TextCircled{5} Sehnenreflexe und  \TextCircled{6}
Pyramidenzeichen  helfen dann bei der Abklärung, ob die Lähmung zentralen oder peripheren
Ursprungs ist.

% im Kopf, nicht auf dem Vordruck!

\begin{table}[H]
\FormatTable
\begin{tabularx}{\linewidth}{X|X|X}\toprule\addlinespace
\TextCircled{1} \EE{K}raft
&
\TextCircled{2} \EE{To}nus
&
\TextCircled{3} \EE{T}rophik
\\\addlinespace\midrule\addlinespace
\TextCircled{4} \EE{F}ein\EE{mot}ilität
&
\TextCircled{5} \EE{S}ehnen\EE{r}eflexe
&
\TextCircled{6} \EE{Py}ramiden\EE{z}eichen
\\\addlinespace\midrule\addlinespace
\TextCircled{7} Zerebelläre Zeichen
&
\TextCircled{8} Frontale Zeichen
&
\TextCircled{9} Bewegungsversuche
\\\addlinespace\toprule\addlinespace
Meninigismus, Lasègue
&
Hirnnerven
&
Sensibilität
\\\addlinespace\bottomrule
\end{tabularx}
\end{table}

\subparagraph{Im Detail:} 
Die ersten 2~$\times$~3 (Kraft, Tonus, Trophik;
Feinmotorik, Sehnenreflexe, Pyramidenzeichen) sollen das
\EE{Hauptsymptom} suchen und topographisch zuordnen: Die
\EE{Lähmung/Parese}. Diese muss zuallererst definiert
werden: 

\begin{itemize}
  \item Welche Muskeln sind schwach? (\EE{Seitenvergleich!})
  \item Kraftgrad?
  \item Zentrale oder periphere Parese?
  
  Zuordnung durch To, T, FMOT, SR, PYZ
  \item Muster (Tetra-, Para-, Hemi-, Monoparese =
  4\,--\,2\,--\,1-Extremitäten oder einzelne musekeln?)
\end{itemize}

\subparagraph{Paresen}
Als Paresen werden Störungen der Motorik einzelner Muskeln
bzw. Muskelgruppen bezeichnet, die mit Lähmungen verbunden
sind. Man unterscheidet \EE{spastische} und \EE{schlaffe}
\EE{Lähmungen} (erhöhter bzw. verringerter Muskeltonus). Bei
vollständigem Funktionsausfall großer Muskelgruppen spricht man von Plegie oder Paralyse (Hemiplegie,
Tetraplegie).
}
{}
{}
{}
{}






\begin{table}[H]
\FormatTable
\begin{tabularx}{\linewidth}{XXX}\toprule\addlinespace
\TabHeading{DD}
&
\TabHeading{zentrale vs.}
&
\TabHeading{periphere Parese}
\\\addlinespace\midrule\addlinespace
\TabSub{Lokalisation}
&
(Mono-), Hemi-, Para-, Tetraparese
&
jeweils innervierte Muskeln
\\\addlinespace
\TabSub{Tonus}
&
spastisch, pseudoschlaff
&
schlaff
\\\addlinespace
\TabSub{Trophik}
&
Inaktivitätsatrophie
&
umschriebene Atrophie
\\\addlinespace
\TabSub{Feinmotorik}
&
gestört
&
gestört
\\\addlinespace
\TabSub{Sehnenreflexe}
&
gesteigert
&
abgeschwächt, fehlend
\\\addlinespace
\TabSub{Pyramidenzeichen}
&
\EE{positiv}
&
\EE{negativ}
\\\addlinespace\bottomrule
\end{tabularx}
\end{table}

Die Zeichen einer zentralen Störung benötigen z.\,T. eine
gewisse Zeit um sich zu entwickeln, im akuten Stadium
gleicht daher eine zentrale Parese teilweise einer
peripheren.


% \Text{Kraftgrade}
% {~
% 
% }
% {}
% {}
% {}
% {}

\subparagraph{Seitenvergleich} Der Seitenvergleich ist im
Rahmen der neurologischen Untersuchung sehr wichtig.
Besonders bei Untersuchungen, deren Befunde von vornherein
sehr individuell sind (Reflexantwort, \ldots), liefert er
einen guten Orientierungs- und Referenzpunkt. Grundsätzlich
sollen Untersuchungen wie z.\,B. die Kraftprüfung und die
Sehnenreflex-Untersuchung immer im \EE{direkten
Seitenvergleich} vorgenommen werden!


\subsubsection{Untersuchungsgang}
\paragraph{Allgemeines zum Untersuchungsgang}

Alle Untersuchungen werden liegend, beim bis auf die
Unterwäsche entkleideteten Patienten im Seitenvergleich
durchgeführt.






% \Text{5 Sehnenreflexe (gesteigert, abgeschwächt)}
% {5a Bizepssehenreflex
% 5b Radiusperiostrefelex
% 5cTrizepssehenreflex
% 5d Patellarsehenenreflex
% 5e Patellarsehnen mit Jendrasek
% 5f Achillessehnenreflex}
% {}
% {}
% {}
% {}
% 
% 
% 
% \Text{6 Pyramidenbahnzeichen}
% {6a Trömner
% 6b Babinski}
% {}
% {}
% {}
% {}
% 
% 
% 
% \Text{7 Nervendehnungsschmerzen}
% {7aMeningismus
% 7b Lasegue}
% {}
% {}
% {}
% {}











% \Text{Rigor}{%
% Rigor ist eine besondere Form der Tonuserhöhung, der in
% Beugern und Streckern gleichermaßen auftritt. Es handelt
% sich dabei um eine dauerhafte Widerstandserhöhung, die bei
% der passiven Prüfung ruckartig nachgibt und abgehackt
% erscheint: \IX{Zahnradphänomen} 
% (ein typisches Zeichen bei \I[Morbus!Parkinson]{M.
% Parkinson}).
% 
% Bei der passiven Beweglichkeitsprüfung legt der Untersucher
% die Finger einer Hand z.\,B. in die Ellenbeuge und bewegt
% den Unterarm des Patienten mit der anderen Hand. Dabei lässt
% sich das ruckartige Bewegen der Sehnen der Beugemuskulatur
% gut palpieren.
% }{}{}{}{}
% 
% 
% \Text{Tremor}{%
% Als Tremor bezeichnet man rhythmische Zuckungen
% antagonistisch wirkender Muskeln.
% 
% Man differenziert Amplitude und Dauer.
% 
% \begin{itemize}
%   \item Amplitude: grobschlägig -- feinschlägig
%   \item Dauer: ständig, in Ruhe, bei Intention
% \end{itemize}
% 
% Ein Neuauftreten bzw. Zunahme bei zielgerichteten Bewegungen wird bei
% zerebellären Störungen gefunden: Intentionstremor.
% 
% Bei älteren Personen besteht häufig ein feinschlägiger Ruhetremor,
% der bei absichtlichen Bewegungen verschwindet: physiologischer
% Alterstremor.
% 
% Einen feinschlägigen Tremor findet man auch bei Alkoholkrankheit.
% }{}{}{}{}
% 
% 
% 




\Text{\TextCircled{1} Kraftprüfung (Beugung und
Steckung gegen Kraft des Untersuchers oder Schwerkraft)} 
{Vgl. Lehrvideo.

Geprüft wird die Kraftentwicklung einzelner Muskeln bzw.
Muskelgruppen. Der Patient soll Bewegungen in
Funktionsrichtung des jeweils untersuchten Muskels gegen den
Widerstand des Untersuchers mit möglichst maximaler Kraft
ausführen. Es muss darauf geachtet werden, dass die
einzelnen Bewegungen \EE{isoliert} untersucht werden, d.\,h.
die angrenzenden Knochen/Gelenke müssen ggfs. durch den
Untersucher fixiert werden.

Dem Patienten müssen klare Anweisungen gegeben werden; auf
die Formulierung dieser Anweisungen soll daher im Praktikum
geachtet werden.

Der \EE{Seitenvergleich} ist für die Beurteilung wichtig.
Zur Beschreibung der Kraft wird die Einteilung in \T{Kraftgrade} verwendet:

\noindent\begin{minipage}{\linewidth}
\FormatTable
\begin{tabularx}{\linewidth}{lX}\toprule\addlinespace
\multicolumn{2}{c}{\TabHeading{Kraftgrade}}
\\\addlinespace\midrule\addlinespace
0 & keine Reaktion
\\
1 & Muskeln spannen sich an, aber keine Bewegung möglich
\\
2 & Bewegung nach aufheben der Schwerkraft möglich
\\
3 & Bewegung gegen die Schwerkraft möglich
\\
4 & Bewegung gegen leichten Widerstand möglich
\\
5 & Bewegung gegen erheblichen Wiederstand (Untersucher muss sich anstrengen) möglich
\\\addlinespace\bottomrule
\end{tabularx}
\end{minipage}

\bigskip

\noindent\begin{minipage}{\linewidth}
\CorrectionItemizeBox

\begin{itemize}
  \item \EE{Vorstrecken} der Arme
  (\I{Armvorhalteversuch}, \I{AVV}; mit geschlossenen
  Augen, gegen die Schwerkraft)
  \item \II{Positionsversuch} (\I{PV}) der Beine (mit
  geschlossenen Augen)
  \begin{multicols}{2}[\item \EE{Kennmuskel der
  \E{oberen Extremität}}]
  \begin{compactitem}
    \item Faustschluß
    \item Fingerglieder strecken
    \item Fingerspreizen
    \item Beugen im Handgelenk
    \item Stecken im Handgelenk\columnbreak
    \item Beugen im Ellenbogen
    \item Strecken im Ellenbogen
    \item Abduktion im Schultergelenk
    \item Adduktion im Schultergelenk
  \end{compactitem}\end{multicols}
  \begin{multicols}{2}[\item \EE{Kennmuskel der \E{unteren
  Extremität}}]
  \begin{compactitem}
    \item Plantarflexion der Zehen
    \item Dorsalextension der Zehen
    \item Plantarflexion der Sprunggelenke
    \item Dorsalextension der Sprunggelenke\columnbreak
    \item Streckung im Kniegelenk
    \item Beugung im Kniegelenk
    \item Abduktion im Hüftgelenk
    \item Adduktion im Hüftgelenk
  \end{compactitem}
  \end{multicols}
\end{itemize}
\end{minipage}


}{}{}{}{}


\begin{PictureSeries}{}{}
\PictureSeriesRow{2}
{./images/pku/clean-img142.jpg}
{}
{./images/pku/clean-img143.jpg}
{}
{}
{}
{}
{}
\end{PictureSeries}

 

\Text{\TextCircled{2} Muskeltonus}
{Vgl. Lehrvideo.

\Cave{Im Rahmen der Tonusprüfung werden die Gelenke oft
stark beansprucht!}

Man beurteile: Vorhanden, abgeschwächt, schlaff, fehlend,
Spastik?, Rigor?

% Schnelle Bewegungen lassen eher eine Spastik erkennen,
% langsame einen Rigor.

\begin{itemize}
  \item Untersuchung der OE (Ellenbogen und
  Handgelenk)
  \item Untersuchung der UE (Knie und Hüftgelenk)
\end{itemize}
}
{}
{}
{}
{}






\Text{\TextCircled{3} Trophik}
{Vgl. Lehrvideo.

Man beurteile: Atrophie, Faszikulationen, Hautveränderungen.

Schonhaltungen oder Lähmungen (\I{Paresen}) führen sehr
rasch an der betroffenen Extremität zu Muskelatrophie. Man beurteilt die
Muskulatur daher immer im \EE{Seitenvergleich} (ggfs.
Umfangbestimmung mit Maßband). Bei länger währender
Immobilität können \I{Beugekontrakturen} (Verkürzung der
Muskeln durch Atrophie) entstehen.

An der Hand gelten den \EE{Thenar-} und
\EE{Hypothenarmusklen} besondere Beachtung.
}
{}
{}
{}
{}



\Text{\TextCircled{4} Feinmotorik}
{\I{Diadochokinese}: Fähigkeit, rasch
aufeinander folgende Bewegungen, beispielsweise Pronation
und Supination, wiederholt auszuführen. Fehlt diese
Fähigkeit, spricht man von einer
\I[Adiadochokinese]{Adiadochokinese}, bei einer alleinigen Verlangsamung der Bewegungen von einer
\I{Bradydiadochokinese}. Ist die Bewegung darüber
hinausgehend eingeschränkt, wird dies als
\I{Dysdiadochokinese} bezeichnet.
Vgl. Lehrvideo. 

\begin{itemize}
  \item Fingertapping
  \item Klavier spielen
  \item Hände rasch abwechseln im und gegen den
  Uhrzeigersinn drehen (\enquote*{(Nicht-EU-konforme)
  Glühbirne ein- und ausschrauben})
\end{itemize}
}
{}
{}
{}
{}

\Text{\TextCircled{5} Sehnenreflexe}
{Da die Reflexe individuell sehr unterschiedlich ausgeprägt
sind, gibt der \EE{Seitenvergleich} einen wichtigen
Referenzpunkt um Auffälligkeiten abzuschätzen. 
Die Übersicht \FullPrettyRef{fig:bilderserie-reflexe}, gibt
eine Übersicht über die wichtigsten peripheren Reflexe.

Bei schwachen Reflexen kann man sich mit dem
\T[Jendrassik-Handgriff]{Jendrassik'schen-Handgriff}
behelfen: Der Jendrassik-Handgriff ist ein nach dem
ungarischen Neurologen Ernö Jendrassik (\VonBis{1858}{1921})
benanntes Untersuchungsmanöver.
Der Patient greift dabei mit den Händen ineinander und zieht
kräftig mit beiden vor dem Oberkörper angewinkelten Armen.
Durch die muskuläre Anspannung der Arme kommt es durch
Bahnung zu einer \E{Steigerung der Reflexerregbarkeit}.
% 
% \includegraphics[width=0.45\linewidth]{./images/pku/clean-img149.jpg}
\bigskip
}{}
{./images/pku/clean-img149.jpg}
{Jendrassik-Handgriff}
{}


\begin{PictureSeriesWide}{}{Bilderserie:
Reflexe\label{fig:bilderserie-reflexe}}

\PictureSeriesRowWide{3}
{./images/pku/clean-img145.jpg}
{\TabSub{Bizepssehnenreflex}:
Sehne mit Zeigefinger gut tasten. Hammerschlag auf den in
der Ellenbeuge des Patienten gelegten Zeigefinger.
\E{Ergebnis: Beugung im Ellbogengelenk.}}
{./images/pku/clean-img144.jpg}
{\TabSub{Radius-Periost-Reflex}:
Der Unterarm wird gebeugt mit den Handflächen nach unten
gelagert. 
Einen Finger ca. 5\CM* oberhalb des Handgelenks auf den
Radius legen und mit dem Reflexhammer auf den Finger
klopfen. 
\E{Ergebnis: Normal ist eine leichte Beugebewegung des
Unterarms.}}
{./images/pku/clean-img146.jpg}
{\TabSub{Trizepssehnenreflex} 
Bei angewinkeltem Unterarm und abgewinkeltem Oberarm wird auf
die Sehne des
\E{Musculus triceps brachii}
kurz oberhalb des \E{Olekranons} ein Schlag
ausgeführt.  
\E{Ergebnis: Streckung des Unterarms.}}
{}
{}
\PictureSeriesRowWide{3}
{./images/pku/clean-img147.jpg}
{\TabSub{Patellarsehnenreflex}:
Bei leicht gebeugtem Bein erfolgt ein Schlag auf die Sehne des M.
quadriceps femoris unterhalb der Patella.
\E{Ergebnis: Ruckartige Streckung des Beines im Kniegelenk.}}
{./images/pku/Jendrassik-cropped}
{\TabSub{Patellarsehnenreflex mit Jendrassik'schem Handgriff}:
Der Patient greift mit den Händen ineinander und zieht
kräftig mit beiden vor dem Oberkörper angewinkelten Armen.}
{./images/pku/clean-img148.jpg}
{\TabSub{Achillessehnenreflex}:
Schlag auf die Achillessehne ca. 5\CM* oberhalb des
Fersenbeins bei abduziertem und abgewinkeltem Bein.
\E{Ergebnis: leichte Plantarflexion des Fußes im
Sprunggelenk.}}
{}
{}
\end{PictureSeriesWide}
\index{Musculus!triceps brachii}\index{Olekranon}

% \Layout{57}{}{%
% 
% }{}
% {%
% \begin{tabularx}{\linewidth}{XXX}
% \includegraphics[width=\linewidth]{./images/pku/clean-img145.jpg}
% &
% \includegraphics[width=\linewidth]{./images/pku/clean-img144.jpg}
% &
% \includegraphics[width=\linewidth]{./images/pku/clean-img146.jpg}
% \\
% \TabSub{Bizepssehnenreflex}   
% 
% Sehne mit Zeigefinger gut tasten. Hammerschlag auf den in
% der Ellenbeuge des Patienten gelegten Zeigefinger.
% 
% Ergebnis: Beugung im Ellbogengelenk.
% &
% \TabSub{Radius-Periost Reflex} 
% 
% Der Unterarm wird gebeugt mit den Handflächen nach unten
% gelagert.
% 
% Einen Finger ca. 5\CM* oberhalb des Handgelenks auf den
% Radius legen und mit dem Reflexhammer auf den Finger
% klopfen. 
% 
% Ergebnis: Normal ist eine leichte Beugebewegung des
% Unterarms.
% &
% 
% \TabSub{Trizepssehnenreflex}
% 
% Bei angewinkeltem Unterarm und abgewinkeltem Oberarm wird auf
% die Sehne des
% \I{Musculus triceps brachii} kurz oberhalb des \I[Olekranon]{Olekranons}
% ein Schlag ausgeführt.
% 
% Ergebnis: Streckung des Unterarms.
% \\\addlinespace
% \includegraphics[width=\linewidth]{./images/pku/clean-img147.jpg}
% &
% \includegraphics[width=\linewidth]{./images/pku/Jendrassik-cropped}
% &
% \includegraphics[width=\linewidth]{./images/pku/clean-img148.jpg}
% \\
% \TabSub{Patellarsehnenreflex}
% 
% Bei leicht gebeugtem Bein erfolgt ein Schlag auf die Sehne des M.
% quadriceps femoris unterhalb der Patella.
% 
% Ergebnis: Ruckartige Streckung des Beines im Kniegelenk.
% &
% \TabSub{Patellarsehnenreflex mit Jendrassik'schem Handgriff}
% 
% Der Patient greift mit den Händen ineinander und zieht
% kräftig mit beiden vor dem Oberkörper angewinkelten Armen.
% &
% \TabSub{Achillessehnenreflex}
% 
% Schlag auf die Achillessehne ca. 5\CM* oberhalb des
% Fersenbeins bei abduziertem und abgewinkeltem Bein.
% 
% Ergebnis: leichte Plantarflexion des Fußes im Sprunggelenk.
% \end{tabularx}
% }
% {Bilderserie: Reflexe\label{tab:bilderserie-reflexe}}
% {DEMAW}

\Text{\TextCircled{6} Pyramidenbahnzeichen}{%
Eine akute Pyramidenbahnläsion im Gehirn oder Rückenmark
  (z.\,B. bei Schlaganfall) führt zu schlaffen Lähmungen, zu
  gesteigerten und pathologischen Reflexen. Ein klassischer
  Vertreter dieser Reflexe ist der \T{Babinski-Reflex}.

\subparagraph{Babinski-Reflex} Er ist ein sicheres
Zeichen bei zentraler Lähmung (bei Säuglingen ist dieser
Reflex allerdings physiologisch positiv).


Es wird mit einem zugespitzten Gegenstand (z.\,B.
Reflexhammerhandgriff, Daumennagel, Kugelschreiber etc.) kräftig entlang
der lateralen Fußsohle von der Ferse zu den Zehen und dann
in Richtung großer Zehe gestrichen.

Als positiv wird die tonische \EE{Extension der Großzehe}
und \EE{Plantarflexion der restlichen Zehen} mit
\EE{Spreizphänomen} gewertet.


% % Trömner ist KEIN Pyramidenbahnzeichen!
% \subparagraph{Trömner-Reflex} Beim Schnippen oder
%   Anschlagen der palmaren Seite des Endgliedes des
%   Mittelfingers kommt es zu einer Flexion
%   der Fingerendglieder inkl. des Daumens. Ist diese
%   Reaktion deutlich gesteigert, gilt sie als Zeichen für
%   eine Schädigung der zentralen motorischen Bahnsysteme.
}{}{}{}{}

\begin{PictureSeries}{}{Babinski-Reflex}
\PictureSeriesRow{2}
{./images/pku/clean-img150.jpg}
{}
{./images/pku/clean-img151.jpg}
{}
{}
{}
{}
{}
\PictureSeriesRow{2}
{./images/pku/clean-img152}
{Babinski positiv}
{./images/pku/clean-img153}
{Babinski negativ}
{}
{}
{}
{}
\end{PictureSeries}

% \Text{Sensibilität}{%
% Die Oberflächensensibilität kann orientierend durch
% gleichzeitiges Bestreichen beider Extremitäten mit allen
% Fingerkuppen geprüft werden. Der Patient soll dabei die
% Augen geschlossen halten, um die optische Kontrolle
% auszuschließen. Man bittet den Patienten auf
% unterschiedliche Wahrnehmungen zu achten
% (Seitenungleichheit, unterschiedliche
% Empfindungsqualitäten).
% 
% Neben Störungen der Wahrnehmung kommen auch abnorme Reizerscheinungen
% vor = Parästhesien (Mißempfindungen):
% 
% 
% \begin{itemize}
%   \item Kribbeln, \enquote{Ameisenlaufen}
%   \item Prickeln, Brennen
%   \item elektrisierende Schmerzen
% \end{itemize}
% }{}{}{}{}

\Text{Nervendehnungsschmerz (Lasègue-Zeichen)}{%
Nervendehnungsschmerzen können bei Dehnung von Nerven, 
Nervenwurzeln, Rückenmark oder Meningen auftreten.

Das gestreckte Bein des liegenden Patienten wird im
Hüftgelenk passiv (!) gebeugt. 

\begin{PictureSeriesWide}{}{Lasègue-Zeichen}
\PictureSeriesRowWide{3}
{./images/pku/clean-img158.jpg}
{}
{empty}
{}
{empty}
{}
{}
{}
\end{PictureSeriesWide}


Ist das \T{Lasègue-Zeichen}
positiv, so kommt es zu Schmerzen lumbosakral, im Gesäß und
im Bereich des dorsalen Oberschenkels, welche bis zur
Kniekehle ausstrahlen\footnote{Ist der Schmerz eher diffus,
kann es sich um einen \I{Pseudo-Lasègue} handeln,
der andere Ursachen (Gelenksschäden, Verkürzung der
Muskulatur, \ldots) hat \cite{Zeiler2001}.}, die Beugung ist
nicht oder nur eingeschränkt möglich. Es muss immer der \EE{Beugungswinkel} angegeben werden, ab dem das
Lasègue-Zeichen positiv ist! \citep{Zeiler2001}
Wird das \E{Kniegelenk schmerzreflektorisch gebeugt}, so
entspricht dies dem \T{Kernig-Zeichen}.
}{}{}{}{}





% \Text{Feinmotorik (Diadochokinese)}{%
% Die Feinmotorik wird mit schnellen Wechselbewegungen durch
% abwechselnde Pronation und Supination der Hände (Glühbirnen
% ein- und ausdrehen) untersucht.
% 
% Diese Untersuchung wird mit beiden Händen gleichzeitig
% durchgeführt.
% 
% Weiters wird der Daumen-Zeigefinger-Versuch gemacht, bei dem
% eine länger dauernde Schraubbewegung (Uhr aufziehen)
% demonstriert werden soll.
% 
% Ist der Bewegungsablauf normal und rhythmisch spricht man
% von Eudiadochokinese, eine verlangsamte Bewegung nennt man
% Bradydiadochokinese, eine völlig unkoordinierte Bewegung
% Adiadochokinese.
% }{}{}{}{}
% 
% 
% 
% 
% \begin{table}[H]
% \begin{tabularx}{\linewidth}{XX}
% \includegraphics[width=\linewidth]{./images/pku/clean-img154.jpg}  
% &
% \includegraphics[width=\linewidth]{./images/pku/clean-img155.jpg} 
% \\
% \centering  Schraubbewegung von Daumen
% und Zeigefinger 
% &
% Gegensinnige
% Schraubbewegung der Hände
% \\
% \end{tabularx}
% \end{table}
% 
% 
% \Text{Koordinationsprüfung (Finger-Nase-Versuch)}{%
% Im Rahmen der Prüfung der Koordinationsfähigkeit (Kleinhirn)
% führt man den Finger-Nase-Versuch durch. Dabei soll der
% Patient mit geschlossenen Augen und von der Position des
% gestreckten Armes ausgehend mit der Zeigefingerspitze in
% einer bogenförmigen Bewegung die Nase treffen.
% 
% \begin{table}[H]
% \begin{tabularx}{\linewidth}{XX}
% \includegraphics[width=\linewidth]{./images/pku/clean-img156.jpg}
% &
% \includegraphics[width=\linewidth]{./images/pku/clean-img157.jpg}
% \\
% \end{tabularx}
% \end{table}
% 
% Alternativ kann der Knie-Ferse-Versuch durchgeführt werden.
% Dabei wird im Liegen mit der Ferse das Knie des anderen
% Beines berührt und davon ausgehend die Ferse entlang der
% Tibiakante (Schienbein) nach distal bewegt.
% }{}{}{}{}
% 
% 






\Text{Überprüfung eines möglichen Meningismus}
{Der Meningismus bezeichnet das Symptom der schmerzhaften
Nackensteifigkeit bei Reizungen und Erkrankungen der
Hirnhäute (Notfall!).

Er ist eine reflektorische Verspannung der Nackenmuskulatur
als Reaktion auf den Schmerz, bei tiefer Bewusstlosigkeit
(Koma) löst er sich wieder. Zu einem meningealen Reizsyndrom
gehören ferner Übelkeit bis zum Erbrechen, Licht- und
Geräuschempfindlichkeit (Photo-/ Phonophobie).

% Das \II{Brudzinski-Zeichen} wird am entspannt auf dem Rücken
% liegenden Patienten getestet. Es gilt als positiv, wenn bei
% passivem Vorbeugen des Kopfes reflektorisch die Beine in
% Hüft- und Kniegelenken angewinkelt werden. Es tritt
% insbesondere bei Meningitis, aber auch bei Enzephalitis oder
% einer Subarachnoidalblutung auf. 

\Cave{Die Überprüfung auf Meninigismus-Zeichen ist besonders
bei fiebernden Patienten  wichtig, um eine eventuelle
Meningitis zu erkennen!} 
}{}{}{}{}

\begin{PictureSeriesWide}{}{Meningismus}
\PictureSeriesRowWide{3}
{./images/pku/clean-img159.jpg}
{Kopf im Liegen anheben und Kinn zur Brust führen,
\EE{passive} (!) Bewegung. }
{./images/pku/clean-img160.jpg}
{Seitliche Rotation}
{./images/pku/clean-img161.jpg}
{}
{}
{}
\end{PictureSeriesWide}





\cleardoublepage
\subsection{Blickdiagnosen}

\begin{WidePar}
\begin{multicols}{3}

\noindent\begin{minipage}{\linewidth}
\includegraphics[width=\linewidth]{./images/pku/clean-img162.jpg}

\Captionof{figure}{\I{Palmarerythem}}
\end{minipage}

\noindent\begin{minipage}{\linewidth}
\includegraphics[width=\linewidth]{./images/pku/clean-img163.jpg} 

\Captionof{figure}{\I{Uhrglasnägel}}
\end{minipage}

\noindent\begin{minipage}{\linewidth}
\includegraphics[width=\linewidth]{./images/pku/clean-img164} 

\Captionof{figure}{\I{Malignes Melanom}}
\end{minipage}

\noindent\begin{minipage}{\linewidth}
\includegraphics[width=\linewidth]{./images/pku/clean-img99.jpg} 

\Captionof{figure}{\I{Erythema migrans}}
\end{minipage}

\noindent\begin{minipage}{\linewidth}
\includegraphics[width=\linewidth]{./images/pku/clean-img100.jpg}

\includegraphics[width=\linewidth]{./images/pku/clean-img101}
 
\includegraphics[width=\linewidth]{./images/pku/Herpes-zoster-3.jpg} 

\Captionof{figure}{\I{Herpes zoster}}
\end{minipage}

\noindent\begin{minipage}{\linewidth}
\includegraphics[width=\linewidth]{./images/pku/clean-img170.jpg}

\includegraphics[width=\linewidth]{./images/pku/clean-img171.jpg} 

\Captionof{figure}{\I{Psoriasis}}
\end{minipage}

\noindent\begin{minipage}{\linewidth}
\includegraphics[width=\linewidth]{./images/pku/clean-img165.jpg}

\Captionof{figure}{\I{Tiefe Beinvenenthrombose}}
\end{minipage}

\noindent\begin{minipage}{\linewidth}
\includegraphics[width=\linewidth]{./images/pku/clean-img166.jpg}

\includegraphics[width=\linewidth]{./images/pku/clean-img167.jpg} 

\Captionof{figure}{\I{Varikositas}}
\end{minipage}

\noindent\begin{minipage}{\linewidth}
\includegraphics[width=\linewidth]{./images/pku/clean-img168.jpg}

\Captionof{figure}{\I{Beinödeme}}
\end{minipage}



\noindent\begin{minipage}{\linewidth}
\includegraphics[width=\linewidth]{./images/pku/clean-img169.jpg} 

\Captionof{figure}{\I{Dekubitus}}
\end{minipage}



\noindent\begin{minipage}{\linewidth}
\includegraphics[width=\linewidth]{./images/pku/clean-img172.jpg}

\Captionof{figure}{\I{Gangrän}}
\end{minipage}

\noindent\begin{minipage}{\linewidth}
\includegraphics[width=\linewidth]{./images/pku/clean-img173.jpg} 

\Captionof{figure}{\I{Arterielles Ulcus}}
\end{minipage}

\noindent\begin{minipage}{\linewidth}
\includegraphics[width=\linewidth]{./images/pku/clean-img174.jpg}

\Captionof{figure}{\I{Venöses Ulcus}}
\end{minipage}

\noindent\begin{minipage}{\linewidth}
\includegraphics[width=\linewidth]{./images/pku/clean-img175}

\Captionof{figure}{\I{Rheumatoide Arthritis}}
\end{minipage}

\noindent\begin{minipage}{\linewidth}
\includegraphics[width=\linewidth]{./images/pku/clean-img176.jpg} 

\Captionof{figure}{\I{Dupuytren'sche Kontraktur}}
\end{minipage}
\end{multicols}
\end{WidePar}

\clearpage
\section{Anhang}

\subsection{Der Notfallpatient}

Beim Notfallpatienten gelten andere Erfordernisse als bei
der möglichst umfassenden und ausführlichen Statuserhebung.
Hier ist die Evaluation der lebenswichtigen Funktionen und
Screening auf akut lebensbedrohliche Diagnosen, sowie die
Erstellung von Verdachts- und Differentialdiagnosen zur
weiteren Abklärung vorrangig. 

Um diesen Ansprüchen gerecht zu werden existieren, je nach
Situation, verschiedene Schemata, welche z.\,T. auch in
diversen standardisierten Kurssystemen
(ATLS{\texttrademark}, ITLS{\texttrademark},
PHTLS{\texttrademark}, AMLS{\texttrademark},
ERC{\texttrademark}, \ldots) integriert sind. Basis für die
meisten Herangehensweisen ist das \T{ABC}-Schema (Airway,
Breathing, Circulation), mitunter erweitert um die
Buchstaben \I{DE} (Disability, Exposure/Environment). Die
Details, welche Untersuchungen bei den jeweiligen Punkten
durchgeführt werden sollen, unterscheiden sich dabei teils
erheblich, und sind auch vom Setting, der Ausstattung und
der Ausbildung des Personals abhängig.
In der \FullPrettyRef{tab:notfallpatient-erstuntersuchung}
ist als Beispiel das Schema der ARGE Arbeits- und
Ausbildungsstandards für den Sanitätdienst ({\ARGE})
dargestellt (das gesamte Werk ist unter \url{http://www.aass.at}
kostenlos beziehbar).

% Grundsätzlich gilt, je akut lebensbedrohlicher die Erkrankung, umso
% kürzer die Anamnese und Untersuchung. Lebensrettende
% Sofortmaßnahmen siehe die Lehrveranstaltungen
% \enquote{Erste Hilfe} und
% \enquote{Notfallmedizin}. Vgl. auch
% \cite{TertialeNotfallIntensivmed1} Hamp und Weidenauer:
% \emph{Lehrbuch Tertiale Notfall- und Intensivmedizin}.

\begin{table}[H]
\caption{Erstuntersuchung des Notfallpatienten, nach
\cite{Aass-1.0var1}}
\label{tab:notfallpatient-erstuntersuchung}

\begin{WidePar}

\CorrectionItemizeBox\renewcommand{\itemhook}{%
%   \setlength{\topsep}{0pt}%
  \setlength{\itemsep}{0.004\baselineskip}
%   \setlength{\parskip}{\baselineskip}
  \setlength{\leftmargin}{1.4em}
  \setlength{\labelwidth}{0.9em}
}
\noindent\begin{minipage}[t]{\linewidth -
{2\fboxrule} - {2\fboxsep}}
\noindent\begin{minipage}[t]{\linewidth}
{%\vspace{-2em}%
\footnotesize
% \begin{tabulary}{\linewidth}[H]{cp{8em}p{8.5cm}}
\begin{tabularx}{\linewidth}[H]{p{2cm}XX}
\toprule\addlinespace
\noindent~
&
\noindent\TabHeading{Beurteilung / Untersuchung}
&
\noindent\TabHeading{Befunde}
\\\addlinespace\toprule\addlinespace
\noindent\TabHeading{Szeneüberblick m. Schutz}
&
\begin{minipage}[t]{\linewidth}
\begin{itemize}
  \item Sicherheit, Selbstschutz; Patientenanzahl; 
  \item Umgebung, Ort,
Zeit, Gefahrenzonen, Umstände,
\item Weitere Resourcen erforderlich? Unfallmechanismus,
Großschaden
\item Lagemedlung erforderlich
\item {Trauma?} {Mechanismus?}
\end{itemize}
\end{minipage}
&
\noindent\color{MuwBlueDark}\begin{minipage}[t]{\linewidth}
\begin{itemize}
  \item[\SymbolPkuNormal] Umgebung sicher
%   \item[\SymbolPkuVariante]
%   \item[\SymbolFindingPathological] 
\end{itemize}
\end{minipage}
\\\addlinespace\toprule\addlinespace
\noindent\TabHeading{Eindruck}
&
\begin{minipage}[t]{\linewidth}
\begin{itemize}
  \item Alter, Geschlecht, \EE{Hautfarbe}, Gesichtsausdruck,
  \EE{Haltung}, spontane Bewegungen,  Sprache
  \item Offensichtliche Verletzungen, Blutungen
  \item Sonstige \EE{Auffälligkeiten}
\end{itemize}
\end{minipage}
&
\noindent\color{MuwBlueDark}\begin{minipage}[t]{\linewidth}
\begin{itemize}
%   \item[\SymbolPkuNormal] 
  \item[\SymbolPkuVariante] Blässe
  \item[\SymbolFindingPathological] Schohnhaltung, Zynose
\end{itemize}
\end{minipage}
\\\addlinespace\midrule\addlinespace
\noindent\TabHeading{Bewusstsein}
&
\begin{minipage}[t]{\linewidth}
\begin{itemize}
  \item \EE{Bewusstseinsgrad} (WASB)
\end{itemize}
\end{minipage}
&
\noindent\color{MuwBlueDark}\begin{minipage}[t]{\linewidth}
\begin{itemize}
  \item[\SymbolPkuNormal] \textbf{W}ach
  \item[\SymbolPkuVariante] Reaktion auf \textbf{A}nsprache
  \item[\SymbolFindingPathological] Reaktion auf
  \textbf{S}chmerzreiz, \textbf{B}ewusstlos
\end{itemize}
\end{minipage}
\\\addlinespace\midrule\addlinespace
\noindent\TabHeading{Hauptbeschwerde}
&
\begin{minipage}[t]{\linewidth}
\begin{itemize}
  \item 
  \EE{Leitsymptom(e)}
\end{itemize}
\end{minipage}
&
\noindent\color{MuwBlueDark}\begin{minipage}[t]{\linewidth}
\end{minipage}
\\\addlinespace\toprule\addlinespace
\noindent\T{A}

\TabHeading{Atemweg}
&
\begin{minipage}[t]{\linewidth}
\begin{itemize}
  \item *Inspektion der Atemwege (\EE{Mund}, \EE{Nase})
  \item \EE{Atemgeräusche}
\end{itemize}
\end{minipage}
&
\noindent\color{MuwBlueDark}\begin{minipage}[t]{\linewidth}
\begin{itemize}
  \item[\SymbolPkuNormal] Atemwege frei
%   \item[\SymbolPkuVariante]
  \item[\SymbolFindingPathological] Atemwege verlegt
\end{itemize}
\end{minipage}
\\\addlinespace\midrule\addlinespace
\noindent\T{B}

\TabHeading{Atmung}
&
\begin{minipage}[t]{\linewidth}
\begin{itemize}
  \item Schätzen: \EE{Atemfrequenz}, \EE{-tiefe}
  \item *\EE{Sp\ce{O2}}
  \item Atemgeräusche
  \item Inspektion der \EE{Thorax-Atembewegungen}, ggfs. 
  Atemprobe
  \item \E{Zeichen der
  Atemnot}: Hautfarbe, Bauch-/Thorax\-bewegungen,
  Anstrengung beim Atmen, Atemhilfsmuskulatur, \ldots
  \item \EE{Auskultation}: Seitenvergleich, Lungenbasen
\end{itemize}
\end{minipage}
&
\noindent\color{MuwBlueDark}\begin{minipage}[t]{\linewidth}
\begin{itemize}
  \item[\SymbolPkuNormal] Atmung normfrequent, normale
  seitengleiche Atemexkursioen, vesikuläres seitengleiches
  Atemgeräusch ohne Nebengeräusche, Eupnoe, Sp\ce{O2}
  \VonBis{94}{100}\PC*
%   \item[\SymbolPkuVariante]
%   \item[\SymbolFindingPathological] 
\end{itemize}
\end{minipage}
\\\addlinespace\midrule\addlinespace
\noindent\T{C}

\TabHeading{Kreislauf}
&
\begin{minipage}[t]{\linewidth}
\begin{itemize}
  \item \EE{Hautfarbe}, \EE{-temperatur}
  \item Schweiß
  \item \EE{Rekap}-Zeit
  \item \EE{Radialispuls}: Stärke, Rhythmus, Abschätzen der
  Frequenz
\end{itemize}
\end{minipage}
&
\noindent\color{MuwBlueDark}\itshape\begin{minipage}[t]{\linewidth}
\begin{itemize}
  \item[\SymbolPkuNormal] Haut rosig und warm, Rekap-Zeit
  unter 2\Sec*, Radialispuls rhythmisch und gut palp.,
  normfrequent
%   \item[\SymbolPkuVariante]
%   \item[\SymbolFindingPathological] 
\end{itemize}
\end{minipage}
\\\addlinespace\cmidrule{2-3}\addlinespace
\noindent\T{STU}

\TabHeading{*Schnelle Trauma-Untersuchung}
&
\begin{minipage}[t]{\linewidth}
\begin{itemize}
  \item Inspektion
  und Abtasten von
  \begin{compactenum}
    \item \EE{Kopf} inkl. Ohren,
    \item \EE{Hals},
    \item \EE{Thorax},
    \item \EE{Bauch},
    \item \EE{Becken} und
    \item \EE{Oberschenkel}
  \end{compactenum}
\end{itemize}
\end{minipage}
&
\noindent\color{MuwBlueDark}\begin{minipage}[t]{\linewidth}
\begin{itemize}
  \item[\SymbolPkuNormal] Kein Hinweis auf Blutungen oder
  akut lebensbedrohliche Verletzungen
%   \item[\SymbolPkuVariante]
%   \item[\SymbolFindingPathological] 
\end{itemize}
\end{minipage}
\\\addlinespace\midrule\addlinespace
\noindent\T{D}

\TabHeading{Schnelle\newline Untersuchung}
&
\begin{minipage}[t]{\linewidth}
\begin{itemize}
  \item Vitalwerte (\EE{HF}, \EE{RR})
  \item Neurologisch:
  \begin{itemize}
    \item \EE{Orientierung} oder *\EE{GCS}
    \item \EE{Pupillen}
    \item \EE{Kraftprobe OE}, *UE
    \item *Kann der
    Patient \EE{Hände} und \EE{Füße} spüren und aktiv
    bewegen?
  \end{itemize}
  \item *\EE{Blutzuckermessung}
\end{itemize}
\end{minipage}
&
\noindent\color{MuwBlueDark}\begin{minipage}[t]{\linewidth}
\begin{itemize}
  \item[\SymbolPkuNormal] Vitalwerte im Normbereich,
  Allseits orientiert, Pupillen rund, mittelweit, isokor,
  prompte konsensuelle Lichtreaktion, Kraftprobe OE/UE
  seitengleich u. kräftig, Sensorik u. Motorik der Pripherie
  intakt.
%   \item[\SymbolPkuVariante]
%   \item[\SymbolFindingPathological] 
\end{itemize}
\end{minipage}
\\\addlinespace\midrule\addlinespace
\noindent\T{E}

\TabHeading{Erweiterte\newline Untersuchung}

% {\tiny (Wenn Einschätzung der vitalen
% Bedrohung oder Hauptbeschwerde unklar ist.)}
&
\begin{minipage}[t]{\linewidth}
\begin{itemize}
  \item (Fremd-) \EE{Anamnese} 
  \item Einschätzen der \EE{Umgebung}
  \item Weitere relevante \EE{Untersuchungen} (EKG,
  Ultraschall, Röntgen, \ldots)
  \item Arbeitsdiagnose erstellen und Differentialdiagnosen
  ausschließen
\end{itemize}
\end{minipage}
&
\noindent\color{MuwBlueDark}\begin{minipage}[t]{\linewidth}
% \begin{itemize}
%   \item Transportentscheidung 
%   \item Spezielle Maßnahmen gemäß Verdachtsdiagnose
% \end{itemize}
~
\end{minipage}
% \\\addlinespace\bottomrule\addlinespace
% \TabHeading{Beurteilung}
% &
% \begin{minipage}[t]{\linewidth}
% \begin{itemize}
%   \item \EE{Einschätzen der vitalen Bedrohung}
%   ständig während des gesamten Blocks
%   \item Vorläufige Verdachtsdiagnose formulieren
% \end{itemize}
% \end{minipage}
% &
% \noindent\color{MuwBlueDark}\begin{minipage}[t]{\linewidth}
% \begin{itemize}
%   \item {\TxMassVit}
%   \item Reihenfolge der weiteren Untersuchungen und
%   Maßnahmen festlegen
% \end{itemize}
% \end{minipage}
\\\addlinespace\bottomrule
\end{tabularx}

\begin{center}
%   \EE{Der Einschätzungsblock muss in regelmäßigen Abständen
% in angemessenem Umfang wiederholt
% werden (Verlaufskontrolle)!} \tiny

{ Mit einem * versehene Punkte werden nur durchgeführt,
wenn sie \E{situationsentsprechend} sind. 
% 
% Bei \Ca{absolut zeitkritischen Patienten}, bei denen bereits die 
% ABC-Einschätzung ergibt dass ein Transport nicht aufschiebbar ist, 
% kann u.\,U. schon \Ca{nach dem Punkt C eine vorgezogene Strategie- 
% und Transportentscheidung} getroffen werden. 
% Die Maßnahmen von D und E sollen dann, sofern möglich, während des 
% Transportes erfolgen. 
}
\end{center}
}
\end{minipage}
\end{minipage}%\vspace{-\baselineskip}

\end{WidePar}
\end{table}
% 
% \begin{table}[H]
% \FormatTable
% \begin{supertabular}{m{67.85mm}m{93.43mm}}
% \toprule\addlinespace
% \multicolumn{2}{m{163.28mm}}{
% \centering \TabHeading{Untersuchungsparameter für
% Notfallpatienten}} 
% \\\addlinespace\midrule\addlinespace
% Bewusstseinslage
%  &
%  Reaktion auf Ansprechen, Berührung
% und Schmerzreize zur Beurteilung der Komatiefe 
% \\\addlinespace\midrule\addlinespace
% Atemwege
%  &
% Atemgeräusche (Stridor, Gurgeln, \ldots), Gestik / Haltung
% \\\addlinespace\midrule\addlinespace
% Atmung
%  &
% Atemfrequenz, -tiefe, 
% Atemgeräusche, Auskultation der Lunge
% \\\addlinespace\midrule\addlinespace
% Herz-Kreislauf
%  &
%  Blutdruck, Pulsfrequenz,
% Pulsqualität, Arrhythmie, Herzgeräusche
% \\\addlinespace\midrule\addlinespace
% Aussehen,
% Geruch
%  &
% Zyanose, Blässe, Hautturgor,
% Verletzungen
% 
% (inkl. Blutungen)
% 
%  Urämie, Alkohol,
% Ketonkörper
% \\\addlinespace\midrule\addlinespace
% ZNS
%  &
%  Pupillenform, Pupillenreaktion auf
% Licht, Meningismus, Paresen, Tonus, Reflexe
% \\\addlinespace\midrule\addlinespace
% Abdomen
%  &
%  Druckschmerz, Abwehrreaktion,
% Darmgeräusche, Organvergrößerungen
% \\\addlinespace\midrule\addlinespace
% Weiterführende
% Untersuchungen
%  &
%  Labor (Blutzucker, Blutgase,
% Elektrolyte, etc.), EKG, bildgebende Verfahren
% \\\addlinespace\bottomrule
% \end{supertabular}
% \end{table}


\begin{Literary}

\fullcite{Aass-1.0var1}

\fullcite{AassMk-2013var1}

\fullcite{TertialeNotfallIntensivmed1}

\fullcite{SemmelABC1}

\end{Literary}

\printbibliography

\ChapterEnd